\PassOptionsToPackage{dvipsnames}{xcolor}
\documentclass[preprint]{aastex702}
\usepackage{showyourwork}
% Margin icons on the verified equations. When every check passes, check_equations.py writes
% output/equation_icons.tex: an \addvalue{<label>_check}{<notebook>} line per registry equation.
% Each such \label then gets \EquationCheckSymbol in the margin, linking to its notebook; other
% labels keep showyourwork's icons (the GitHub logo on figures).
\newcommand\EquationCheckSymbol{\faCheckCircle}
\definecolor{equation-check}{rgb}{0.10,0.50,0.22}  % green (#1A7F37), distinct from showyourwork's blue icons
\newcommand\EquationCheckIcon{{\color{equation-check}\EquationCheckSymbol}\xspace}
\makeatletter
\ifdefined\usevalue  % the compile pass (showyourwork's preprocessing pass has its own label format)
  \LetLtxMacro\eqcheck@sywlabelformat\showkeyslabelformat
  \renewcommand\showkeyslabelformat[1]{%
    \IfEq*{\usevalue{#1_check}}{KEYNOTFOUND}{\eqcheck@sywlabelformat{#1}}{%
      \normalfont\raisebox{0pt}[6pt][0pt]{%
        \href{\GitHubURL/blob/\GitHubSHA/\usevalue{#1_check}}{\color{equation-check}\EquationCheckSymbol}}}}
\fi
\makeatother
\IfFileExists{output/equation_icons.tex}{\input{output/equation_icons.tex}}{}
% Notebooks in notebooks/ that accompany a section. Each gets a GitHub icon in the margin at the
% section's \label, linking to the notebook (showyourwork draws one at every \label with a
% <label>_script value), and the section cites it in a footnote with \NotebookLink.
\providecommand\addvalue[2]{}  % undefined in showyourwork's preprocessing pass
\newcommand\NotebookLink[1]{\href{\GitHubURL/blob/\GitHubSHA/notebooks/#1}{\texttt{\detokenize{#1}}}}
\addvalue{subsec:ft_cosine_script}{notebooks/ft_spot_rotation.ipynb}
\addvalue{subsec:ft_envelope_script}{notebooks/ft_squared_trapezoid.ipynb}
\addvalue{subsec:averaging_script}{notebooks/averaging_random_params.ipynb}
\addvalue{subsec:kernel_anal_case_script}{notebooks/closed_form_kernel_n2.ipynb}
\addvalue{subsec:harmonic_n2_script}{notebooks/lower_envelope.ipynb}
\addvalue{app:fourier_trapezoid_script}{notebooks/ft_trapezoid.ipynb}

\usepackage{theorem}
\usepackage{amsmath}
\usepackage{amssymb}
\usepackage{mathtools}
\usepackage{enumitem}
\usepackage{graphicx}
\usepackage{mathrsfs}
\usepackage{booktabs}
\usepackage{float}
\usepackage{morefloats}
\usepackage{etoolbox}
\usepackage{multirow}
\usepackage{rotating}
\usepackage{longtable}
\usepackage{xspace}
\usepackage{tikz}
\usepackage{listings}
% \hypersetup{hypertexnames=false}

% Python code listings (Appendix E)
\lstset{
  language=Python,
  basicstyle=\footnotesize\ttfamily,
  keywordstyle=\bfseries,
  commentstyle=\itshape\color{gray},
  stringstyle=\color{gray},
  showstringspaces=false,
  columns=fullflexible,
  keepspaces=true,
  breaklines=true,
  frame=single,
  framerule=0.3pt,
  rulecolor=\color{lightgray},
  xleftmargin=1em,
  aboveskip=1em,
  belowskip=1em,
}

\newcommand{\FT}[1]{\hat{#1}}
\newcommand{\E}{\mathbb{E}}
\newcommand{\Cov}{\mathrm{Cov}}
\newcommand{\Var}{\mathrm{Var}}
\newcommand{\sinc}{\mathrm{sinc}}

\newcommand{\kepler}{{\it Kepler}\xspace}
\newcommand{\tess}{{\it TESS}\xspace}
\newcommand{\ktwo}{{\it K2}\xspace}

\newcommand{\xxx}[1]{{\color{red}{[#1]}}}

% Kepler-63 year colors
\definecolor{colorYear1}{named}{blue}
\definecolor{colorYear2}{named}{red}
\definecolor{colorYear3}{HTML}{a333d4}
\definecolor{colorYear4}{HTML}{24bf3c}
\definecolor{colorYearAll}{named}{black}
\newcommand{\code}[1]{{\texttt{#1}}}

\newcounter{ftprop}
\renewcommand{\theftprop}{FT\arabic{ftprop}}

\newcounter{ftpair}
\renewcommand{\theftpair}{P\arabic{ftpair}}

\begin{document}

\title{\sc Physically Motivated Gaussian Processes for Stellar Variability I: \\
A Starspot Kernel with Differential Rotation}
\shorttitle{starspot gaussian process}
\shortauthors{birky et al.}

\correspondingauthor{Jess Birky}
\email{jess@birky.net}

\author[0000-0002-7961-6881]{Jessica Birky}
\affiliation{Department of Astronomy, University of Washington, 3910 15th Avenue NE, Seattle, WA 98195, USA}
\affiliation{Homer L. Dodge Department of Physics and Astronomy, University of Oklahoma, Norman, OK 73019, USA}
\email{jess@birky.net}
\author[0000-0002-0802-9145]{Eric Agol}
\affiliation{Department of Astronomy, University of Washington, 3910 15th Avenue NE, Seattle, WA 98195, USA}
\email{agol@uw.edu}
\author[0000-0001-6487-5445]{Rory Barnes}
\affiliation{Department of Astronomy, University of Washington, 3910 15th Avenue NE, Seattle, WA 98195, USA}
\email{rkb9@uw.edu}
\author[0000-0002-0637-835X]{James Davenport}
\affiliation{Department of Astronomy, University of Washington, 3910 15th Avenue NE, Seattle, WA 98195, USA}
\email{jrad@uw.edu}
\author[0000-0001-9590-2274]{Sean Matt}
\affiliation{Homer L. Dodge Department of Physics and Astronomy, University of Oklahoma, Norman, OK 73019, USA}
\email{sean@ou.edu}

\begin{abstract}
We present a novel, physically motivated Gaussian process (GP) kernel designed to model starspot variability from photometric lightcurves. Using an analytic flux model, we derive an analytic kernel function that describes the temporal correlations induced by rotating starspots as a function of spot properties and stellar rotation: equitorial rotation period ($P_{\rm eq}$), differential rotation ($\kappa$), inclination angle ($I$), spot lifetimes ($\ell_{\rm spot}$), spot emergence and decay timescales ($\tau_{\rm spot}$), maximum spot size ($\alpha_{\rm max}$), spot contrast ($f_{\rm spot}$), and spot emergence rate ($\dot{N}_{\rm spot}$). We thoroughly validate our derivation using \code{SymPy} and compare our analytic solutions against to numerical simulations. Furthermore, we interpret how each parameter affects the kernel structure, showing that $\alpha_{\rm max}$, $f_{\rm spot}$, and $\dot{N}_{\rm spot}$ are degenerate parameters controlling the amplitude of the correlations (which can be re-parameterized into a single parameter $\sigma_k$), while $P_{\rm eq}$, $\kappa$, and $I$ control the periodicity and harmonic structure, and $\ell_{\rm spot}$, $\tau_{\rm em}$, and $\tau_{\rm dec}$ control the decay of the correlations over time. Our implementation is available in the package \code{spotgp}, which uses \code{JAX} for efficient automatic differentiation and GPU acceleration. Furthermore, our implementation makes use of a sparse banded matrix solver (with bandwidth $b$), reducing the computational complexity of GP inference from $\mathcal{O}(N^3)$ to $\mathcal{O}(Nb^2)$ for $N$ data points. Using simulated lightcurves, we use Markov chain Monte Carlo (MCMC) and Cramer-Rao bound analysis to systematically explore how inference performs for different parameter regimes, and observational conditions (cadence, baseline, and noise level). Finally, our simulated inference demonstrates that \code{spotgp} is a highly promising tool for extracting differential rotation and inclinations just from \emph{individual} lightcurves of stars. \\
\end{abstract}

% % ========================================================================================
\section{Introduction} \label{sec:introduction}

High-precision photometric surveys such as \emph{Kepler} \citep{borucki_kepler_2010}, \emph{K2} \citep{howell_k2_2014}, and \emph{TESS} \citep{ricker_transiting_2014} have been transforming our understanding of fundamental stellar physics, revealing the effects of stellar variability and dynamics for stars across a wide range of masses, ages, and Galactic environments. In particular, precise time-domain data have enabled constraints on the rotation periods of thousands of stars \citep{mcquillan_rotation_2014,gordon_stellar_2021,boyle_tess_rotation_2026}, which has been crucial for understanding the evolution of stellar angular momentum \citep{weber_angular_1967,skumanich_time_1972}, the surface variability and magnetic activity of stars \citep{santos_starspot_2017,forgacs-dajka_time-dependent_2021}, and the stellar dynamo process \citep{brun_differential_2017,kapyla_simulations_2023,noraz_impact_2022}.

Over time, stars lose angular momentum through magnetized stellar winds, causing them to spin down as they age. The rate of this spin-down depends on the star's mass, age, and magnetic field strength, and is driven by the torque from stellar winds flowing along open magnetic field lines, a process known as \emph{magnetic braking}. Early work established that $P_{\rm rot} \propto t^{-1/2}$ \citep{skumanich_time_1972}, opening up the field of gyrochronology for inferring stellar ages from rotation periods \citep{barnes_rotational_2003,barnes_ages_2007}. Further work has developed empirical relations for how magnetic braking torque depends on stellar properties \citep{kawaler_angular_1988,sills_angular_2000,reiners_radius-dependent_2012,vansaders_fast_2013,matt_mass-dependence_2015,amard_impact_2020,ireland_effect_2021}. However open questions remain: recent studies have discovered evidence for weakened magnetic braking in rapidly rotating old field stars \citep{van_saders_weakened_2016,saunders_stellar_2024}, suggesting that the braking process may be more complex than originally thought.

Much of our knowledge of stellar rotation comes from photometric variability driven by starspots. Low-mass stars ($M \lesssim 1.5 \, M_{\odot}$) possess outer convective envelopes in which turbulent motions interact with rotation and magnetic fields to drive a stellar dynamo \citep{brun_differential_2017}. Differential rotation (the latitudinal variation in angular velocity) plays a fundamental role in this process: it shears poloidal magnetic field lines into a toroidal configuration (the $\Omega$-effect), which is a key mechanism driving the activity cycle \citep{noraz_impact_2022}. Numerical simulations of rotating convection predict that the magnitude and geometry of differential rotation depend on the stellar mass, rotation rate, and convective structure \citep{augustson_convection_2012,guerrero_differential_2013,kapyla_simulations_2023}. In particular, these simulations predict a transition from solar-like differential rotation (fast equator, slow poles) to anti-solar differential rotation (fast poles, slow equator) as the Rossby number $Ro = P_{\rm rot}/\tau_{\rm conv}$ increases past $Ro \sim 1$. Many starspot properties, including spot lifetimes \citep{giles_kepler_2017}, latitude distributions, and emergence patterns, are also governed by $Ro$, making it a highly important dimensionless parameter for characterizing the stellar dynamo regime \citep{gastine_solarlike_2013,noraz_hunting_2022}. However, this predicted transition in differential rotation remains largely untested observationally, because measuring differential rotation for large samples of stars has been extremely challenging.

Existing techniques for measuring stellar differential rotation include: Doppler imaging \citep{collier_cameron_differential_2007}, the Fourier transform method for spectral line profiles \citep{reiners_radius-dependent_2012}, asteroseismology \citep{benomar_asteroseismic_2018}, and transit spot-crossing \citep{morris_starspots_2017}. However, each have each been limited to samples of $\sim$tens of stars, and are restricted to specific stellar types (e.g., fast rotators, bright stars, or transiting systems). On the other hand, the high-precision, long-baseline photometric datasets from \kepler and \tess encode information on the spot variability for hundreds of thousands of stars, which has yet to be fully leveraged. However, extracting differential rotation from disk-integrated photometry is notoriously difficult: \citet{aigrain_testing_2015} conducted a blind hare-and-hounds exercise on 1,000 simulated \kepler-like lightcurves and found \emph{little correlation between the reported and simulated values of the differential rotation shear}. \citet{morris_solar_2019} reinforced this conclusion using the Sun itself as a benchmark, finding no detectable differential rotation signal even in a noise-free solar lightcurve reconstructed from archival sunspot records. These results suggest that the stochastic nature of spot emergence at a broad range of latitudes causes the measured period to reflect an average over active latitudes, rather than revealing the latitude-dependent rotation law.

Modeling the magnetohydrodynamical processes that accurately reproduce the convection, magnetic fields, and dynamo process of low-mass stars is inherently complex and computationally expensive \citep{miesch_threedimensional_2000,miesch_structure_2008,charbonneau_dynamo_2020}. As an alternative, analytic spot models \citep{kipping_analytic_2012,luger_starry_2019} treat the stellar surface as hosting an arbitrary number of circular spots that emerge and decay at different latitudes, each modulating the disk-integrated flux according to its projected area and contrast. These models have been successful in matching the statistical properties of observed lightcurves \citep{aigrain_testing_2015}, while providing a physically interpretable framework for the spot geometry.

Although there are approaches for forward modeling surfaces with spots, inferring information about \emph{individual} star spots is an ill-posed problem with many degeneracies \citep{luger_starry_2019,basri_information_2020}. \citet{basri_information_2020} showed that the single-dip ratio (a metric for whether a lightcurve shows one or two dips per rotation) is only weakly sensitive to differential rotation shear and is strongly confounded by spot lifetime, concluding that differential rotation may be ``nearly impossible to accurately infer from light curves alone unless spots live for many rotations.''
A Gaussian process \citep[GP;][]{rasmussen_gaussian_2006} offers an alternative approach that treats starspot variability as a stochastic process and models the \emph{distribution} of spot properties rather than individual spots.
GPs have become a standard tool for modeling stellar rotation \citep{angus_inferring_2018,foreman-mackey_fast_2017,gordon_stellar_2021}, but the most widely used quasi-periodic (QP) kernel has no parameter that encodes differential rotation, meaning that any differential rotation signal is absorbed into the other hyperparameters in an uncontrolled way. \citealt{nicholson_quasi-periodic_2022} systematically tested the mapping between QP kernel hyperparameters and physical spot properties, finding that while the rotation period and evolution timescale are reliably recovered, the harmonic complexity parameter has no intuitive physical interpretation. These results collectively motivate the development of a physically grounded GP kernel: one whose hyperparameters map directly to measurable stellar properties including differential rotation.

In this paper, we derive a physically motivated GP kernel for stellar rotation from first principles. Starting from an analytic starspot model \citep{kipping_analytic_2012}, we obtain a closed-form kernel function that naturally encodes the effects of spot geometry, evolution, and differential rotation. We validate the analytic kernel against numerical Monte Carlo simulations, demonstrate that it captures harmonic structure and non-sinusoidal periodicity absent in standard quasi-periodic and simple harmonic oscillator kernels, and show that it provides significant computational speedups over numerical approaches. The paper is organized as follows: Section~\ref{sec:lightcurve} describes the lightcurve model, Section~\ref{sec:kernel_derivation} derives the GP kernel, Section~\ref{sec:mcmc} presents the \code{JAX} implementation of the kernel and MCMC inference framework, Section~\ref{sec:results} validates the analytic kernel against numerical results, Section~\ref{sec:discussion} discusses the physical interpretability of our results with intuitive visualizations, and the Appendix Section provides detailed derivations and \code{sympy} validation notebooks. Equations marked with \EquationCheckIcon in the margin are verified with \code{SymPy} against their derivation and, for the \code{spotgp} functions listed in Appendix~\ref{app:equation_checks}, numerically against \code{spotgp}; the icon links to the verifying notebook, and Appendix~\ref{app:equation_checks} lists all verified equations. A companion paper \citep[][hereafter Paper~II]{birky_starspot_2026b} extends the kernel with additive composite terms for several coexisting spot populations and for non-spot variability, and applies the full framework to the four-year \kepler lightcurve of Kepler-63.

\vspace{\baselineskip}
% \clearpage
% ========================================================================================
\section{Lightcurve Model} \label{sec:lightcurve}

We use the analytic starspot model defined in \citet{kipping_analytic_2012}. The flux model for $N_{\rm spot}$ spots with contrast $f_k$ is:
\begin{equation}
    F(t) = 1 - \frac{1}{\pi} \sum_{k=1}^{N_{\rm spot}} A_k(t) \cdot (1 - f_{\rm spot,k}),
    \label{eqn:flux_model}
\end{equation}
where $A_k(t)$ is the projected area of spot $k$ at time $t$ and $f_{\rm spot,k}$ is the contrast of the spot (where $f_{\rm spot} = 0$ represents a totally black spot, $0 < f_{\rm spot} < 1$ represents a partially dark spot, $f_{\rm spot} = 1$ represents no contrast, and $f_{\rm spot} > 1$ represents a bright spot or stellar faculae). The flux deficit at a given time is the sum of the contributions from all visible spots on the surface of the star, where each contribution is determined by the projected area of the spot and its contrast relative to the stellar photosphere. 

% ------------------------------------------------------------------
\subsection{Single-Spot Projected Area}

For each individual spot, the flux deficit is determined by the projected area of the spot on the visible hemisphere of the star. For ease of notation we drop the spot index $k$ and write the projected area of a single spot as $A(t)$, however these equations are applied to each spot $k$ with its own time-dependent area $A_k(t)$ and angle $\beta_k(t)$. This projected area of starspot $k$ at time $t$ is given by:
\begin{equation}
    A_k(t) 
        = \mathbb{R}\Big[ \cos^{-1}[\cos\alpha_k\csc\beta_k]
        + \cos\beta_k\sin\alpha_k \Xi_k - \cos\alpha_k\sin\beta_k \Psi_k \Big]
\end{equation}
where the intermediate terms are
\begin{align}
    \Xi_k &= \sin\alpha_{k}\cos^{-1}[-\cot\alpha_{k}\cot\beta_{k}],\\
    \Psi_k &= \sqrt{1-\cos^2\alpha_k\csc^2\beta_k}
\end{align}

Expanded out, $A(t)$ becomes a piecewise expression with three cases:
\begin{equation}
  A_{k}(t) = \left\{ \begin{array}{ll}
    \pi \sin^2\alpha_{k} \cos \beta_{k} & (0< \beta_{k} < \pi/2 -\alpha_{k}) \\[6pt]
    \cos^{-1}[\cos \alpha_{k} \csc \beta_{k}]  +\cos \beta_{k} \sin^2 \alpha_{k} \cos^{-1}[-\cot \alpha_{k} \cot \beta_{k}] & \\
    \hspace{4cm} - \cos \alpha_{k} \sin \beta_{k} \sqrt{1-\cos^2 \alpha_{k} \csc^2 \beta_{k}} & (\pi/2 -\alpha_{k} < \beta_{k} < \pi/2 +\alpha_{k}) \\[6pt]
    0 & (\pi/2 +\alpha_{k} < \beta_{k} < \pi)
  \end{array} \right.
  \label{eq:A_exact}
\end{equation}
where $\alpha_k$ is the angular radius of the spot and $\beta_k$ is the angle between the spot normal and the line of sight \citep{kipping_analytic_2012}. The three cases correspond to distinct geometric configurations of the spot on the stellar disk: fully visible, partiay occulted by the limb, and fully hidden on the far side, as shown in Figure~\ref{fig:spot_geometry}. Paper~II presents a self-contained derivation of Eq.~\eqref{eq:A_exact} as a surface integral over the spot cap.

\begin{figure}[ht!]
\begin{center}
    \script{plot_spot_geometry.py}
    \includegraphics[width=\linewidth]{figures/spot_geometry.pdf}
    \caption{Illustration of the three geometric regimes for the projected area of a circular starspot on a rotating star. Each panel shows the stellar disk as seen by the observer, with latitude and longitude grid lines, the spot (dark red), the angular radius $\alpha_k$, and the angle $\beta_k$ between the spot normal and the line of sight. \textbf{Left:} The spot is fully visible on the near hemisphere. \textbf{Center:} The spot straddles the stellar limb and is partially occulted. \textbf{Right:} The spot is on the far side and hidden from view.}
    \label{fig:spot_geometry}
\end{center}
\end{figure}

The angle between spot normal and the line of sight for a single spot at a given time is:
\begin{equation}
    \cos \beta(t) =  \cos I \sin\Phi + \sin I \cos\Phi\cos\Lambda(t),
    \label{eq:cos_beta}
\end{equation}
where $I$ is the inclination of the stellar rotation axis, $\Phi$ is the latitude of the spot, and $\Lambda(t)$ is the longitude of the spot as a function of time. We emphasize the distinction between the two frames related by Eq.~\eqref{eq:cos_beta}: $\Phi$ is a \emph{stellar} coordinate, measured from the stellar equator in the rotation frame, with $\Phi \in [-\pi/2, \pi/2]$ and $\Phi = +\pi/2$ the pole inclined toward the observer; $I \in [0, \pi/2]$ is the angle between the rotation axis and the line of sight, so that the sub-observer point lies at latitude $\pi/2 - I$. Defining $\Phi$ in the rotation frame is what makes the differential rotation law $\omega_0(\Phi)$ (Eq.~\ref{eq:kernel_full}) well posed, since latitudinal shear is an intrinsic property of the star and cannot depend on the viewing geometry. Consequently the inclination enters the model not through the coordinate $\Phi$ itself, but through the observer-dependent visibility mask it imposes on that coordinate: since $\cos\beta$ ranges over $[\sin(\Phi - I),\, \sin(\Phi + I)]$ as the star rotates, spots at $\Phi \geq I$ are circumpolar and permanently visible, spots at $\Phi \leq -I$ never rotate into view, and only the band $|\Phi| < I$ produces rotationally modulated flux. The time dependence of $\beta(t)$ arises from the rotation of the star, which causes the spot to move in longitude and therefore change its angle relative to the line of sight. The projected area $A(t)$ is then determined by both $\alpha(t)$ and $\beta(t)$ according to the piecewise function in Eqn~\ref{eq:A_exact}.

When the spot angular radius $\alpha_k$ is small compared to 1~radian (typical for Sun-like stars where spots subtend $\lesssim\!6^\circ$), we can make the approximation $\sin\alpha_k \approx \alpha_k$, so $\sin^2\!\alpha_k \approx \alpha_k^2$. In this limit:
\begin{itemize}
    \item The \textbf{fully visible} case simplifies to $A_k \approx \pi\alpha_k^2\cos\beta_k$.
    \item The \textbf{partial-visibility window} has width $2\alpha_k$ and becomes vanishingly narrow, so the spot transitions nearly instantaneously from fully visible to fully hidden as it crosses the limb.
    \item The \textbf{geometric projection effect} is purely $\cos\beta$: the same projection factor as a flat disk viewed at angle $\beta$. Higher-order corrections from spherical curvature (the $\Xi_k$, $\Psi_k$ terms) vanish at leading order in $\alpha_k$.
\end{itemize}
This allows us to replace the entire three-case piecewise function with a single expression:
\begin{equation}
    A(t) \approx \pi\, \alpha^2(t) \; \cdot \; \max\!\big\{\cos\beta(t),\; 0\big\}
    \label{eq:A_smallspot}
\end{equation}
where $\alpha(t)$ is the angular radius of the spot (whose time evolution is described in Section~\ref{sec:spot_profile}). Figure~\ref{fig:projected_area} illustrates this approximation. The left panel shows the size of spots with three different angular radii $\alpha_{\max}$ on the stellar disk at $\beta = 0$ (disk center), providing a visual reference for the spot-to-star size ratio. The right panel compares the exact projected area $A(\beta)$ from Eq.~\eqref{eq:A_exact} (solid lines) with the small-spot approximation $\max\{\cos\beta, 0\}$ (dashed line). For small spots ($\alpha_{\max} \lesssim 15^\circ$), the exact and approximate curves are nearly indistinguishable, validating the factorization. For larger spots ($\alpha_{\max} \sim 30^\circ$), the transition at the limb ($\beta = \pi/2$) becomes more gradual in the exact solution, but the approximation remains reasonable.

\begin{figure}[ht!]
\begin{center}
    \script{plot_projected_area.py}
    \includegraphics[width=0.85\linewidth]{figures/fig_projected_area.pdf}
    \caption{\textit{Left:} Stellar disk viewed face-on, with circles indicating the projected size of spots with angular radii $\alpha_{\max} = 5^\circ$, $15^\circ$, and $30^\circ$ centered at $\beta = 0$ (disk center). \textit{Right:} Normalized projected area $A(\beta)/A(0)$ as a function of the angle $\beta$ between the spot normal and the line of sight. Solid colored lines show the exact piecewise expression (Eq.~\ref{eq:A_exact}) for each $\alpha_{\max}$; the dashed black line shows the small-spot approximation $\max\{\cos\beta,\,0\}$ (Eq.~\ref{eq:A_smallspot}). The two converge for $\alpha_{\max} \lesssim 15^\circ$.}
    \label{fig:projected_area}
\end{center}
\end{figure}

The key consequence for the kernel derivation is that the projected area now \emph{factorizes} into a product of two functions:
\begin{equation}
    A(t) \propto \Gamma(t) \cdot \mathcal{V}(t),
\end{equation}
where we define $\Gamma(t) \equiv \alpha^2(t)/\alpha_{\max}^2$ as the normalized squared spot-size envelope (with unit peak), and $\mathcal{V}(t) \equiv \max\{\cos\beta(t),\, 0\}$ as the spot visibility function. Figure~\ref{fig:spot_decomposition} illustrates this decomposition. 
In this paper we proceed to derive the GP kernel under the simplifying assumptions of small spot approximation and negligible limb darkening. These assumptions allow us to solve for an analytic, closed-form solution for the kernel function, however we will refer the reader to Paper II for the numerical approach for how to compute the full visibility function including these terms.


% ------------------------------------------------------------------
\subsection{The Emergence and Decay Profile of Spots} \label{sec:spot_profile}

In this work, we assume that the angular size of the spot, $\alpha(t)$, evolves over time as the spot emerges and decays. Here we test two different parameterizations of the spot evolution: a piecewise linear function and an exponential growth/decay function. In both cases, the spot reaches its maximum angular size $\alpha_{\mathrm{max}}$ at time $t_{\mathrm{ref}}$, which is the reference time for the spot. The spot emerges over a timescale $\tau_{\rm em}$ and decays over a timescale $\tau_{\rm dec}$.

\textbf{Trapezoidal emergence/decay profile:} Following \citet{kipping_analytic_2012}, we define a piecewise linear function with an emergence timescale $\tau_{\rm em}$, a plateau duration $\ell_{\rm spot}$, and a decay timescale $\tau_{\rm dec}$. The spot holds its maximum angular size $\alpha_{\mathrm{max}}$ over a plateau centered on $t_{\mathrm{ref}}$, which is the reference time for the spot. The Heaviside step function $\mathsf{H}(x)$ is used to ensure that the spot only contributes to the flux deficit during its emergence and decay phases. The spot angular radius is given by:

\begin{equation}
    \frac{\alpha_{\rm trap}(t)}{\alpha_{\mathrm{max}}} 
        = \mathcal{I}^{-1} [\Delta t_1 \mathsf{H}(\Delta t_1) - \Delta t_2 \mathsf{H}(\Delta t_2)]  - \mathcal{E}^{-1} [\Delta t_3 \mathsf{H}(\Delta t_3) - \Delta t_4 \mathsf{H}(\Delta t_4)],
    \label{eqn:trapezoid}
\end{equation} 
where $\Delta t_1 = t - (t_{\mathrm{ref}} - \ell_{\rm spot}/2 - \tau_{\rm em})$, $\Delta t_2 = t - (t_{\mathrm{ref}} - \ell_{\rm spot}/2)$, $\Delta t_3 = t - (t_{\mathrm{ref}} + \ell_{\rm spot}/2)$, and $\Delta t_4 = t - (t_{\mathrm{ref}} + \ell_{\rm spot}/2 + \tau_{\rm dec})$. The normalization factors $\mathcal{I} = \tau_{\rm em}$ and $\mathcal{E} = \tau_{\rm dec}$ are the maximum values of the two bracketed ramp functions, which ensures that the peak of the trapezoidal function is normalized to 1, so that $\alpha_{\rm trap}(t)$ reaches $\alpha_{\mathrm{max}}$ at its maximum. The trapezoidal profile is a simple, piecewise linear model that captures the general shape of spot emergence and decay, with a linear rise to maximum size, a plateau, and a linear decay back to zero. For equal emergence and decay timescales, $\tau_{\rm em} = \tau_{\rm dec} \equiv \tau_{\rm spot}$, it reduces to the symmetric trapezoid of Eq.~\eqref{eq:alpha_piecewise} (with $t_{\rm ref} = 0$).

\textbf{Exponential emergence/decay profile:} An alternative, model for the spot-size evolution (similar to \citealt{aigrain_testing_2015}) uses exponential emergence and decay rather than linear ramps. The spot angular radius is:
\begin{equation}
    \alpha_{\rm exp}(t) = \alpha_{\max} \times
    \begin{cases}
        e^{(t - t_1)/\tau_{\rm em}} & t < t_1 \quad\text{(exponential rise)} \\[4pt]
        1 & t_1 \leq t \leq t_2 \quad\text{(plateau)} \\[4pt]
        e^{-(t - t_2)/\tau_{\rm dec}} & t > t_2 \quad\text{(exponential decay)}
    \end{cases}
    \label{eq:alpha_exp}
\end{equation}
where $t_1$ and $t_2 = t_1 + \tau_{\rm plat}$ mark the start and end of the plateau phase.
The three parameters are: $\tau_{\rm em}$ (emergence $e$-folding time), $\tau_{\rm plat}$ (plateau duration), and $\tau_{\rm dec}$ (decay $e$-folding time).
Without loss of generality we center the envelope by setting $t_1 = 0$, $t_2 = \tau_{\rm plat}$.

\begin{figure}[ht!]
\begin{center}
    \script{plot_spot_decomposition_panel.py}
    \includegraphics[width=\linewidth]{figures/spot_decomposition_panel.pdf}
    \caption{Decomposition of the single-spot projected area signal $A(t)$ into the normalized squared spot-size envelope $\Gamma(t) = \alpha^2(t)/\alpha_{\max}^2$ (top) and the visibility function $\mathcal{V}(t) = \max\{\cos\beta(t),\, 0\}$ (middle), with their product $A(t) \propto \Gamma(t) \cdot \mathcal{V}(t)$ (bottom). \textit{Left:} trapezoidal envelope (Eq.~\ref{eqn:trapezoid}). \textit{Right:} exponential envelope (Eq.~\ref{eq:alpha_exp}). Both use the same rotation period $P = 3$\,d, spot latitude $\Phi = 30^\circ$, and edge-on inclination $I = 90^\circ$.}
    \label{fig:spot_decomposition}
\end{center}
\end{figure}

% ------------------------------------------------------------------
\subsection{Lattitudinal and Longitudinal Evolution of Spots}

The position of a starspot on the surface is described by its latitude ($\Phi$) and longitude ($\Lambda$). Similarly to the Sun, the latitude of the spot is assumed to be constant over time, while the longitude evolves due to the rotation of the star:
\begin{align}
    \Phi &= \Phi_{\mathrm{ref}}, \\
    \Lambda(t) &= \Lambda_{\mathrm{ref}} + \frac{2\pi (t-t_{\mathrm{ref}})}{P(\Phi)}, 
\end{align}
We assume a Sun-like differential rotation profile, where the rotation period $P$ at latitude $\Phi_k$ depends on  the latitudinal shear $\kappa$ and the rotation period at the equator of the star $P_{\rm eq}$:
\begin{equation}
    P(\Phi_k) = \frac{P_{\mathrm{eq}}}{1 - \kappa \sin^2\Phi_k},
\end{equation}
where $\kappa = 0$ corresponds to solid body rotation and $\kappa > 0$ corresponds to solar-like differential rotation (equator rotates faster than poles), while $\kappa < 0$ corresponds to anti-solar differential rotation (poles rotate faster than equator). Or in other words, $\kappa$ quantifies the relative difference in rotation rate between the equator and the poles, $\kappa = (\Omega_{\rm eq} - \Omega_{\rm pole}) / \Omega_{\rm eq}$ for equatorial rotation rate $\Omega_{\rm eq} = 2\pi / P_{\rm eq}$ and polar rotation rate $\Omega_{\rm pole} = 2\pi / P_{\rm pole}$. Figure \ref{fig:differential_rotation} illustrates the effect of differential rotation on the rotational velocity at different latitudes on the stellar surface.

\begin{figure}[ht!]
\begin{center}
    \script{plot_differential_rotation.py}
    \includegraphics[width=0.8\linewidth]{figures/differential_rotation.pdf}
    \caption{Illustration of differential rotation on a stellar sphere viewed at inclination $I = 70^\circ$. Arrows show the rotational velocity at each latitude, with length proportional to the angular velocity $\omega(\Phi) = 2\pi/P(\Phi)$. \textit{Left:} Solar-like differential rotation ($\kappa = 1.0$), where the equator rotates faster than the poles. \textit{Right:} Anti-solar differential rotation ($\kappa = -1.0$), where the poles rotate faster than the equator.}
    \label{fig:differential_rotation}
\end{center}
\end{figure}

% \clearpage
% ------------------------------------------------------------------
\subsection{Summary of Parameters}

Tables~\ref{tab:star_parameters} and \ref{tab:spot_parameters} summarize the physical parameters of the lightcurve model. Table~\ref{tab:star_parameters} lists the stellar parameters that are held fixed for all spots on the star, while Table~\ref{tab:spot_parameters} lists the parameters that can vary for each individual spot.

\begin{center}
    \begin{longtable}{rcl} 
    \hline
        \multicolumn{3}{c}{\textit{Star variables}} \\
    \hline
        $P_{\rm eq}$        &&  rotation period at equator \\
        $\kappa$            &&  latitudinal differential rotation shear \\
        $I$                 &&  inclination angle between the rotation axis and the line of sight ($0$, $\pi/2$) \\
        $N_{\rm spot}$      &&  total number of spots observed within the simulated time $T_{\rm obs}$ \\
        $l_{\rm spot}$      &&  lifetime of each spot \\
        $\tau_{\rm spot}$   &&  emergence and decay timescale of each spot  \\
    \hline
    \caption{Stellar parameters and their definitions.}
    \label{tab:star_parameters}
    \end{longtable} 
\end{center}

\begin{center}
    \begin{longtable}{rcl} 
    \hline
        \multicolumn{3}{c}{\textit{Individual spot variables}} \\
    \hline
        $k$                     && index number for each spot \\
        $A_k(t)$                && projected area of star spot \\
        $\beta_k(t)$           && angle between spot normal and line of sight \\
        $\alpha_k(t)$          && angular size of spot (emerges and decays) \\ 
        $\Lambda_k(t)$          && starspot longitude (0, $2\pi$) \\
        $\Lambda_{\rm ref,k}$          && reference longitude of spot (0, $2\pi$) \\
        $\Phi_{\rm ref,k}$            && reference latitude of spot ($-\pi/2$, $\pi/2$) \\
        $t_{\rm ref,k}$     && time at which spot reaches maximum size \\
        $\alpha_{\rm max,k}$ && maximum angular size of spot \\
        $f_{\rm spot,k}$    && spot-star contrast (0 = black spot, 1 = no contrast) \\
    \hline
    \caption{Individual spot parameters and their definitions.}
    \label{tab:spot_parameters}
    \end{longtable} 
\end{center}

Figure~\ref{fig:lightcurve_demo} shows an example lightcurve generated from the starspot model, with contributions from each individual spot shown in color and the cumulative flux deficit shown in black. 

% ------------------------------------------------------------------------
\begin{figure}[ht!]
\begin{center}
    \script{plot_lightcurve_animation.py}
    \includegraphics[width=\linewidth]{figures/starspot_animation_solar_dr_frame.pdf}
    \includegraphics[width=\linewidth]{figures/starspot_animation_anti_dr_frame.pdf}
    \caption{Examples of simulated lightcurves generated from the starspot model from the \citet{kipping_analytic_2012} model. The top panel shows the case with solar-like differential rotation ($\kappa=0.5$), while the bottom panel shows the case with anti-solar differential rotation ($\kappa=-0.5$). The contributions from each individual spot are shown in color, while the cumulative flux deficit is shown in black. The lightcurve exhibits quasi-periodic variability due to the rotation of the star and the emergence and decay of spots.}
    \label{fig:lightcurve_demo}
\end{center}
\end{figure}

\clearpage
% \vspace{\baselineskip}
% ========================================================================================
\section{Gaussian Process} \label{sec:gp}

\subsection{Motivation for Gaussian Processes}

In practice, extracting the physical parameters of a star from its lightcurve is challenging due to the complex and stochastic nature of starspot-induced variability. The observed lightcurve is a superposition of signals from multiple spots, each emerging at different times and locations on the stellar surface. Even with very high quality data, performing model inversion to infer the underlying spot properties is inherently challenging for at least several reasons:
\begin{itemize}
    \item \textbf{Degeneracies:} Multiple combinations of spot properties can produce similar lightcurves, leading to degeneracies in the inferred parameters (Figure~\ref{fig:degeneracies}).
    \item \textbf{Non-uniqueness:} The same lightcurve can be explained by different spot configurations, making it difficult to uniquely determine the spot properties (Figure~\ref{fig:nonuniqueness}). 
    \item \textbf{Noise and systematics:} Real lightcurves are affected by instrumental noise, stellar granulation, and other sources of variability that can obscure the starspot signal and further complicate the inference of spot properties.
\end{itemize}

\begin{figure}[ht!]
    \centering
    \script{plot_degeneracies.py}
    \includegraphics[width=\linewidth]{figures/fig_degeneracies.pdf}
    \caption{Illustration of degeneracies in starspot lightcurve modeling.}
    \label{fig:degeneracies}
\end{figure}

\begin{figure}[h!]
    \centering
    \script{plot_nonuniqueness.py}
    \includegraphics[width=\linewidth]{figures/fig_nonuniqueness.pdf}
    \caption{Illustration of non-uniqueness in starspot lightcurve modeling. All lightcurves have the same parameters: $P_{\rm eq} = 5$\,d, $\kappa = 0$, $I = 80^\circ$, $N_{\rm spot} = 10$, $l_{\rm spot} = 8$\,d, $\tau_{\rm spot} = 3$\,d, yet the resulting lightcurves differ due to the random positions and emergence times of individual spots.}
    \label{fig:nonuniqueness}
\end{figure}

Given these challenges, a more practical approach is to not just model the properties of \emph{individual} spots, but to instead model the \emph{statistical properties} of the starspot population as a whole. This is where Gaussian processes (GPs) come in: they provide a flexible and powerful framework for modeling the covariance structure of the lightcurve without needing to explicitly model each individual spot. By deriving a physically motivated kernel function that captures the key features of starspot variability, we can use GPs to infer the underlying physical parameters of the star and its spots from the observed lightcurve, while marginalizing over the complex and degenerate spot configurations. This allows us to extract meaningful information about the star's rotation period, differential rotation, spot lifetimes, and emergence/decay timescales as a distribution over the entire spot population, rather than trying to pin down the properties of each individual spot.

% ------------------------------------------------------------------------
\subsection{Mathematical Formulation of Gaussian Processes}
Formally, a Gaussian process is defined by a multivariate Gaussian:
\begin{equation}
    \mathcal{GP}(\mathbf{x}) \sim \mathcal{N}(m(\mathbf{x}), k(\mathbf{x}, \mathbf{x}'))
\end{equation}
where the mean function, $m(\mathbf{x})$, represents the expected value of the function at each input point $\mathbf{x}$, while the covariance function, $k(\mathbf{x}, \mathbf{x}')$ models the correlation between function values at different input points $\mathbf{x}$ and $\mathbf{x}'$ and quantifies the smoothness and spatial properties of the underlying function. 

This GP, defined by a mean function and a covariance function, represents a probability distribution over functions and can be used to model the lightcurve data. Given some lightcurve data $(\mathbf{x}, \mathbf{y})$, the log marginal likelihood (or evidence) of a GP is given by:
\begin{equation} \label{eqn:marginal}
    \log P(\textbf{y} | \textbf{x}, \theta) = 
        -\frac{1}{2} \textbf{y}^T K^{-1} \textbf{y}
        -\frac{1}{2} \log \vert K \vert
        -\frac{n}{2} \log 2\pi
\end{equation}

This distribution allows us to make predictions at new test points, where the  predictive mean and variance at a given test point $x^*$ can be computed as:
\begin{align}
    \mu &= k(x_*,\mathbf{x}) k(\mathbf{x},\mathbf{x})^{-1} \mathbf{y} \\
    \sigma^2 &= k(x_*,x_*) - k(x_*,\mathbf{x}) k(\mathbf{x},\mathbf{x})^{-1} k(\mathbf{x},x_*)
\end{align}

For this particular application, the mean function $m(t)$ is determined by the expected flux deficit from the average spot coverage. In practice, we subtract the mean from the observed data and set $m(t) = 0$. The covariance function $k(\mathbf{x}, \mathbf{x}')$ is derived in Section~\ref{sec:kernel_derivation}, and encodes the physical properties of the star and its spots. By maximizing the log marginal likelihood with respect to the hyperparameters $\theta$ of the kernel, we can infer the underlying physical parameters of the star and its spots from the observed lightcurve data. For simplification we assume that the distribution of starspot properties does not evolve with time. In other words, the kernel is \emph{stationary} and depends only the time lag $\tau$ between two points.
With these basic assumptions, we move on to how to compute the kernel function numerical and analytically. However, for a more comprehensive introduction to Gaussian processes and their properties, see \citet{rasmussen_gaussian_2006}. 

% ------------------------------------------------------------------------
\subsection{The Covariance Matrix and Kernel Function} \label{sec:covariance_matrix_numerical} 

A brute-force approach to modeling the covariance structure of the lightcurve is to directly compute the covariance matrix $K$ from a large sample of simulated lightcurves generated from the starspot model.
Assuming we have 1-dimensional training data, where $\mathbf{x} = {t_1,\dots t_n}$ represent the times and $\mathbf{y} = {F_1,\dots F_n}$ represent the fluxes of the lightcurve, then we want to solve for covariance matrix:
\begin{equation}
    K = 
      \left[ {\begin{array}{ccc}
        {\rm Var}(F_1)  & \cdots & {\rm Cov}(F_1,F_n) \\
        \vdots  & \ddots & \vdots\\
        {\rm Cov}(F_n,F_1)  & \cdots & {\rm Var}(F_n) \\
      \end{array} } \right]
\end{equation}
Each element of the covariance matrix $K$ is the equivalent to the covariance between $F_i$ and $F_j$
\begin{align}
    K_{ij} = {\rm Cov}(F_i, F_j) 
\end{align}
The first term is the expectation of the flux product at two different times $t$ and $t+\tau$. Flux at a given time is given by the \citet{kipping_analytic_2012} lightcurve model (Eqn.~\ref{eqn:flux_model}). The random variables of this process are the reference time, latitude, and longitude of each spot ($t_{\rm ref,k}$, $\Lambda_{\rm ref,k}$, and $\Phi_{\rm ref,k}$ respectively).

The covariance matrix can then be estimated over a large sample of simulations with  random variables (e.g., a uniform distribution with $t_{\rm ref,k} \in [0, T]$, $\Lambda_{\rm ref,k} \in [0, 2\pi]$, and $\Phi_{\rm ref,k} \in [-\pi/2, \pi/2]$). The covariance can then be estimated as the sample covariance of the fluxes at two different times $t$ and $t+\tau$:
\begin{align}
    \mathcal{K}_{\rm num}(\tau) 
        &= \mathrm{Cov}[F(t), F(t+\tau)]  \\
        &= \frac{1}{N - 1} \sum_{i=1}^N  (F(t_i) - \mathbb{E}[F]) (F(t_i+\tau )- \mathbb{E}[F]).
    \label{eq:numerical_kernel}
\end{align}
In practice, the numerical approach is computationally expensive, since it requires generating a large number of lightcurves and computing the covariance for each pair of time points. Moreover, it does not provide a closed-form expression for the kernel function, which limits its interpretability and makes it difficult to use in a Gaussian process regression framework.
In the following sections, we derive an analytical expression for the kernel $\mathcal{K}(\tau)$ (Section~\ref{sec:kernel_derivation}) and use the numerical estimate for validation.

% \vspace{\baselineskip}
\clearpage
% ========================================================================================
\section{Deriving the Analytical Kernel for Starspot Variability} \label{sec:kernel_derivation}

% ------------------------------------------------------------------------
The goal of this section is to derive an analytical expression for the kernel function $\mathcal{K}(\tau)$, which describes the covariance of the flux at two different times $t$ and $t+\tau$ due to starspot variability. We first start with the analytic flux model defined in Equation~\ref{eqn:flux_model} and then derive the kernel function by computing the autocorrelation of the flux signal. In this section we will show the derivation for the simplified case assuming the small spot approximation and no limb darkening, however the derivation can be extended to these more complex cases, which we derive in Paper II.

The result is a physically motivated kernel function that captures the key features of starspot-induced variability, such as the rotation period, spot lifetime, and emergence/decay timescales. The kernel function depends on the underlying physical properties of the star and its spots, which are encapsulated in the hyperparameters:
\[
\theta = \{P_{\rm eq}, \kappa, I, \dot{N}_{\rm spot}, l_{\rm spot}, \tau_{\rm spot}, \alpha_{\rm max} \}.
\]
This means that the kernel can be used in a Gaussian process regression framework to model and predict stellar variability, while also allowing for inference on the physical properties of the star and its spots through the hyperparameters. In the following sections, we give an overview of the Fourier transform notation and identities used in the derivation, and then we present the step-by-step derivation of the kernel function $\mathcal{K}(\tau)$ from the flux model. 

% ------------------------------------------------------------------------
\subsection{Fourier Transform Notation and Identities}

We use the Fourier transform convention $\FT{f}(\omega) = \int_{-\infty}^{\infty} f(t)\,e^{-i\omega t}\,dt$ with inverse $f(t) = (2\pi)^{-1}\int \FT{f}(\omega)\,e^{i\omega t}\,d\omega$ throughout. Table~\ref{tab:notation_fourier} summarizes the notation.

\begin{table}[ht!]
\centering
\caption{Notation used in Section~\ref{sec:kernel_derivation}.}
\label{tab:notation_fourier}
\begin{tabular}{cl}
    \hline
    \textbf{Symbol} & \textbf{Description} \\
    \hline
    $t$ & Time \\
    $\tau$ & Time lag \\
    $\omega$ & Angular frequency \\
    $\omega_0$ & Fundamental angular frequency, $2\pi/T$ \\
    $T$ & Period of a periodic function \\
    $n$ & Harmonic index (integer) \\
    $f(t),\, g(t)$ & Functions in the time domain \\
    $\FT{f}(\omega),\, \FT{g}(\omega)$ & Fourier transforms of $f$ and $g$ \\
    $c_n$ & Fourier series coefficients \\
    $\delta(\cdot)$ & Dirac delta function \\
    $(f \ast g)(t)$ & Convolution of $f$ and $g$ \\
    $R_f(\tau)$ & Autocorrelation function of $f$ \\
    $\varepsilon(\omega)$ & Energy spectral density, $|\FT{f}(\omega)|^2$ \\
    $k(\tau)$ & Autocovariance kernel \\
    $A(t)$ & Single-spot flux signal \\
    % $\mathrm{rect}_W(t)$ & Rectangular pulse of width $W$ \\
    % $\mathrm{sinc}(x)$ & Sinc function, $\sin(x)/x$ \\
    % $\mathrm{sgn}(\cdot)$ & Sign function \\
    % $a,\, b$ & Arbitrary constants \\
    $t_0$ & Fixed time shift \\
    \hline
\end{tabular}
\end{table}


The derivation relies on three standard results. First, a periodic function with period $T$ and Fourier coefficients $c_n = T^{-1}\int_0^T f(t)\,e^{-in\omega_0 t}\,dt$ (where $\omega_0 = 2\pi/T$) has a Fourier transform given by the weighted Dirac comb $\FT{f}(\omega) = 2\pi \sum_n c_n\,\delta(\omega - n\omega_0)$.
Second, the convolution theorem: convolution in the time domain corresponds to multiplication in the frequency domain (\ref{ft:convolution}), and a product in the time domain becomes a convolution in the frequency domain (\ref{ft:multiplication}).
Third, the Wiener--Khinchin theorem: the autocorrelation $R_f(\tau) = \int f(t)\,f(t+\tau)\,dt$ is the inverse Fourier transform of the power spectral density $S_f(\omega) = |\FT{f}(\omega)|^2$:
\begin{equation}
    R_f(\tau) = \frac{1}{2\pi}\int_{-\infty}^{\infty} |\FT{f}(\omega)|^2\, e^{i\omega\tau}\, d\omega.
    \label{eq:wk}
\end{equation}
Together, these allow us to obtain the autocovariance kernel $\mathcal{K}(\tau)$ from the Fourier transform of the single-spot signal $A(t)$. In the Appendix, Table~\ref{tab:fourier_properties} lists the Fourier transform properties and Table~\ref{tab:fourier_pairs} lists the specific transform pairs used in this paper.


% ------------------------------------------------------------------------
\subsection{Approach to Derivation}

By Eq.~\eqref{eq:A_smallspot} and the flux model (Eq.~\ref{eqn:flux_model}), a single dark spot dims the star by $A(t)/\pi = \alpha_{\max}^2\,\Gamma(t)\,\mathcal{V}(t)$. The constant $\alpha_{\max}^2$ passes unchanged through every Fourier transform below, so in this section we write $A(t) \equiv \Gamma(t)\,\mathcal{V}(t)$ for the unit-amplitude single-spot signal; the kernel $\mathcal{K}(\tau)$ of the flux is $\alpha_{\max}^4$ times its autocorrelation. The autocorrelation function of function $A(t)$ is equivalent to taking the convolution between $A(t)$ and a lagged version of the function. By the convolution property (\ref{ft:convolution}), convolution in the time domain is multiplication in the frequency domain. Here $\mathcal{F}\{A(t)\} = \hat{A}(\omega)$ denotes the Fourier transform of a function $A(t)$, and $\mathcal{F}^{-1}\{\hat{A}(\omega)\} = A(t)$ denotes the inverse Fourier transform.
So the Fourier transform of the kernel function is:
\begin{align}
    \mathcal{F}\{\mathcal{K}(\tau) \}
        &=  \alpha_{\max}^4\, A(t) \ast A(t - \tau) \\
        &=  \alpha_{\max}^4\, \hat{A}(\omega) \cdot \hat{A}(\omega) \\
        &= \alpha_{\max}^4\, |\hat{A}(\omega)|^2 \\
        & = \alpha_{\max}^4\, S_A(\omega)
\end{align}
The term $S_A(\omega)$ is the \emph{power spectral density} of the single-spot signal. By the Wiener--Khinchin theorem (Eq.~\ref{eq:wk}), the autocorrelation function is the inverse Fourier transform of the power spectral density:
\begin{equation}
    \mathcal{K}(\tau) \overset{\mathcal{F}} \longleftrightarrow \alpha_{\max}^4\, S_A(\omega)
\end{equation}
\begin{equation}
    \mathcal{K}(\tau) = \frac{\alpha_{\max}^4}{2\pi} \int_{-\infty}^\infty S_A(\omega) e^{i\omega t} \, d\omega .
\end{equation}
Hence to solve for $\mathcal{K}(\tau)$, we need to first solve for the Fourier transform of a single spot, $\hat{A}(\omega) = \mathcal{F}\{ A(t) \}$. 
Referring back to Equation~\ref{eq:A_smallspot}, ignoring the edge cases (for small spot sizes), the unit-amplitude signal $A(t)$ is the product of the normalized squared spot-size envelope $\Gamma(t) = \alpha^2(t)/\alpha_{\max}^2$ and the visibility function $\mathcal{V}(t) = \mathrm{max}\{\cos\beta(t), 0 \}$. The visibility function is a truncated cosine, which can be written as the product of a square wave (\ref{fp:sqwave}) and cosine wave (\ref{fp:cosine}): $\mathcal{V}(t) = \sqcap(t) \cdot C(t)$.
By the multiplication property (\ref{ft:multiplication}), a product in the time domain becomes a convolution in the frequency domain:
\begin{align}
    \mathcal{F}\{ A(t) \} &= \mathcal{F}\{ \Gamma(t) \cdot \mathcal{V}(t) \} \\
        &= \mathcal{F}\{ \Gamma(t) \} \ast \mathcal{F}\{ \mathcal{V}(t)  \} \\
        \hat{A}(\omega) &= \hat{\Gamma}(\omega) \ast \hat{\mathcal{V}}(\omega)
        \label{eqn:ft_At}
\end{align}

As we show in the next section, the Fourier transform of the spot evolution $\hat{A}(\omega)$ can be solved analytically for both of the components $\hat{\Gamma}(\omega)$ and $\hat{\mathcal{V}}(\omega)$. In Section \ref{subsec:ft_cosine} we derive the Fourier transform of the truncated cosine $\hat{\mathcal{V}}(\omega)$ (representing the variability as the star rotates), and in Section \ref{subsec:ft_envelope} we derive the Fourier transform of the squared spot-size envelope $\hat{\Gamma}(\omega)$ (representing the growth and decay of the spot). Finally, in Section \ref{subsec:kernel_full} we combine these results to obtain the final expression for the kernel $\mathcal{K}(\tau)$.

% Overview diagram: analytic vs numerical kernel derivation
% Include in ms.tex with % Overview diagram: analytic vs numerical kernel derivation
% Include in ms.tex with % Overview diagram: analytic vs numerical kernel derivation
% Include in ms.tex with \input{overview}
% Requires: \usepackage{tikz} in the preamble

\usetikzlibrary{arrows.meta, positioning, calc, backgrounds, fit, shapes.geometric}

\begin{figure*}[ht!]
\centering
\begin{tikzpicture}[
    % Node styles
    box/.style={
        rectangle, rounded corners=4pt, draw=#1, fill=#1!8,
        line width=1.2pt, text width=3.0cm, minimum height=1.6cm,
        align=center, font=\small
    },
    box/.default=black,
    analytic/.style={box=blue!70!black},
    numerical/.style={box=red!70!black},
    result/.style={
        rectangle, rounded corners=6pt, draw=black!70!green, fill=green!8,
        line width=1.8pt, text width=3.6cm, minimum height=1.8cm,
        align=center, font=\small\bfseries
    },
    verify/.style={
        diamond, draw=black!60!green, fill=green!15,
        line width=1.4pt, text width=1.8cm,
        align=center, font=\small\bfseries, inner sep=2pt,
        aspect=1.8
    },
    arr/.style={
        -{Stealth[length=3mm, width=2mm]}, line width=1.0pt, color=#1
    },
    arr/.default=black!60,
    label/.style={font=\normalsize, text=black!60, align=center},
    rowlabel/.style={font=\small\bfseries\scshape, rotate=90, anchor=center},
]

% ═══════════════════════════════════════════════════════════════
% Row 1: Analytic derivation (top)
% ═══════════════════════════════════════════════════════════════

\node[analytic] (A1) {
    Define flux model\\[2pt]
    $A(t) = \Gamma(t) \cdot \mathcal{V}(t)$
};

\node[analytic, right=1.1cm of A1] (A2) {
    Fourier transforms\\[2pt]
    $\hat{A}(\omega) = \hat{\Gamma} \ast \hat{\mathcal{V}}$
};

\node[analytic, right=1.1cm of A2] (A3) {
    Average over\\spot parameters\\[2pt]
    $\langle|\hat{A}|^2\rangle = \sum_n |c_n|^2 S_\Gamma$
};

\node[result, right=1.1cm of A3] (A4) {
    Analytic kernel\\[2pt]
    $\scriptstyle k(\tau) = \alpha_{\max}^4 R_\Gamma \!\sum_n |c_n|^2 \cos(n\omega_0\tau)$
};

% Arrows (analytic row)
\draw[arr=blue!70!black] (A1) -- (A2);
\draw[arr=blue!70!black] (A2) -- (A3);
\draw[arr=blue!70!black] (A3) -- (A4);

% Step labels
\node[label, above=0.15cm of A1] {\textit{Step 1}};
\node[label, above=0.15cm of A2] {\textit{Step 2}};
\node[label, above=0.15cm of A3] {\textit{Step 3}};
\node[label, above=0.15cm of A4] {\textit{Step 4}};

% Row label (centered vertically with the row)
\node[rowlabel, blue!70!black] at ($(A1.west)+(-0.9cm,0)$) {Analytic};

% ═══════════════════════════════════════════════════════════════
% Row 2: Numerical verification (bottom)
% ═══════════════════════════════════════════════════════════════

\node[numerical, below=2.8cm of A1] (N1) {
    Flux model\\[2pt]
    $F(t) = 1 - \frac{1}{\pi}\sum_k A_k(t)$
};

\node[numerical, right=1.1cm of N1, text width=3.4cm] (N2) {
    Simulate $N$ lightcurves\\[2pt]
    random $t_{\rm ref}$, $\Lambda_{\rm ref}$, $\Phi_{\rm ref}$
};

\node[result, right=1.1cm of N2, fill=red!8, draw=red!70!black] (N3) {
    Numerical kernel\\[2pt]
    $k_{\rm num}(\tau) = \frac{1}{N}\sum_i \mathrm{ACF}_i(\tau)$
};

% Arrows (numerical row)
\draw[arr=red!70!black] (N1) -- (N2);
\draw[arr=red!70!black] (N2) -- (N3);

% Step labels
\node[label, above=0.15cm of N1] {\textit{Step 1}};
\node[label, above=0.15cm of N2] {\textit{Step 2}};
\node[label, above=0.15cm of N3] {\textit{Step 3}};

% Row label (centered vertically with the row)
\node[rowlabel, red!70!black] at ($(N1.west)+(-0.9cm,0)$) {Numerical};

% ═══════════════════════════════════════════════════════════════
% Convergence / verification
% ═══════════════════════════════════════════════════════════════

% Place the verification node between analytic result and numerical result
\node[verify] (V) at ($(A4)!0.5!(N3)$) {
    Verify\\$k \approx k_{\rm num}$
};

\draw[arr=black!60!green, line width=1.2pt] (A4) -- (V);
\draw[arr=black!60!green, line width=1.2pt] (N3) -- (V);

% Shared starting point arrow
\draw[arr=black!40, dashed, line width=0.8pt]
    (A1.south) -- ++(0, -0.7cm) -| (N1.north);

\node[label, font=\scriptsize\itshape] at ($(A1.south)!0.5!(N1.north) + (-1.1cm, 0)$) {same model};

\end{tikzpicture}
\caption{Overview of the kernel derivation strategy. \textit{Top row (blue):} The analytic path derives $k(\tau)$ from the single-spot flux model $A(t) = \Gamma(t)\cdot\mathcal{V}(t)$ by (1) computing Fourier transforms of the envelope and visibility terms, (2) averaging the energy spectral density over random spot parameters ($t_{\rm ref}$, $\Lambda_{\rm ref}$, $\Phi_{\rm ref}$), and (3) applying the Wiener--Khinchin theorem. \textit{Bottom row (red):} The numerical path simulates $N$ lightcurves from the same flux model with randomly drawn spot parameters and estimates the kernel as the sample-averaged autocorrelation. \textit{Center (green):} The two paths converge, verifying the analytic kernel against the numerical estimate.}
\label{fig:overview}
\end{figure*}

% Requires: \usepackage{tikz} in the preamble

\usetikzlibrary{arrows.meta, positioning, calc, backgrounds, fit, shapes.geometric}

\begin{figure*}[ht!]
\centering
\begin{tikzpicture}[
    % Node styles
    box/.style={
        rectangle, rounded corners=4pt, draw=#1, fill=#1!8,
        line width=1.2pt, text width=3.0cm, minimum height=1.6cm,
        align=center, font=\small
    },
    box/.default=black,
    analytic/.style={box=blue!70!black},
    numerical/.style={box=red!70!black},
    result/.style={
        rectangle, rounded corners=6pt, draw=black!70!green, fill=green!8,
        line width=1.8pt, text width=3.6cm, minimum height=1.8cm,
        align=center, font=\small\bfseries
    },
    verify/.style={
        diamond, draw=black!60!green, fill=green!15,
        line width=1.4pt, text width=1.8cm,
        align=center, font=\small\bfseries, inner sep=2pt,
        aspect=1.8
    },
    arr/.style={
        -{Stealth[length=3mm, width=2mm]}, line width=1.0pt, color=#1
    },
    arr/.default=black!60,
    label/.style={font=\normalsize, text=black!60, align=center},
    rowlabel/.style={font=\small\bfseries\scshape, rotate=90, anchor=center},
]

% ═══════════════════════════════════════════════════════════════
% Row 1: Analytic derivation (top)
% ═══════════════════════════════════════════════════════════════

\node[analytic] (A1) {
    Define flux model\\[2pt]
    $A(t) = \Gamma(t) \cdot \mathcal{V}(t)$
};

\node[analytic, right=1.1cm of A1] (A2) {
    Fourier transforms\\[2pt]
    $\hat{A}(\omega) = \hat{\Gamma} \ast \hat{\mathcal{V}}$
};

\node[analytic, right=1.1cm of A2] (A3) {
    Average over\\spot parameters\\[2pt]
    $\langle|\hat{A}|^2\rangle = \sum_n |c_n|^2 S_\Gamma$
};

\node[result, right=1.1cm of A3] (A4) {
    Analytic kernel\\[2pt]
    $\scriptstyle k(\tau) = \alpha_{\max}^4 R_\Gamma \!\sum_n |c_n|^2 \cos(n\omega_0\tau)$
};

% Arrows (analytic row)
\draw[arr=blue!70!black] (A1) -- (A2);
\draw[arr=blue!70!black] (A2) -- (A3);
\draw[arr=blue!70!black] (A3) -- (A4);

% Step labels
\node[label, above=0.15cm of A1] {\textit{Step 1}};
\node[label, above=0.15cm of A2] {\textit{Step 2}};
\node[label, above=0.15cm of A3] {\textit{Step 3}};
\node[label, above=0.15cm of A4] {\textit{Step 4}};

% Row label (centered vertically with the row)
\node[rowlabel, blue!70!black] at ($(A1.west)+(-0.9cm,0)$) {Analytic};

% ═══════════════════════════════════════════════════════════════
% Row 2: Numerical verification (bottom)
% ═══════════════════════════════════════════════════════════════

\node[numerical, below=2.8cm of A1] (N1) {
    Flux model\\[2pt]
    $F(t) = 1 - \frac{1}{\pi}\sum_k A_k(t)$
};

\node[numerical, right=1.1cm of N1, text width=3.4cm] (N2) {
    Simulate $N$ lightcurves\\[2pt]
    random $t_{\rm ref}$, $\Lambda_{\rm ref}$, $\Phi_{\rm ref}$
};

\node[result, right=1.1cm of N2, fill=red!8, draw=red!70!black] (N3) {
    Numerical kernel\\[2pt]
    $k_{\rm num}(\tau) = \frac{1}{N}\sum_i \mathrm{ACF}_i(\tau)$
};

% Arrows (numerical row)
\draw[arr=red!70!black] (N1) -- (N2);
\draw[arr=red!70!black] (N2) -- (N3);

% Step labels
\node[label, above=0.15cm of N1] {\textit{Step 1}};
\node[label, above=0.15cm of N2] {\textit{Step 2}};
\node[label, above=0.15cm of N3] {\textit{Step 3}};

% Row label (centered vertically with the row)
\node[rowlabel, red!70!black] at ($(N1.west)+(-0.9cm,0)$) {Numerical};

% ═══════════════════════════════════════════════════════════════
% Convergence / verification
% ═══════════════════════════════════════════════════════════════

% Place the verification node between analytic result and numerical result
\node[verify] (V) at ($(A4)!0.5!(N3)$) {
    Verify\\$k \approx k_{\rm num}$
};

\draw[arr=black!60!green, line width=1.2pt] (A4) -- (V);
\draw[arr=black!60!green, line width=1.2pt] (N3) -- (V);

% Shared starting point arrow
\draw[arr=black!40, dashed, line width=0.8pt]
    (A1.south) -- ++(0, -0.7cm) -| (N1.north);

\node[label, font=\scriptsize\itshape] at ($(A1.south)!0.5!(N1.north) + (-1.1cm, 0)$) {same model};

\end{tikzpicture}
\caption{Overview of the kernel derivation strategy. \textit{Top row (blue):} The analytic path derives $k(\tau)$ from the single-spot flux model $A(t) = \Gamma(t)\cdot\mathcal{V}(t)$ by (1) computing Fourier transforms of the envelope and visibility terms, (2) averaging the energy spectral density over random spot parameters ($t_{\rm ref}$, $\Lambda_{\rm ref}$, $\Phi_{\rm ref}$), and (3) applying the Wiener--Khinchin theorem. \textit{Bottom row (red):} The numerical path simulates $N$ lightcurves from the same flux model with randomly drawn spot parameters and estimates the kernel as the sample-averaged autocorrelation. \textit{Center (green):} The two paths converge, verifying the analytic kernel against the numerical estimate.}
\label{fig:overview}
\end{figure*}

% Requires: \usepackage{tikz} in the preamble

\usetikzlibrary{arrows.meta, positioning, calc, backgrounds, fit, shapes.geometric}

\begin{figure*}[ht!]
\centering
\begin{tikzpicture}[
    % Node styles
    box/.style={
        rectangle, rounded corners=4pt, draw=#1, fill=#1!8,
        line width=1.2pt, text width=3.0cm, minimum height=1.6cm,
        align=center, font=\small
    },
    box/.default=black,
    analytic/.style={box=blue!70!black},
    numerical/.style={box=red!70!black},
    result/.style={
        rectangle, rounded corners=6pt, draw=black!70!green, fill=green!8,
        line width=1.8pt, text width=3.6cm, minimum height=1.8cm,
        align=center, font=\small\bfseries
    },
    verify/.style={
        diamond, draw=black!60!green, fill=green!15,
        line width=1.4pt, text width=1.8cm,
        align=center, font=\small\bfseries, inner sep=2pt,
        aspect=1.8
    },
    arr/.style={
        -{Stealth[length=3mm, width=2mm]}, line width=1.0pt, color=#1
    },
    arr/.default=black!60,
    label/.style={font=\normalsize, text=black!60, align=center},
    rowlabel/.style={font=\small\bfseries\scshape, rotate=90, anchor=center},
]

% ═══════════════════════════════════════════════════════════════
% Row 1: Analytic derivation (top)
% ═══════════════════════════════════════════════════════════════

\node[analytic] (A1) {
    Define flux model\\[2pt]
    $A(t) = \Gamma(t) \cdot \mathcal{V}(t)$
};

\node[analytic, right=1.1cm of A1] (A2) {
    Fourier transforms\\[2pt]
    $\hat{A}(\omega) = \hat{\Gamma} \ast \hat{\mathcal{V}}$
};

\node[analytic, right=1.1cm of A2] (A3) {
    Average over\\spot parameters\\[2pt]
    $\langle|\hat{A}|^2\rangle = \sum_n |c_n|^2 S_\Gamma$
};

\node[result, right=1.1cm of A3] (A4) {
    Analytic kernel\\[2pt]
    $\scriptstyle k(\tau) = \alpha_{\max}^4 R_\Gamma \!\sum_n |c_n|^2 \cos(n\omega_0\tau)$
};

% Arrows (analytic row)
\draw[arr=blue!70!black] (A1) -- (A2);
\draw[arr=blue!70!black] (A2) -- (A3);
\draw[arr=blue!70!black] (A3) -- (A4);

% Step labels
\node[label, above=0.15cm of A1] {\textit{Step 1}};
\node[label, above=0.15cm of A2] {\textit{Step 2}};
\node[label, above=0.15cm of A3] {\textit{Step 3}};
\node[label, above=0.15cm of A4] {\textit{Step 4}};

% Row label (centered vertically with the row)
\node[rowlabel, blue!70!black] at ($(A1.west)+(-0.9cm,0)$) {Analytic};

% ═══════════════════════════════════════════════════════════════
% Row 2: Numerical verification (bottom)
% ═══════════════════════════════════════════════════════════════

\node[numerical, below=2.8cm of A1] (N1) {
    Flux model\\[2pt]
    $F(t) = 1 - \frac{1}{\pi}\sum_k A_k(t)$
};

\node[numerical, right=1.1cm of N1, text width=3.4cm] (N2) {
    Simulate $N$ lightcurves\\[2pt]
    random $t_{\rm ref}$, $\Lambda_{\rm ref}$, $\Phi_{\rm ref}$
};

\node[result, right=1.1cm of N2, fill=red!8, draw=red!70!black] (N3) {
    Numerical kernel\\[2pt]
    $k_{\rm num}(\tau) = \frac{1}{N}\sum_i \mathrm{ACF}_i(\tau)$
};

% Arrows (numerical row)
\draw[arr=red!70!black] (N1) -- (N2);
\draw[arr=red!70!black] (N2) -- (N3);

% Step labels
\node[label, above=0.15cm of N1] {\textit{Step 1}};
\node[label, above=0.15cm of N2] {\textit{Step 2}};
\node[label, above=0.15cm of N3] {\textit{Step 3}};

% Row label (centered vertically with the row)
\node[rowlabel, red!70!black] at ($(N1.west)+(-0.9cm,0)$) {Numerical};

% ═══════════════════════════════════════════════════════════════
% Convergence / verification
% ═══════════════════════════════════════════════════════════════

% Place the verification node between analytic result and numerical result
\node[verify] (V) at ($(A4)!0.5!(N3)$) {
    Verify\\$k \approx k_{\rm num}$
};

\draw[arr=black!60!green, line width=1.2pt] (A4) -- (V);
\draw[arr=black!60!green, line width=1.2pt] (N3) -- (V);

% Shared starting point arrow
\draw[arr=black!40, dashed, line width=0.8pt]
    (A1.south) -- ++(0, -0.7cm) -| (N1.north);

\node[label, font=\scriptsize\itshape] at ($(A1.south)!0.5!(N1.north) + (-1.1cm, 0)$) {same model};

\end{tikzpicture}
\caption{Overview of the kernel derivation strategy. \textit{Top row (blue):} The analytic path derives $k(\tau)$ from the single-spot flux model $A(t) = \Gamma(t)\cdot\mathcal{V}(t)$ by (1) computing Fourier transforms of the envelope and visibility terms, (2) averaging the energy spectral density over random spot parameters ($t_{\rm ref}$, $\Lambda_{\rm ref}$, $\Phi_{\rm ref}$), and (3) applying the Wiener--Khinchin theorem. \textit{Bottom row (red):} The numerical path simulates $N$ lightcurves from the same flux model with randomly drawn spot parameters and estimates the kernel as the sample-averaged autocorrelation. \textit{Center (green):} The two paths converge, verifying the analytic kernel against the numerical estimate.}
\label{fig:overview}
\end{figure*}


Figure~\ref{fig:overview} provides an overview of the derivation strategy. The analytic path (top row) derives $\mathcal{K}(\tau)$ through Fourier analysis, while a parallel numerical path (bottom row) estimates the kernel via Monte Carlo simulation. Finally, we compare the numerical and analytical kernels as a validation test of the derivation.

% -----------------------------------------------
\subsection{Fourier Transform: The Rotation Term $\mathcal{V}(t)$} \label{subsec:ft_cosine}

As given in Equation~\ref{eq:cos_beta}, the drop in flux due to a spot rotating into and out of view is determined by the angle $\beta$ between the spot normal and the line of sight.\footnote{Symbolic and numerical check: \NotebookLink{ft_spot_rotation.ipynb}} The spot is visible when $\cos\beta > 0$, and the visibility function is $\mathcal{V}(t) = \max\{\cos\beta(t),\,0\}$. The angle $\beta$ depends on the spot latitude $\Phi$, the stellar inclination $I$, and the rotation period $P$ of the star. 

Figure~\ref{fig:visibility} shows how the signal varies as a function of spot latitude $\Phi$ and stellar inclination $I$. For a star viewed equator-on ($I = 90^\circ$, left panel of Figure~\ref{fig:visibility}), the visibility function looks like a half-truncated cosine wave, where the spot rotates into view for half a rotation period and then rotates out of view for the other half. The different colors show the relative change in flux for spots at different latitudes $\Phi$, with spots at higher latitudes having smaller amplitudes due to their smaller projected areas.

Figure~\ref{fig:visibility} also shows how the signal varies as a function of stellar inclination $I$ for a fixed spot latitude $\Phi = 60^\circ$ (right panel). For a star viewed equator-on ($I = 90^\circ$), the visibility function is a half-truncated cosine wave (red). For a star viewed at intermediate inclination ($60^\circ < I < 90^\circ$), the visibility function is a semi-truncated cosine wave (blue). For a star viewed closer to pole-on ($I = 30^\circ$), the visibility function is a full cosine wave which never rotates out of view, but varies in amplitude due to the changing projected area of the spot (orange). Finally, for a star viewed pole-on ($I = 0^\circ$), the visibility function is a constant (blue) since the spot is always visible with the same projected area.
\begin{figure}[ht!]
    \centering
    \script{plot_spot_visibility.py}
    \includegraphics[width=\linewidth]{figures/fig_visibility.pdf}
    \caption{Visibility function $\mathcal{V}(t) = \max\{\cos\beta(t),\,0\}$ over three rotation periods.
    \textit{Left:} varying spot latitude $\Phi$ at fixed edge-on inclination $I=90^\circ$.
    \textit{Right:} varying stellar inclination $I$ at fixed latitude $\Phi=60^\circ$.
    The amplitude and duty cycle of $\mathcal{V}(t)$ set the Fourier coefficients $|c_n|^2$ that enter the kernel.}
    \label{fig:visibility}
\end{figure}

The visibility function $\mathcal{V}(t) = \max\{\cos\beta(t),\,0\}$ is periodic with the rotation period $P = 2\pi/\omega_0$, so it has a discrete Fourier series (Section~\ref{sec:kernel_derivation}):
\begin{equation}
    \mathcal{V}(t) = \sum_{n=-\infty}^{\infty} c_n\, e^{in\omega_0 t},
    \label{eq:cap_fourier_series}
\end{equation}
and its Fourier transform is a summation over Dirac functions (\ref{fp:exp}):
\begin{equation}
    \FT{\mathcal{V}}(\omega) = 2\pi \sum_{n=-\infty}^{\infty} c_n\, \delta(\omega - n\omega_0).
    \label{eq:cap_hat}
\end{equation}

\textbf{Solution for the edge-on case}: For the edge-on case ($I = 90^\circ$), $\cos\beta(t) = \cos\Phi\cos(\omega_0 t)$, and the visibility function reduces to $\mathcal{V}(t) = \cos\Phi \cdot \max\{\cos(\omega_0 t),\, 0\}$. The Fourier coefficients of the half-rectified cosine $\max\{\cos\theta,\,0\}$ are (see Appendix~\ref{app:fourier_visibility} for the full derivation):
\begin{equation}
    c_n =
    \cos\Phi
    \begin{cases}
        1 / \pi & n = 0 \\[4pt]
        1 / 4 & n = \pm 1 \\[4pt]
        \displaystyle \frac{(-1)^{|n|/2+1}}{\pi(n^2 - 1)} & n = \pm 2, \pm 4, \pm 6, \ldots \\[6pt]
        0 & n = \pm 3, \pm 5, \pm 7, \ldots
    \end{cases}
    \label{eq:cn_edgeon}
\end{equation}
All coefficients are real and satisfy $c_{-n} = c_n$, as expected for a real, even function.

\textbf{Solution for general inclinations}: For general inclination, $\cos\beta(t) = b_0 + b_1\cos(\omega_0 t)$ where $b_0 = \cos I\sin\Phi$ and $b_1 = \sin I\cos\Phi$. The spot is visible when $|\omega_0 t| < \theta_{\rm vis} = \arccos(-b_0/b_1)$, and the general Fourier coefficients (derived in Appendix~\ref{app:fourier_visibility}) are:
\begin{equation}
    \boxed{
    c_n = \frac{1}{2\pi}\int_{-\theta_{\rm vis}}^{\theta_{\rm vis}} \big(b_0 + b_1\cos\theta\big)\, e^{-in\theta}\, d\theta
    = \frac{1}{\pi}\left[\frac{b_0\sin(n\,\theta_{\rm vis})}{n} + \frac{b_1}{2}\left(\frac{\sin\big((n-1)\theta_{\rm vis}\big)}{n-1} + \frac{\sin\big((n+1)\theta_{\rm vis}\big)}{n+1}\right)\right]
    }
    \label{eq:cn_general}
\end{equation}
with $c_0 = (b_0\,\theta_{\rm vis} + b_1\sin\theta_{\rm vis})/\pi$ and $c_{\pm 1} = (b_0\sin\theta_{\rm vis} + b_1[\theta_{\rm vis} + \sin\theta_{\rm vis}\cos\theta_{\rm vis}]/2)/\pi$ obtained by taking the appropriate limits. These are not separate cases but removable singularities of Eq.~\eqref{eq:cn_general}, recovered using $\sin(n\theta_{\rm vis})/n \to \theta_{\rm vis}$ as $n \to 0$ and $\sin[(n-1)\theta_{\rm vis}]/(n-1) \to \theta_{\rm vis}$ as $n \to 1$.

In general, it is difficult to solve for a closed-form solution to the integral of the full kernel function in Equation~\ref{eq:kernel_full} due to the complexity of the latitude-dependent terms. However, for specialized cases (e.g., a star with edge-on inclination, solid body rotation, and uniform latitude spots), the kernel reduces to a simple, closed-form solution. Truncating to second harmonic order, the Fourier transform of the visibility function is:
\begin{equation}
    \FT{\mathcal{V}}(\omega) \approx 2\pi\left[
        c_0\,\delta(\omega)
        + c_1\big[\delta(\omega - \omega_0) + \delta(\omega + \omega_0)\big]
        + c_2\big[\delta(\omega - 2\omega_0) + \delta(\omega + 2\omega_0)\big]
    \right].
    \label{eq:cap_hat_explicit}
\end{equation}

As we will show in Section~\ref{subsec:kernel_combined_single_spot}, the kernel $\mathcal{K}(\tau)$ is dominated by the first two harmonics, so this approximation is sufficient for our purposes (though in our \code{spotgp} implementation, this Fourier series can be expanded to an arbitrary number of terms as specified by the user).

% ------------------------------------------------------------------------
\subsection{Fourier Transform: Spot Evolution Envelope Term $\Gamma(t)$} \label{subsec:ft_envelope}

In this section we show how to derive the Fourier transform of the spot evolution envelope, which captures the growth and decay of the spot over time.\footnote{Symbolic and numerical checks: \NotebookLink{ft_squared_trapezoid.ipynb} for the squared trapezoid $\Gamma(t)$, and \NotebookLink{ft_trapezoid.ipynb} for the trapezoid $\alpha(t)$ itself.} Here we present several common spot evolution profiles, but the same approach can be applied to any arbitrary profile $\alpha(t)$ by performing the appropriate integral. These profiles include: (1) squared trapezoidal \citep{kipping_analytic_2012}, (2) exponential \citep{aigrain_testing_2015}, and (3) skewed normal \citep{forgacs-dajka_time-dependent_2021}. The trapezoidal and exponential profiles are simple functions which produce an analytic kernel function. On the other hand, the skewed normal profile is a more flexible profile that can capture asymmetric growth and decay timescales and reproduce the observed spot lifetimes of the Sun \citep{forgacs-dajka_time-dependent_2021}. 

In the Results and Discussion (Sections~\ref{sec:results}--\ref{sec:discussion}) we show a detailed analysis of how the properties of the trapezoidal profile translate to the structure of the kernel function, as its analytic solution is easily interpretable. While perhaps not the most realistic model, the trapezoidal parameterization provides an intuitive insight into how the kernel behaves in different limits. We include the implementations for all three kernels in our \code{spotgp} implementation (as well as the flexibility to define custom envelope functions) so that users can experiment with different assumptions for spot lifetimes and decay timescales.

\textbf{(1) Squared trapezoidal profile for $\Gamma(t)$:} In the first example, we consider the trapezoidal profile (Equation~\ref{eqn:trapezoid}) from \citet{kipping_analytic_2012}. Appendix~\ref{app:fourier_trapezoid}, shows how we solve for the Fourier transform of the normalized squared envelope $\Gamma(t) = \alpha^2(t)/\alpha_{\max}^2$: Since $\alpha(t)$ is a trapezoid with peak $\alpha_{\max}$, the normalized squared envelope $\Gamma(t) = \alpha^2(t)/\alpha_{\max}^2$ has unit peak height, the same support $\ell_{\rm spot} + 2\tau_{\rm spot}$, the same plateau duration $\ell_{\rm spot}$, but parabolic (rather than linear) ramps. Since $\Gamma(t)$ is real and even, its Fourier transform is purely real and can be computed via the cosine transform $\FT{\Gamma}(\omega) = 2\int_0^\infty \Gamma(t)\cos(\omega t)\,dt$, splitting into a plateau integral and a parabolic-ramp integral with two integrations by parts (see Appendix~\ref{app:fourier_trapezoid}):
\begin{equation}
    \FT{\Gamma}(\omega) = \frac{4}{\tau_{\rm spot}^2\,\omega^3}
    \left[
        \tau_{\rm spot}\,\omega\cos\!\left(\frac{\omega\,\ell_{\rm spot}}{2}\right)
        + \sin\!\left(\frac{\omega\,\ell_{\rm spot}}{2}\right)
        - \sin\!\left(\frac{\omega\,\ell_{\rm spot}}{2} + \omega\,\tau_{\rm spot}\right)
    \right].
    \label{eq:Gamma_hat}
\end{equation}
A visualization of this profile is shown in the left panel of Figure~\ref{fig:envelope_kernel_comparison}.

\textbf{(2) Exponential profile for $\Gamma(t)$:}
For the piecewise exponential envelope (Eq.~\ref{eq:alpha_exp}), setting $t_1 = 0$ and $t_2 = \tau_{\rm plat}$, the normalized squared envelope $\Gamma(t) = \alpha_{\rm exp}^2(t)/\alpha_{\max}^2$ is
\begin{equation}
    \Gamma(t) =
    \begin{cases}
        e^{2t/\tau_{\rm em}} & t < 0 \\
        1 & 0 \leq t \leq \tau_{\rm plat} \\
        e^{-2(t - \tau_{\rm plat})/\tau_{\rm dec}} & t > \tau_{\rm plat}
    \end{cases}
    \label{eq:Gamma_exp}
\end{equation}
with doubled exponential rates compared to $\alpha(t)$.
Because $\Gamma(t)$ is not symmetric when $\tau_{\rm em} \neq \tau_{\rm dec}$, the Fourier transform is complex.
Splitting $\FT{\Gamma}(\omega) = \FT{\Gamma}_R(\omega) - i\,\FT{\Gamma}_I(\omega)$ and evaluating each piece integral analytically:
\begin{equation}
    \FT{\Gamma}_R(\omega) =
        \frac{2\tau_{\rm em}}{4 + \omega^2\tau_{\rm em}^2}
        + \frac{\sin(\omega\tau_{\rm plat})}{\omega}
        + \frac{2\tau_{\rm dec}\cos(\omega\tau_{\rm plat})}{4 + \omega^2\tau_{\rm dec}^2}
        - \frac{\omega\tau_{\rm dec}^2\sin(\omega\tau_{\rm plat})}{4 + \omega^2\tau_{\rm dec}^2}
    \label{eq:Gamma_hat_exp_R}
\end{equation}
\begin{equation}
    \FT{\Gamma}_I(\omega) =
        -\frac{\omega\tau_{\rm em}^2}{4 + \omega^2\tau_{\rm em}^2}
        + \frac{1 - \cos(\omega\tau_{\rm plat})}{\omega}
        + \frac{\omega\tau_{\rm dec}^2\cos(\omega\tau_{\rm plat})}{4 + \omega^2\tau_{\rm dec}^2}
        + \frac{2\tau_{\rm dec}\sin(\omega\tau_{\rm plat})}{4 + \omega^2\tau_{\rm dec}^2}
    \label{eq:Gamma_hat_exp_I}
\end{equation}
The full derivation, including the direct convolution for $R_\Gamma(\tau)$, is analogous to the trapezoidal case (Appendix~\ref{app:fourier_trapezoid}), and a figure of the function is shown in the middle panel of Figure~\ref{fig:envelope_kernel_comparison}.

\textbf{(3) Skewed Normal profile for $\Gamma(t)$:}
\citet{forgacs-dajka_time-dependent_2021} modeled sunspot group area evolution using a skew-normal function (their Eq.~1), which naturally captures the observed asymmetry between rapid emergence and slower decay phases of sunspot lifetimes. We adopt this as a third envelope option. The normalized squared envelope is
\begin{equation}
    \Gamma(t) \;\propto\; \exp\!\left(-\frac{t^2}{2\sigma_{\rm sn}^2}\right) \left[1 + \mathrm{erf}\!\left(\frac{n_{\rm sn}\,t}{\sigma_{\rm sn}\sqrt{2}}\right)\right],
    \label{eq:Gamma_skewnorm}
\end{equation}
rescaled so that the peak value equals unity. The two free parameters are $\sigma_{\rm sn}$ (the scale parameter in days, controlling the overall spot lifetime) and $n_{\rm sn}$ (the dimensionless skewness parameter). When $n_{\rm sn} = 0$ the envelope reduces to a symmetric Gaussian; $n_{\rm sn} < 0$ gives rapid rise and slow decay (typical of solar spots), while $n_{\rm sn} > 0$ gives the opposite asymmetry.
Unlike the trapezoidal and exponential profiles, the skew-normal $\Gamma(t)$ does not yield a closed-form autocorrelation $R_\Gamma(\tau)$ or Fourier transform $\widehat{\Gamma}(\omega)$, so we therefore compute both numerically via FFT-based methods in our \code{spotgp} implementation. This function is shown in the right panel of Figure~\ref{fig:envelope_kernel_comparison}.

\begin{figure}[ht!]
\begin{center}
    \script{plot_envelope_kernel_comparison.py}
    \includegraphics[width=\linewidth]{figures/fig_envelope_kernel_comparison.pdf}
    \caption{Comparison of the three spot envelope models: trapezoidal (blue), exponential (green), and skew-normal (red). \textit{Left:} Normalized squared envelopes $\Gamma(t)$. \textit{Middle:} Energy spectral densities $S_\Gamma(\omega)$ (normalized, log scale), showing the different high-frequency decay rates set by envelope smoothness. \textit{Right:} Normalized envelope autocorrelation $R_\Gamma(\tau)/R_\Gamma(0)$. All three implementations are included in our \code{spotgp} package, along with the functionality to define custom envelope profiles.}
    \label{fig:envelope_kernel_comparison}
\end{center}
\end{figure}

% -----------------------------------------------
\subsection{Combined Kernel for a Single Spot} \label{subsec:kernel_combined_single_spot}

Substituting the delta train~\eqref{eq:cap_hat} into Eq.~\eqref{eqn:ft_At} and using the sifting property of the Dirac delta (\ref{fp:delta}, \ref{ft:freqshift}), $f(\omega) \ast \delta(\omega - a) = f(\omega - a)$:
\begin{equation}
    \FT{A}(\omega) = \sum_n c_n\, \FT{\Gamma}(\omega - n\omega_0)
    \label{eq:A_hat_sum}
\end{equation}
Explicitly to second harmonic order, using linearity (\ref{ft:linearity}):
\begin{equation}
    \boxed{
    \FT{A}(\omega) =
        c_0\,\FT{\Gamma}(\omega)
        + c_1\big[\FT{\Gamma}(\omega - \omega_0) + \FT{\Gamma}(\omega + \omega_0)\big]
        + c_2\big[\FT{\Gamma}(\omega - 2\omega_0) + \FT{\Gamma}(\omega + 2\omega_0)\big]
    }
    \label{eq:A_hat}
\end{equation}
where $c_0$, $c_1$, $c_2$ are the Fourier coefficients of the visibility function given by Eq.~\eqref{eq:cn_general} (all real, with $c_{-n} = c_n$). Each term has a clear physical interpretation in the time domain:
\begin{itemize}[nosep]
    \item The \textbf{$n=0$} (DC ``direct current'') term, $c_0\,\Gamma(t)$, is the non-oscillatory component: the mean flux deficit that tracks the spot-size envelope. It is largest when the spot is visible for most of the rotation (e.g., a high-latitude spot on a low-inclination star), contributing a smooth, aperiodic dimming.
    \item The \textbf{$n=1$} (fundamental) term, $2c_1\,\Gamma(t)\cos(\omega_0 t)$, is the dominant periodic component at the stellar rotation frequency. It produces the familiar once-per-rotation brightness dip as the spot rotates in and out of view.
    \item The \textbf{$n=2$} (second harmonic) term, $2c_2\,\Gamma(t)\cos(2\omega_0 t)$, captures the non-sinusoidal shape of the visibility function (the flattened profile near disk center and sharp limb transitions), which a pure cosine cannot reproduce.
\end{itemize}
Figure~\ref{fig:fourier_decomposition} shows the Fourier decomposition of the single-spot signal $A(t) = \Gamma(t)\,\mathcal{V}(t)$ for three stellar inclinations at a fixed latitude $\Phi = 30^\circ$, with the same spot decay envelope $\Gamma(t)$. Here we see that terms up to second harmonic order are sufficient for capturing the shape of the signal for varying $I$.

\begin{figure}[ht!]
    \centering
    \script{plot_fourier_decomposition.py}
    \includegraphics[width=\linewidth]{figures/fig_fourier_decomposition.pdf}
    \caption{Fourier decomposition of the single-spot signal $A(t) = \Gamma(t) \cdot\mathcal{V}(t)$ for a spot at latitude $\Phi = 30^\circ$.
    \textit{Left:} Time-domain signal for three stellar inclinations, with the normalized squared envelope $\Gamma(t) = \alpha^2(t)/\alpha_{\max}^2$ shown as a dashed gray line.
    \textit{Right:} Normalized squared Fourier coefficients $|c_n|^2 / \sum |c_n|^2$ of the visibility function for each inclination.}
    \label{fig:fourier_decomposition}
\end{figure}

% Figure~\ref{fig:kernel_harmonic_components} illustrates how the individual $n=0$, $n=1$, and $n=2$ harmonic components combine to produce the full kernel shape for a single spot at a fixed latitude $\Phi = 30^\circ$. Each panel decomposes the normalized kernel $\mathcal{K}(\tau)/\mathcal{K}(0)$ (Eq.~\ref{eq:kernel_truncated}) for a different stellar inclination. The fractional power in each harmonic is defined as
% \begin{equation}
%     f_n = \frac{(2 - \delta_{n0})\,|c_n|^2}{\sum_{m=0}^{N}(2 - \delta_{m0})\,|c_m|^2},
%     \label{eq:fractional_power}
% \end{equation}
% where the factor $(2 - \delta_{n0})$ accounts for the positive and negative frequency contributions at each harmonic. At low inclination ($I = 30^\circ$), the DC term dominates ($f_0 \sim 89\%$), producing a kernel that decays nearly monotonically with only weak periodic modulation. As the inclination increases toward edge-on ($I = 90^\circ$), the fundamental becomes the leading component ($f_1 \sim 50\%$), and the kernel develops pronounced oscillations at the rotation period. The second harmonic remains a minor but non-negligible contribution ($f_2 \sim 9\%$ at $I = 90^\circ$), sharpening the peaks and introducing structure at half the rotation period. The shaded region under the $n=0$ component highlights the aperiodic ``floor'' of the kernel set by the DC visibility power.

% \begin{figure}[p]
%     \centering
%     \script{plot_kernel_harmonic_components.py}
%     \includegraphics[width=0.9\linewidth]{figures/fig_kernel_harmonic_components.pdf}
%     \caption{Harmonic decomposition of the kernel $\mathcal{K}(\tau)/\mathcal{K}(0)$ (Eq.~\ref{eq:kernel_truncated}) for a single spot at latitude $\Phi = 30^\circ$ at three stellar inclinations ($I = 30^\circ$, $60^\circ$, $90^\circ$), with fixed spot parameters ($P_{\rm eq} = 5$\,d, $\ell_{\rm spot} = 15$\,d, $\tau_{\rm spot} = 3$\,d, $\kappa = 0$). The black curve is the full kernel; colored curves show the individual $n=0$ (DC, blue), $n=1$ (fundamental, orange), and $n=2$ (second harmonic, green) components. The shaded blue region marks the DC contribution. Inset annotations give the fractional power $f_n$ (Eq.~\ref{eq:fractional_power}) in each harmonic. For pole-on viewing the kernel is dominated by the aperiodic DC term, while for edge-on viewing the fundamental produces strong periodic oscillations.}
%     \label{fig:kernel_harmonic_components}
% \end{figure}


% -----------------------------------------------
\subsection{Averaging Over Random Spot Parameters} \label{subsec:averaging}

The expression for $\FT{A}(\omega)$ in Eq.~\eqref{eq:A_hat_sum} describes a \emph{single} spot with a particular reference time $t_{\rm ref}$, reference longitude $\Lambda_{\rm ref}$, and latitude $\Phi_{\rm ref}$. To obtain the kernel for a \emph{population} of spots, we must average the power spectral density $|\FT{A}(\omega)|^2$ over these random variables.\footnote{Symbolic and numerical check: \NotebookLink{averaging_random_params.ipynb}} We treat $t_{\rm ref}$ and $\Lambda_{\rm ref}$ as uniformly distributed, and $\Phi_{\rm ref}$ as drawn from some latitude distribution $p(\Phi)$.

\subsubsection{Effect of $t_{\rm ref}$: time shift of the envelope}

Shifting the envelope by $t_{\rm ref}$ gives $\Gamma(t - t_{\rm ref})$, whose Fourier transform picks up a phase by the time-shifting property (\ref{ft:timeshift}, \ref{fp:delta}):
\begin{equation}
    \FT{\Gamma}(\omega;\, t_{\rm ref}) = \FT{\Gamma}(\omega)\, e^{-i\omega\, t_{\rm ref}}.
\end{equation}
This phase propagates to $\FT{A}(\omega)$, but since $|e^{-i\omega\, t_{\rm ref}}|^2 = 1$, it \textbf{drops out} of the power spectral density $|\FT{A}(\omega)|^2$. The kernel is therefore independent of $t_{\rm ref}$: any spot reference time yields the same autocorrelation.

\subsubsection{Effect of $\Lambda_{\rm ref}$: longitude phase}

Shifting the initial longitude $\Lambda_{\rm ref}$ is equivalent to a time shift of $\Lambda_{\rm ref}/\omega_0$ in the periodic visibility function. By the time-shifting property (\ref{ft:timeshift}), the Fourier series coefficients acquire a harmonic-dependent phase:
\begin{equation}
    c_n(\Lambda_{\rm ref}) = c_n\, e^{in\Lambda_{\rm ref}},
    \label{eq:cn_phase}
\end{equation}
where $c_n \equiv c_n(\Lambda_{\rm ref} = 0)$. Substituting into Eq.~\eqref{eq:A_hat_sum}, the full single-spot spectrum is:
\begin{equation}
    \FT{A}(\omega) = e^{-i\omega t_{\rm ref}} \sum_n c_n\, e^{in\Lambda_{\rm ref}}\, \FT{\Gamma}(\omega - n\omega_0).
    \label{eq:A_hat_random}
\end{equation}

The power spectral density is:
\begin{equation}
    |\FT{A}(\omega)|^2 = \sum_n \sum_m c_n\, c_m^*\; e^{i(n-m)\Lambda_{\rm ref}}\; \FT{\Gamma}(\omega - n\omega_0)\, \FT{\Gamma}^*(\omega - m\omega_0).
    \label{eq:esd_double_sum}
\end{equation}
Averaging over $\Lambda_{\rm ref}$ uniformly distributed on $[0, 2\pi]$:
\begin{equation}
    \left\langle e^{i(n-m)\Lambda_{\rm ref}} \right\rangle_{\Lambda_{\rm ref}}
    = \frac{1}{2\pi}\int_0^{2\pi} e^{i(n-m)\Lambda_{\rm ref}}\, d\Lambda_{\rm ref}
    = \delta_{nm}.
    \label{eq:longitude_orthogonality}
\end{equation}
This \textbf{cancels out all cross terms} between different harmonics. Only the diagonal ($n = m$) terms survive:
\begin{equation}
    \boxed{
    \left\langle |\FT{A}(\omega)|^2 \right\rangle_{t_{\rm ref},\, \Lambda_{\rm ref}}
    = \sum_{n=-\infty}^{\infty} |c_n|^2\; S_\Gamma(\omega - n\omega_0),
    }
    \label{eq:esd_averaged}
\end{equation}
where $S_\Gamma(\omega) \equiv |\FT{\Gamma}(\omega)|^2$ is the power spectral density of the squared envelope, and each harmonic of the visibility function contributes independently to the power spectrum, with weight $|c_n|^2$. 
% The random longitude destroys any coherent interference between different harmonics, leaving only incoherent (power-additive) contributions.

\subsubsection{From averaged PSD to the kernel}

Applying the Wiener--Khinchin theorem (Eq.~\ref{eq:wk}) to the averaged power spectral density:
\begin{equation}
    \mathcal{K}(\tau) = \frac{\alpha_{\max}^4}{2\pi}\int_{-\infty}^{\infty} \left\langle |\FT{A}(\omega)|^2 \right\rangle\; e^{i\omega\tau}\, d\omega
    = \alpha_{\max}^4 \sum_n |c_n|^2 \cdot \frac{1}{2\pi}\int_{-\infty}^{\infty} S_\Gamma(\omega - n\omega_0)\, e^{i\omega\tau}\, d\omega.
\end{equation}
Substituting $u = \omega - n\omega_0$:
\begin{equation}
    \frac{1}{2\pi}\int_{-\infty}^{\infty} S_\Gamma(u)\, e^{i(u + n\omega_0)\tau}\, du
    = e^{in\omega_0\tau} \cdot \underbrace{\frac{1}{2\pi}\int_{-\infty}^{\infty} S_\Gamma(u)\, e^{iu\tau}\, du}_{R_\Gamma(\tau)},
\end{equation}
where $R_\Gamma(\tau)$ is the autocorrelation of the normalized squared envelope $\Gamma(t) = \alpha^2(t)/\alpha_{\max}^2$. Therefore:
\begin{equation}
    \mathcal{K}(\tau) = \alpha_{\max}^4\, R_\Gamma(\tau) \sum_{n=-\infty}^{\infty} |c_n|^2\, e^{in\omega_0\tau}.
    \label{eq:kernel_complex}
\end{equation}

Since the $c_n$ are real with $c_{-n} = c_n$, the complex exponentials combine into cosines:
\begin{equation}
    \boxed{
    \mathcal{K}(\tau) = \alpha_{\max}^4\, R_\Gamma(\tau) \left[ |c_0|^2 + 2\sum_{n=1}^{\infty} |c_n|^2\, \cos(n\omega_0\tau) \right]
    }
    \label{eq:kernel_single_spot}
\end{equation}

The kernel factorizes into two interpretable components:
\begin{itemize}[nosep]
    \item $R_\Gamma(\tau)$: the \textbf{normalized envelope autocorrelation} of $\Gamma(t) = \alpha^2(t)/\alpha_{\max}^2$, which sets the overall decay timescale of correlations (governed by spot lifetime $\ell_{\rm spot}$ and rise/decay time $\tau_{\rm spot}$).
    \item The \textbf{cosine sum}: a quasi-periodic modulation at the rotation frequency and its harmonics, with amplitudes set by the squared Fourier coefficients $|c_n|^2$ of the visibility function.
\end{itemize}

Truncating to second harmonic order gets:
\begin{equation}
    \mathcal{K}(\tau) \approx \alpha_{\max}^4\, R_\Gamma(\tau) \left[ |c_0|^2 + 2|c_1|^2\cos(\omega_0\tau) + 2|c_2|^2\cos(2\omega_0\tau) \right],
    \label{eq:kernel_truncated}
\end{equation}
which represents the kernel function for a \emph{single} spot.

% -----------------------------------------------
\subsection{Multi-Spot Kernel} \label{subsec:multi_spot}

The observed flux deficit is the sum over $N_{\rm spot}$ independent spots (Eq.~\ref{eqn:flux_model}). With the small-spot projected area $A_k(t) = \pi\,\alpha_k^2(t)\,\mathcal{V}_k(t)$ (Eq.~\ref{eq:A_smallspot}), the factor of $\pi$ cancels the $1/\pi$ of the flux model:
\begin{equation}
    \delta F(t) = \frac{1}{\pi}\sum_{k=1}^{N_{\rm spot}} A_k(t)\,(1 - f_{{\rm spot},k})
    = \sum_{k=1}^{N_{\rm spot}} (1 - f_{{\rm spot},k})\,\alpha_k^2(t)\,\mathcal{V}_k(t).
\end{equation}
Each term is the single-spot flux deficit $\alpha^2(t)\,\mathcal{V}(t) = \alpha_{\max}^2 A(t)$, whose autocorrelation, averaged over the spot's random parameters, is the single-spot kernel of Eq.~\eqref{eq:kernel_single_spot}. That autocorrelation is integrated over time, so it has units of time. Spots emerge at random reference times $t_{{\rm ref},k}$, which we model as a Poisson process with rate $\dot{N}_{\rm spot}$ [spots/day]. For such a sum of independent pulses at random times (Campbell's theorem), the covariance is the emergence rate times the time-integrated autocorrelation of a single pulse. For spots with identical contrast $f_{\rm spot}$:
\begin{equation}
    \mathcal{K}(\tau) \equiv \Cov[\delta F(t),\, \delta F(t+\tau)] = \dot{N}_{\rm spot}\,(1 - f_{\rm spot})^2\; \left\langle \mathcal{K}(\tau) \right\rangle_{\Phi},
    \label{eq:kernel_Nspot}
\end{equation}
where $\langle \mathcal{K}(\tau) \rangle_\Phi$ is the single-spot kernel (Eq.~\ref{eq:kernel_single_spot}) averaged over the latitude distribution $p(\Phi)$. Equivalently, $N_{\rm spot}$ spots with reference times spread uniformly over a long baseline $T_{\rm obs}$ each contribute $\langle \mathcal{K}(\tau) \rangle_\Phi / T_{\rm obs}$, so that $\dot{N}_{\rm spot} = N_{\rm spot}/T_{\rm obs}$.

% -----------------------------------------------
\subsubsection{Averaging over latitude}

The latitude $\Phi_{\rm ref}$ affects both the Fourier coefficients $c_n$ (through $b_0 = \cos I\sin\Phi$ and $b_1 = \sin I\cos\Phi$) and, in the presence of differential rotation, the rotation frequency $\omega_0(\Phi) = 2\pi(1 - \kappa\sin^2\Phi)/P_{\rm eq}$. Averaging Eq.~\eqref{eq:kernel_single_spot} over $p(\Phi)$:
\begin{equation}
    \left\langle \mathcal{K}(\tau) \right\rangle_\Phi
    = \alpha_{\max}^4 \int_{-\pi/2}^{\pi/2} R_\Gamma(\tau) \left[ |c_0(\Phi)|^2 + 2\sum_{n=1}^{N} |c_n(\Phi)|^2\, \cos\!\big(n\omega_0(\Phi)\,\tau\big) \right] p(\Phi)\, d\Phi.
    \label{eq:kernel_lat_avg}
\end{equation}
In general, this integral must be evaluated numerically.

% -----------------------------------------------
\subsubsection{Full kernel} \label{subsec:kernel_full}

Combining Eqs.~\eqref{eq:kernel_Nspot} and \eqref{eq:kernel_lat_avg}, the complete GP kernel for the flux variability due to a population of independent spots is:
\begin{equation}
    \boxed{
    \mathcal{K}(\tau) = \dot{N}_{\rm spot}\,(1 - f_{\rm spot})^2\,\alpha_{\max}^4
    \int_{-\pi/2}^{\pi/2} R_\Gamma(\tau) \left[\, |c_0(I,\Phi)|^2 + 2\sum_{n=1}^{N} |c_n(I,\Phi)|^2\, \cos\!\big(n\,\omega_0(\Phi)\,\tau\big) \right] p(\Phi)\, d\Phi
    }
    \label{eq:kernel_full}
\end{equation}
where $R_\Gamma(\tau)$ is the normalized autocorrelation of the squared spot-size envelope $\Gamma(t) = \alpha^2(t)/\alpha_{\max}^2$, $c_n(I,\Phi)$ are the Fourier coefficients of the visibility function (Eq.~\ref{eq:cn_general}), $\omega_0(\Phi) = 2\pi(1 - \kappa\sin^2\Phi)/P_{\rm eq}$ is the latitude-dependent rotation frequency, and $p(\Phi)$ is the spot latitude distribution.

The stellar inclination enters the kernel \emph{only} through the visibility coefficients $c_n$; the rotation frequency $\omega_0(\Phi)$ is independent of $I$. Writing $b_0 = \cos I \sin\Phi$ and $b_1 = \sin I \cos\Phi$ out explicitly, the half-width of the visibility window depends on $I$ and $\Phi$ through the single combination $\cot I \tan\Phi$,
\begin{equation}
    \theta_{\rm vis}(I,\Phi) = \arccos\!\big(-\cot I\,\tan\Phi\big),
    \label{eq:theta_vis_inc}
\end{equation}
and the Fourier coefficients of Eq.~\eqref{eq:cn_general} become
\begin{equation}
    c_n(I,\Phi) = \frac{1}{\pi}\left[
        \frac{\cos I\sin\Phi}{n}\,\sin\!\big(n\,\theta_{\rm vis}\big)
        + \frac{\sin I\cos\Phi}{2}\left(
            \frac{\sin\!\big((n-1)\theta_{\rm vis}\big)}{n-1}
          + \frac{\sin\!\big((n+1)\theta_{\rm vis}\big)}{n+1}
        \right)\right].
    \label{eq:cn_inc}
\end{equation}
Equation~\eqref{eq:theta_vis_inc} partitions the stellar surface at the latitudes $\Phi = \pm I$, as described following Eq.~\eqref{eq:cos_beta}: spots at $\Phi \geq I$ are circumpolar ($\theta_{\rm vis} = \pi$), spots at $\Phi \leq -I$ never rotate into view ($\theta_{\rm vis} = 0$, $c_n \equiv 0$), and only the band $|\Phi| < I$ contributes rotationally modulated power, with duty cycle $\Delta = \theta_{\rm vis}/\pi$. We stress that the integral in Eq.~\eqref{eq:kernel_full} runs over the full physical range $[-\pi/2, \pi/2]$ and not the visible range $(-I, \pi/2]$: the two are equivalent because $c_n$ vanishes identically on the hidden cap, and retaining the full range is what keeps $p(\Phi)$ normalized over the whole star. Restricting the bounds and renormalizing $p(\Phi)$ over them would rescale the kernel amplitude by $\pi/(\pi/2 + I)$.

Two limits are instructive. For edge-on viewing ($I = \pi/2$) the $\cot I$ term vanishes, so $\theta_{\rm vis} = \pi/2$ at every latitude and $c_n(\pi/2,\Phi) = g_n\cos\Phi$, recovering the closed-form kernel of Eq.~\eqref{eq:kernel_edgeon_closed} below. For pole-on viewing ($I \to 0$) we have $b_1 \to 0$, the visibility function reduces to the constant $\sin\Phi$ over the northern hemisphere, all $c_{n \geq 1} \to 0$, and the kernel collapses to pure envelope decay with no rotational signal, $\mathcal{K}(\tau) \to \sigma_k^2 R_\Gamma(\tau)/4$ for $p(\Phi) = 1/\pi$. Decreasing the inclination therefore does not reduce the total variance $\mathcal{K}(0)$; it transfers power out of the rotation harmonics and into the aperiodic $n = 0$ term.

The full hyperparameters of this kernel are:
\begin{equation}
    \theta = \{ P_{\rm eq},\; \kappa,\; I,\; \dot{N}_{\rm spot},\; \ell_{\rm spot},\; \tau_{\rm spot},\; \alpha_{\rm max},\; f_{\rm spot} \}.
\end{equation}

The amplitude in Eq.~\eqref{eq:kernel_full} is set by the spot emergence rate $\dot{N}_{\rm spot}$ [spots/day], not by the total number of spots. For a time series of duration $T_{\rm obs}$ the expected number of spots is $N_{\rm spot} = \dot{N}_{\rm spot}\,T_{\rm obs}$, but a longer baseline adds spots without changing the variance of the flux. The amplitude prefactor $\dot{N}_{\rm spot}\,(1 - f_{\rm spot})^2\,\alpha_{\max}^4$ is therefore independent of the observing baseline, an intrinsic stellar property. In our \code{spotgp} implementation the amplitude is specified through \texttt{nspot\_rate} (emergence rate in spots/day), which is the recommended parameterization for inference.

The amplitude prefactor is degenerate in practice (only the product can be constrained from the lightcurve variance) so it can be absorbed into a single amplitude hyperparameter $\sigma_k^2$. Defining the spot contrast-area product as $a_{\rm spot} = (1 - f_{\rm spot})\,\alpha_{\max}^2$, then:
\begin{equation}
    \sigma_k^2 \equiv \dot{N}_{\rm spot}\,a_{\rm spot}^2,
    \label{eq:sigma_k_cspot}
\end{equation}
and the reduced hyperparameter set becomes:
\begin{equation}
    \theta = \{ P_{\rm eq},\, \kappa,\, I,\, \ell_{\rm spot},\, \tau_{\rm spot},\, \sigma_k \}.
\end{equation}

% -----------------------------------------------
\subsubsection{Analytic solution for $\kappa=0$ and $I=90$} \label{subsec:kernel_anal_case}

In some cases (e.g. a solid-body rotator viewed at 90 deg inclination), there is a closed-form analytic solution for the kernel.\footnote{Symbolic and numerical checks: \NotebookLink{closed_form_kernel_n2.ipynb} for $n \leq 2$ and general $I$, and \NotebookLink{closed_form_kernel.ipynb} for all harmonics.} Truncating to second harmonic order ($n \leq 2$) and specializing to solid-body rotation ($\kappa = 0$), the cosines factor out of the latitude integral and the kernel becomes:
\begin{equation}
    \mathcal{K}(\tau) = \sigma_k^2\, R_\Gamma(\tau) \left[
        \langle |c_0|^2 \rangle_\Phi
        + 2\langle |c_1|^2 \rangle_\Phi \cos(\omega_0\tau)
        + 2\langle |c_2|^2 \rangle_\Phi \cos(2\omega_0\tau)
    \right],
    \label{eq:kernel_n2}
\end{equation}
where $\omega_0 = 2\pi/P_{\rm eq}$ and the latitude-averaged squared coefficients are $\langle |c_n|^2 \rangle_\Phi = \int_{-\pi/2}^{\pi/2} |c_n(\Phi)|^2\, p(\Phi)\, d\Phi$, with $c_n(\Phi)$ given by Eq.~\eqref{eq:cn_general}.

For the special case of edge-on viewing ($I = \pi/2$) and a uniform latitude distribution $p(\Phi) = 1/\pi$ on $[-\pi/2,\pi/2]$, all Fourier coefficients are proportional to $\cos\Phi$: $c_n(\Phi) = g_n\cos\Phi$, where $g_0 = 1/\pi$, $g_1 = 1/4$, and $g_2 = 1/(3\pi)$. The latitude integrals evaluate to $\langle |c_n|^2 \rangle_\Phi = g_n^2/2$, giving the fully closed-form kernel:
\begin{equation}
    \mathcal{K}(\tau) = \sigma_k^2\, R_\Gamma(\tau) \left[
        \frac{1}{2\pi^2}
        + \frac{1}{16}\cos(\omega_0\tau)
        + \frac{1}{9\pi^2}\cos(2\omega_0\tau)
    \right]
    \quad (I = \pi/2,\; \kappa = 0).
    \label{eq:kernel_edgeon_closed}
\end{equation}

For general inclination $I$, the latitude integrals $\langle |c_n|^2 \rangle_\Phi$ must be evaluated numerically, since $c_0$ and $c_1$ depend on $\theta_{\rm vis}(\Phi) = \arccos(-\cot I\tan\Phi)$, which prevents elementary integration. However, the numerical quadrature over $\Phi$ is inexpensive (a single 1D integral per harmonic).
This simplified form (\ref{eq:kernel_edgeon_closed}) is applicable when a system is viewed edge-on ($I = \pi/2$) with no differential rotation ($\kappa = 0$) and under the assumption that spots are uniformly distributed across all latitudes, however it is flexible enough to work for different envelope terms $R_\Gamma(\tau)$.

% -----------------------------------------------
\subsection{Power Spectral Density of the Full Kernel} \label{subsec:psd}

By the Wiener--Khinchin theorem (Eq.~\ref{eq:wk}), the kernel $\mathcal{K}(\tau)$ and its power spectral density (PSD) $S_{\mathcal{K}}(\omega)$ form a Fourier transform pair:
\begin{equation}
    S_{\mathcal{K}}(\omega) = \int_{-\infty}^{\infty} \mathcal{K}(\tau)\, e^{-i\omega\tau}\, d\tau.
    \label{eq:psd_def}
\end{equation}
Since $\mathcal{K}(\tau)$ is real and even, $S_{\mathcal{K}}(\omega)$ is real, non-negative, and even in $\omega$. To evaluate the transform, we use the frequency-shift property: for a real, even function $R_\Gamma(\tau)$ with transform $S_\Gamma(\omega) = \mathcal{F}\{R_\Gamma\}$,
\begin{equation}
    \mathcal{F}\{R_\Gamma(\tau)\cos(n\omega_0\tau)\} = \tfrac{1}{2}\big[S_\Gamma(\omega - n\omega_0) + S_\Gamma(\omega + n\omega_0)\big],
\end{equation}
which follows from writing $\cos(n\omega_0\tau) = \tfrac{1}{2}(e^{in\omega_0\tau} + e^{-in\omega_0\tau})$ and applying the frequency-shift property (\ref{ft:freqshift}).

Applying this to the solid-body rotation kernel (Eq.~\ref{eq:kernel_n2}) yields:
\begin{equation}
    S_{\mathcal{K}}(\omega) = \sigma_k^2 \left[
        \langle |c_0|^2 \rangle_\Phi\, S_\Gamma(\omega)
        + \sum_{n=1}^{N} \langle |c_n|^2 \rangle_\Phi \Big[S_\Gamma(\omega - n\omega_0) + S_\Gamma(\omega + n\omega_0)\Big]
    \right],
    \quad (\kappa = 0),
    \label{eq:psd_solid}
\end{equation}
where $S_\Gamma(\omega) = |\FT{\Gamma}(\omega)|^2$ is the energy spectral density of the squared envelope and $\omega_0 = 2\pi/P_{\rm eq}$. The PSD is a sum of shifted copies of the envelope spectrum $S_\Gamma(\omega)$, each centered on a harmonic $n\omega_0$ and weighted by the corresponding latitude-averaged squared Fourier coefficient $\langle |c_n|^2 \rangle_\Phi$. The $n = 0$ term produces a low-frequency component centered at $\omega = 0$, while each $n \geq 1$ term contributes a pair of symmetric peaks at $\omega = \pm n\omega_0$, whose width is set by the envelope spectral width (inversely proportional to the spot lifetime).

For the general case with differential rotation ($\kappa \neq 0$), the latitude-dependent rotation frequency $\omega_0(\Phi)$ prevents factoring the cosine out of the latitude integral. Taking the Fourier transform of the full kernel (Eq.~\ref{eq:kernel_full}):
\begin{equation}
\boxed{
    S_{\mathcal{K}}(\omega) = \sigma_k^2 \int_{-\pi/2}^{\pi/2}
    \left[
        |c_0(\Phi)|^2\, S_\Gamma(\omega)
        + \sum_{n=1}^{N} |c_n(\Phi)|^2 \Big[S_\Gamma\!\big(\omega - n\omega_0(\Phi)\big) + S_\Gamma\!\big(\omega + n\omega_0(\Phi)\big)\Big]
    \right] p(\Phi)\, d\Phi.
}
    \label{eq:psd_full}
\end{equation}
Differential rotation causes the harmonic peaks to broaden: spots at different latitudes contribute power at slightly different frequencies $n\omega_0(\Phi)$, smearing each peak over the range $n\omega_0(0)$ to $n\omega_0(\pi/2) = n\omega_0(0)(1 - \kappa)$.

\vspace{\baselineskip}
% ========================================================================================
\section{Physical Interpretation of Kernel Hyperparameters} \label{sec:hyperparameters}

% \subsection{Information Content of the Mean and Covariance} \label{sec:mean_function}
% Before discussing the individual hyperparameters, we illustrate how the single-spot signal connects to the two statistical quantities modeled by the GP: the mean function (first moment) and the kernel (second moment). Figure~\ref{fig:moment_decomposition} shows 80 realizations of the single-spot signal $A(t) = \Gamma(t)\,\mathcal{V}(t)$ for random initial longitudes $\Lambda_{\rm ref}$, at fixed inclination $I = 60^\circ$ and latitude $\Phi = 30^\circ$. The ensemble mean over all longitudes is the smooth, aperiodic curve $\langle A(t) \rangle_\Lambda = c_0\,\Gamma(t)$: the spot-size envelope scaled by the zeroth Fourier coefficient of the visibility function. This is the first moment of the process: the quantity that enters the GP mean function (Section~\ref{sec:mean_function}). It depends linearly on $c_0$ and on the first moment of the envelope, $\int \Gamma\,dt$.

% The lower panel shows the residuals after subtracting this mean. These oscillatory residuals (produced by the periodic rotation of the spot in and out of view) are what the kernel (second moment) describes. Their variance depends quadratically on the Fourier coefficients ($|c_n|^2$) and on the autocorrelation of the envelope, $R_\Gamma(\tau) = \int \Gamma(t)\,\Gamma(t+\tau)\,dt$. The distinction between these two moments (linear versus quadratic dependence on the visibility coefficients, and first versus second functional of the envelope) is what allows the mean and kernel to constrain different combinations of the amplitude parameters.

% \begin{figure}[ht!]
%     \centering
%     \script{plot_moment_decomposition.py}
%     \includegraphics[width=\linewidth]{figures/fig_moment_decomposition.pdf}
%     \caption{Decomposition of the single-spot signal into first and second moments of the GP, for a spot at latitude $\Phi = 30^\circ$ with $I = 60^\circ$, $P_{\rm eq} = 5$\,d, $\ell_{\rm spot} = 15$\,d, and $\tau_{\rm spot} = 3$\,d. \textit{Top:} 80 realizations of $A(t) = \Gamma(t)\,\mathcal{V}(t)$ for random longitudes (blue), with the ensemble mean $\langle A \rangle_\Lambda = c_0\,\Gamma(t)$ (red). The mean depends linearly on $c_0$ and on $\int\Gamma\,dt$ (first moment of the envelope). \textit{Bottom:} Residuals after subtracting the mean, with $\pm\sigma(t)$ envelope (orange). This spread is what the kernel describes; it depends quadratically on $|c_n|^2$ and on the envelope autocorrelation $R_\Gamma(\tau) = \int\Gamma(t)\,\Gamma(t+\tau)\,dt$ (second moment of the envelope).}
%     \label{fig:moment_decomposition}
% \end{figure}

\subsection{Kernel Components and Hyperparameters} \label{subsec:kernel_components}

The full kernel (Eq.~\ref{eq:kernel_full}) is a sum over rotation harmonics, and the cleanest way to understand what each hyperparameter does is to keep it that way. Writing
\begin{equation}
    \mathcal{K}(\tau) = \sum_{n=0}^{N} \mathcal{K}_n(\tau),
    \qquad
    \mathcal{K}_n(\tau) = \sigma_k^2\, s_n \int_{-\pi/2}^{\pi/2} R_\Gamma(\tau)\; |c_n(I,\Phi)|^2 \cos\!\big(n\,\omega_0(\Phi)\,\tau\big)\; p(\Phi)\, d\Phi,
    \label{eq:kernel_harmonic_sum}
\end{equation}
with $s_0 = 1$ and $s_n = 2$ for $n \geq 1$, each harmonic term is the product of exactly three factors, and every physical hyperparameter enters through exactly one of them:
\begin{itemize}[nosep]
    \item \textbf{The envelope autocorrelation $R_\Gamma(\tau)$} encodes how long spots live. It carries $\ell_{\rm spot}$ and $\tau_{\rm spot}$, and it multiplies every order identically.
    \item \textbf{The geometry weight $|c_n(I,\Phi)|^2$} sets how much variance lands in order $n$. It carries the inclination $I$ and the latitude distribution $p(\Phi)$, and nothing else.
    \item \textbf{The rotation phase $\cos(n\,\omega_0(\Phi)\tau)$} carries $P_{\rm eq}$ and $\kappa$ and sets the positions of the kernel peaks.
\end{itemize}
The amplitude $\sigma_k^2$ multiplies all orders equally. Because the three harmonics are sensitive to disjoint pieces of the physics, they are not three views of the same measurement: which parameters a fit can recover is decided by which orders the data constrain. The subsections below examine each order in turn; all quoted numbers are computed with \code{spotgp} (\code{AnalyticKernel.kernel\_components}, Gauss--Legendre latitude quadrature) for a fiducial model with $\ell_{\rm spot} = 15$\,d, $\tau_{\rm spot} = 3$\,d, $I = 60^\circ$, $P_{\rm eq} = 5$\,d, uniform $p(\Phi)$, varying one parameter at a time.\footnote{The accompanying notebook is \texttt{src/notebooks/harmonics\_plots.ipynb}.}

Each $\mathcal{K}_n$ is separately a valid (positive semi-definite) kernel: the $n = 0$ term has spectrum $|\FT{\Gamma}(\omega)|^2 \langle|c_0|^2\rangle_\Phi \geq 0$, the $n \geq 1$ terms have spectra $[S_\Gamma(\omega - n\omega_0) + S_\Gamma(\omega + n\omega_0)]$ that are shifted copies of the same non-negative function, and a positive mixture over latitude preserves both statements. The harmonics are therefore \emph{independently samplable} stochastic processes whose sum is a draw from the full GP. Figure~\ref{fig:harmonic_timeseries} shows this decomposition in the time domain: the $n = 0$ component is an aperiodic random walk--like dimming set by spot emergence and decay, the $n = 1$ component is the familiar rotational modulation, and the $n = 2$ component is a weaker signal at half the rotation period; their sum is statistically identical to a draw from the full kernel. Figure~\ref{fig:harmonic_decomposition_single} shows the same decomposition of the kernel and its PSD: at the fiducial inclination the $n = 0$, 1, 2 terms contribute $60\%$, $36\%$, and $4\%$ of the variance.

\begin{figure}[p]
    \centering
    \script{plot_harmonic_timeseries_samples.py}
    \includegraphics[width=0.9\linewidth]{figures/fig_harmonic_timeseries_samples.pdf}
    \caption{Harmonic decomposition of GP lightcurve samples ($I = 60^\circ$, $P_{\rm eq} = 5$\,d, $\kappa = 0$, $\ell_{\rm spot} = 15$\,d, $\tau_{\rm spot} = 3$\,d). Each of the top three panels is an independent draw from a single harmonic term $\mathcal{K}_n$ of Eq.~\eqref{eq:kernel_harmonic_sum}: the aperiodic $n = 0$ (spot emergence and decay), the fundamental $n = 1$ (rotational modulation at $P_{\rm eq}$), and the second harmonic $n = 2$ (at $P_{\rm eq}/2$). Annotations give each order's share of the variance, $\mathcal{K}_n(0)/\mathcal{K}(0)$. Because the terms are independently samplable, their sum (fourth panel) is a draw from the full kernel.}
    \label{fig:harmonic_timeseries}
\end{figure}

\begin{figure}[ht!]
    \centering
    \script{plot_kernel_psd_harmonic_decomposition_single.py}
    \includegraphics[width=\linewidth]{figures/fig_kernel_psd_harmonic_decomposition_single.pdf}
    \caption{Harmonic decomposition of the kernel (left) and its power spectral density (right) for the fiducial model ($P_{\rm eq} = 5$\,d, $\kappa = 0.2$, $I = 60^\circ$, $\ell_{\rm spot} = 15$\,d, $\tau_{\rm spot} = 3$\,d). The $n = 0$ term (blue) is the pure envelope autocorrelation: aperiodic, spanning the low-frequency lobe of the PSD; the $n = 1$ term (orange) carries the rotation peak at $1/P_{\rm eq}$; the $n = 2$ term (pink) carries the first-harmonic peak at $2/P_{\rm eq}$. Annotations give the fractional contribution of each order to the variance $\mathcal{K}(0)$.}
    \label{fig:harmonic_decomposition_single}
\end{figure}


% -----------------------------------------------
\subsubsection{The $n=0$ Term: A Direct Readout of the Spot Lifetime} \label{subsec:harmonic_n0}

With no cosine factor, the $n = 0$ term is the envelope autocorrelation times a constant,
$\mathcal{K}_0(\tau) = \sigma_k^2\,\langle|c_0|^2\rangle_\Phi\, R_\Gamma(\tau)$.
Its shape is a direct readout of how long spots live, uncontaminated by anything else: for the trapezoidal envelope it decays through the piecewise-polynomial breakpoints at $\tau = \tau_{\rm spot}$, $\ell_{\rm spot}$, $\ell_{\rm spot} + \tau_{\rm spot}$ and reaches zero \emph{exactly} at $\tau = \ell_{\rm spot} + 2\tau_{\rm spot}$ (21\,d for the fiducial model; Appendix~\ref{app:acf_trapezoid}). The plateau carries most of the leverage: doubling $\ell_{\rm spot}$ from 15 to 30\,d moves the half-height lag from 8.9 to 16.4\,d, while $\tau_{\rm spot}$ enters at roughly half that weight.

Because $P_{\rm eq}$ and $\kappa$ enter only through the cosine factor, the $n = 0$ term is \emph{completely blind} to both: the normalized $\mathcal{K}_0$ is identical for any rotation period or shear. This invariance is what makes it useful: rotation cannot leak into it, so it is the fixed reference against which the decay of the $n \geq 1$ orders is measured (Section~\ref{subsec:harmonic_kappa}). Its amplitude, by contrast, is pure geometry: $\langle|c_0|^2\rangle_\Phi$ grows by a factor ${\sim}5$ from edge-on to $I = 10^\circ$ as more of the spot population becomes permanently visible.

The compact support of $\mathcal{K}_0$ also provides a clean model-checking diagnostic. Unlike the exponential or Gaussian decays of phenomenological kernels, which are strictly positive at all lags, the starspot kernel vanishes identically beyond $\ell_{\rm spot} + 2\tau_{\rm spot}$; any observed correlation at longer lags must come from other sources (activity cycles, instrument systematics), or indicates a longer spot lifetime. Combined with the rotational recurrences of the $n \geq 1$ terms, this sets the number of visible ACF peaks,
\begin{equation}
    N_{\rm peaks} \approx \left\lfloor \frac{\ell_{\rm spot} + 2\tau_{\rm spot}}{P_{\rm eq}} \right\rfloor,
    \label{eq:n_peaks}
\end{equation}
the starspot analog of the quality factor of a damped oscillator kernel (in practice the last peak sits on the tail of $R_\Gamma$ and is barely visible).

Finally, the $n = 0$ harmonic is the point of contact between the kernel and the mean function: $c_0 = \bar{\mathcal{V}}$ enters the mean linearly (Section~\ref{sec:mean_function}) and the kernel quadratically. The $|c_0|^2$ floor keeps the kernel non-negative at all lags: each spot only ever dims the star, and the autocorrelation of a one-signed pulse is non-negative, whereas sample ACFs of real lightcurves dip negative at half-period lags. The difference is an estimator effect: subtracting the finite-window time mean biases the sample ACF downward by roughly $\bar{\mathcal{K}} = T^{-1}\!\int\mathcal{K}\,ds$, which is dominated by the $n = 0$ power. The DC power fraction $f_0 = \mathcal{K}_0(0)/\mathcal{K}(0)$ quantifies how large this effect can be, growing from $0.41$ at $I = 90^\circ$ toward unity pole-on. Since the $n = 0$ fluctuations are genuinely stochastic (the number of visible spots varies), the term belongs in the kernel even for mean-subtracted data; the finite-window bias is better handled by fitting or marginalizing an explicit flux offset (Section~\ref{sec:offsets}) than by removing the term.

In practice, a second and typically larger effect compounds the estimator bias above. That bias is a statistical artifact of not knowing the true baseline flux level: it would persist even for a perfectly instrumented, unfiltered lightcurve, simply because the sample mean used in its place is itself contaminated by the same low-frequency power it is meant to remove. The second effect is different in kind: the lightcurve detrending applied upstream of the GP fit (e.g., a Savitzky--Golay filter with a window of tens of days) is a real spectral filter, applied by the reduction pipeline before the GP ever sees the data, that removes power at exactly the frequencies the DC bump occupies. The bump has half-width $\sim 1/(\ell_{\rm spot} + 2\tau_{\rm spot})$, comparable to or narrower than the typical detrending window, so the filter does not merely attenuate the $n = 0$ term uniformly; it can carve out most of it. Because $\ell_{\rm spot}$ and $\tau_{\rm spot}$ are shared with the $n \geq 1$ harmonics through $R_\Gamma(\tau)$, forcing the kernel to match its full geometric DC prediction against a lightcurve with that power removed pulls the fitted envelope timescale away from its true value to compensate; and because the ratio $\langle|c_0|^2\rangle_\Phi/\langle|c_1|^2\rangle_\Phi$ is a principal lever for constraining the inclination (Section~\ref{subsec:kernel_components}), that constraint degrades along with it.

We therefore do not treat the $n = 0$ amplitude as fixed by the geometry alone. Instead we introduce a free multiplicative scale $s_0$ on the DC term,
\begin{equation}
    \mathcal{K}(\tau) = \sigma_k^2\, R_\Gamma(\tau) \left[ s_0\, \langle|c_0|^2\rangle_\Phi + 2\sum_{n \geq 1} \langle|c_n|^2\rangle_\Phi\, \cos(n\,\omega_0(\Phi)\,\tau) \right],
    \label{eq:kernel_s0}
\end{equation}
with $s_0 \geq 0$ the only requirement for positive semi-definiteness, since it multiplies an already positive-definite sub-kernel. $s_0 = 1$ recovers the standard kernel (Eq.~\ref{eq:kernel_harmonic_sum}) exactly, and $s_0 = 0$ is equivalent to dropping the $n = 0$ term and retaining only $n \geq 1$ harmonics; intermediate $s_0 < 1$ absorbs whatever fraction of the geometric DC power is missing. We do not attempt to separate the two suppression mechanisms above: the estimator bias and the detrending filter both act on the same term with the same sign, so their combined effect is what a single scale factor can hope to capture, not either one individually.
% We apply this parameterization to the Kepler-63 fits of Section~\ref{sec:kepler63} and compare the Bayesian evidence of the free-$s_0$, fixed-$s_0 = 1$, and dropped-$s_0 = 0$ models directly, since all three share the same parameter count modulo $s_0$ itself. This is a diagnostic rather than a fully principled correction: the detrending filter's response is broader than the DC bump and reshapes rather than uniformly rescales it, so a single scalar cannot repair the lost spectral shape; the more rigorous fix is to fold the filter's transfer function directly into the kernel, $\mathcal{K} \to F\mathcal{K}F^\top$.

% -----------------------------------------------
\subsubsection{The $n=1$ Term: Rotation Phase and the Differential-Rotation Gap} \label{subsec:harmonic_n1}

The fundamental multiplies the same envelope by $\cos(\omega_0\tau)$. Under solid-body rotation this adds no new timescale: at integer multiples of the period the cosine is unity, so the recurrence heights of the $n = 1$ term land exactly on the $n = 0$ envelope,
\begin{equation}
    \frac{\mathcal{K}(m P_{\rm eq})}{\mathcal{K}(0)} = \frac{R_\Gamma(m P_{\rm eq})}{R_\Gamma(0)}
    \qquad (\kappa = 0),
    \label{eq:peak_height_ratio}
\end{equation}
which evaluates to $0.741$ and $0.432$ at one and two periods for the fiducial model. When $m P_{\rm eq}$ falls in the linear regime of $R_\Gamma$ ($\tau_{\rm spot} < m P_{\rm eq} < \ell_{\rm spot}$), the peak heights decay linearly,
\begin{equation}
    \frac{\mathcal{K}(mP_{\rm eq})}{\mathcal{K}(0)} = \frac{\ell_{\rm spot} + 2\tau_{\rm spot}/3 - mP_{\rm eq}}{\ell_{\rm spot} + 2\tau_{\rm spot}/5},
    \label{eq:linear_peak_decay}
\end{equation}
a distinctive prediction of the trapezoidal profile (exponential profiles give exponentially decaying peaks). Under solid-body rotation, then, all of the \emph{new} information in the fundamental is in where its peaks sit, not how tall they are: peak positions measure $P_{\rm eq}$, peak heights re-measure the envelope.

Differential rotation is what breaks this tie. Spots at latitude $\Phi$ accumulate phase $\omega_0(\Phi)\tau$, with a characteristic spread across the latitude distribution of
\begin{equation}
    \Delta\phi(\tau) = \frac{2\pi\,|\kappa|\,\tau}{P_{\rm eq}}\,\langle \sin^2\Phi \rangle_\Phi
    \label{eq:phase_spread}
\end{equation}
(the full range across latitudes is $2\pi|\kappa|\tau/P_{\rm eq}$). Because $\cos(\omega_0(\Phi)\tau) = 1$ for every latitude at $\tau = 0$, the shear leaves $\mathcal{K}_1(0)$ and the variance shares untouched at any $\kappa$; what it does instead is pull the recurrence heights \emph{below} the $n = 0$ envelope and skew them toward longer lags (high-latitude spots recur later for $\kappa > 0$). For the fiducial geometry the first recurrence retains only $0.58$ of its solid-body height at $\kappa = 0.4$. \textbf{The gap between the $n = 1$ recurrence heights and the $n = 0$ envelope is the shear; nothing else opens it.} Setting $\Delta\phi \approx \pi$ gives the coherence lag
\begin{equation}
    \tau_{\rm coh} \approx \frac{P_{\rm eq}}{2\,|\kappa|\,\langle \sin^2\Phi \rangle_\Phi},
    \label{eq:coherence_lag}
\end{equation}
beyond which the rotational recurrences are strongly suppressed (roughly halved, for the fiducial latitude distribution) even where $R_\Gamma$ has not yet decayed: differential rotation imposes a second coherence timescale competing with the spot lifetime. The dephasing depends only on $|\kappa|$; the sign of $\kappa$ appears solely in the direction of the peak skew. Figure~\ref{fig:decoherence_samples} shows this decoherence in both the lightcurves and the ACF: as $|\kappa|$ increases, GP samples lose their periodic pattern and the ACF recurrences broaden and wash out, while the $R_\Gamma$ envelope (the $n = 0$ reference) remains unchanged.

\begin{figure}[ht!]
    \centering
    \script{plot_decoherence_samples.py}
    \includegraphics[width=\linewidth]{figures/fig_decoherence_samples.pdf}
    \caption{Decoherence from differential rotation. \textit{Left:} Three GP lightcurve samples drawn from the kernel for each $\kappa$ value. \textit{Right:} Normalized kernel ACF (solid) with the envelope autocorrelation $R_\Gamma(\tau)/R_\Gamma(0)$ (dashed). For $\kappa = 0$ (solid-body rotation), the lightcurves are highly periodic and the ACF recurrences ride the $n = 0$ envelope. As $|\kappa|$ increases, spots at different latitudes rotate at different rates and the periodic signal decoheres: at $|\kappa| = 0.5$ the signal is noticeably aperiodic beyond the first two recurrences ($\tau_{\rm coh} \approx 10$\,d from Eq.~\ref{eq:coherence_lag}), and at $|\kappa| = 1$ only the first survives ($\tau_{\rm coh} \approx 5$\,d). The envelope is unchanged throughout, since it depends only on the spot lifetime parameters; the widening gap between the recurrence heights and the envelope is the shear signal. The decoherence is symmetric in $\pm\kappa$. Parameters: $P_{\rm eq} = 5$\,d, $I = 90^\circ$, $\ell_{\rm spot} = 30$\,d, $\tau_{\rm spot} = 3$\,d.}
    \label{fig:decoherence_samples}
\end{figure}

Two practical caveats follow directly from the factorization. First, short-lived spots bias the recovered period low: at $\ell_{\rm spot} = 5$\,d the first kernel recurrence sits at 4.73\,d for a true 5.0\,d period, because the envelope decays across the peak. Second, when the spot lifetime falls below the rotation period the recurrence lands where the envelope has already died (at $P_{\rm eq} = 20$\,d with a 21\,d lifetime, $\mathcal{K}_1(P_{\rm eq})$ is zero to three decimal places), and no rotational information survives at any inclination. The amplitude of the fundamental is the mirror image of the $n = 0$ term: half the variance equator-on, collapsing to $1.4\%$ by $I = 10^\circ$.

% -----------------------------------------------
\subsubsection{The $n=2$ Term: Evidence That Spots Set} \label{subsec:harmonic_n2}

A spot that never sets has projected area exactly $b_0 + b_1\cos(\omega_0 t)$: a pure sinusoid, with no second harmonic to give. Only the clipping of the visibility function at the limb generates power at $n \geq 2$ (Appendix~\ref{app:fourier_visibility}). The second harmonic is therefore not a smaller copy of the fundamental; it is a statement that spots are disappearing behind the star. The kernel is unambiguous about this: restricting spots to circumpolar latitudes ($\Phi \geq I$; e.g.\ a $60^\circ$--$90^\circ$ band at $I = 60^\circ$) drives $\mathcal{K}_2(0)$ to zero exactly while the fundamental retains $12\%$ of the variance.

The share of variance at $n = 2$ is small everywhere: $9\%$ of the total at best (equator-on), $3.7\%$ at the fiducial $I = 60^\circ$, and ${\sim}10^{-4}$ pole-on, so detecting it is a signal-to-noise question before it is a physics question. When it is detected, it adds two things the fundamental cannot. Its recurrence at $P_{\rm eq}/2$ survives envelopes too short for the fundamental (in the $P_{\rm eq} = 20$\,d example above, where $\mathcal{K}_1$ has died, $\mathcal{K}_2$ still retains $0.43$ of its zero-lag value at half a period). And the ratio of the harmonic weights,
\begin{equation}
    \mathcal{H} \equiv \frac{\langle |c_2|^2 \rangle_\Phi}{\langle |c_1|^2 \rangle_\Phi},
    \label{eq:harmonic_ratio}
\end{equation}
is a clean inclination diagnostic: $\mathcal{H} = (4/(3\pi))^2 \approx 0.18$ for edge-on viewing with uniform latitudes, falling by about $5\times$ the magnitude as $I = 30^\circ$. In the lightcurve, a significant $n = 2$ share manifests as the familiar ``double-dip'' shape: two flux minima per rotation.

The same weights control how deeply the kernel oscillates.\footnote{Numerical check: \NotebookLink{lower_envelope.ipynb}} Writing the extrema of the cosine sum explicitly, the kernel is bounded by $k_{\rm upper}(\tau) = R_\Gamma(\tau)\,[\langle|c_0|^2\rangle + 2\langle|c_1|^2\rangle + 2\langle|c_2|^2\rangle]$ at integer periods and $k_{\rm lower}(\tau) = R_\Gamma(\tau)\,[\langle|c_0|^2\rangle - 2\langle|c_1|^2\rangle + 2\langle|c_2|^2\rangle]$ at half periods, so the fractional oscillation depth
\begin{equation}
    \frac{k_{\rm lower}}{k_{\rm upper}} = \frac{\langle|c_0|^2\rangle - 2\langle|c_1|^2\rangle + 2\langle|c_2|^2\rangle}{\langle|c_0|^2\rangle + 2\langle|c_1|^2\rangle + 2\langle|c_2|^2\rangle}
    \label{eq:envelope_ratio}
\end{equation}
is set entirely by the geometry weights: ${\approx}-0.005$ edge-on (the kernel oscillates essentially to zero at the troughs), $+0.25$ at $I = 60^\circ$, and $+0.77$ at $I = 30^\circ$, where the DC floor dominates and the oscillation is a shallow ripple on an aperiodic decay.

% -----------------------------------------------
\subsubsection{Differential Rotation Across the Harmonics} \label{subsec:harmonic_kappa}

At fixed latitude, the PSD sideband of order $n$ has width proportional to $n\kappa$, and the textbook expectation is that higher orders decohere proportionally faster. The kernel says otherwise, because each order samples latitude through its own weight $|c_n(I,\Phi)|^2$, and the shear $\kappa\sin^2\Phi$ doing the smearing is itself latitude-dependent. Higher orders need spots that set, so $|c_2|^2$ concentrates toward the equator, precisely where the shear is weakest. The two effects run in opposite directions, and which one wins depends on the inclination. Quantifying the dephasing rate of order $n$ by $n\,\sigma_\omega^{(n)}$, where $\sigma_\omega^{(n)}$ is the $|c_n|^2$-weighted spread of $\omega_0(\Phi)$:

\begin{table}[ht!]
    \centering
    \caption{Latitude sampling and dephasing rate of the $n = 1$ and $n = 2$ harmonics for $\kappa = 0.4$ and a uniform latitude distribution. $\langle|\Phi|\rangle_n$ is the $|c_n|^2$-weighted mean absolute latitude; $n\,\sigma_\omega^{(n)}$ [rad\,d$^{-1}$] is the dephasing rate.}
    \label{tab:harmonic_dephasing}
    \begin{tabular}{lccl}
        \hline
        Inclination & $\langle|\Phi|\rangle_1$ / $\langle|\Phi|\rangle_2$ & $n\,\sigma_\omega^{(n)}$, $n=1$ / $n=2$ & Outcome \\
        \hline
        $I = 90^\circ$ & $26.8^\circ$ / $26.8^\circ$ & 0.126 / 0.251 & factor 2, as expected \\
        $I = 60^\circ$ & $31.0^\circ$ / $18.0^\circ$ & 0.139 / 0.140 & cancels almost exactly \\
        $I = 30^\circ$ & $30.7^\circ$ / $\phantom{0}8.4^\circ$ & 0.127 / 0.037 & $n=2$ far more coherent \\
        \hline
    \end{tabular}
\end{table}

Only equator-on does the na\"ive $n$-scaling survive (Table~\ref{tab:harmonic_dephasing}). At the fiducial $I = 60^\circ$ the latitude weighting cancels the $n$-scaling almost exactly, and by $I = 30^\circ$ the second harmonic is the \emph{more} coherent of the two: at $\tau = P_{\rm eq}$ with $\kappa = 0.4$, the $n = 1$ recurrence retains $0.58$ of the $n = 0$ envelope while $n = 2$ retains $0.67$. The same latitude weighting shifts the apparent rates: each order reports the mean rotation rate of the latitudes it samples, $\langle\omega_0\rangle_n = 0.975$, $1.101$, $1.192$\,rad\,d$^{-1}$ for $n = 0, 1, 2$ at $I = 60^\circ$, $\kappa = 0.4$ (the $n = 0$ weighting samples the highest latitudes, the $n = 2$ weighting the lowest). The spread of apparent periods \emph{between} the orders is therefore itself a measurement of the shear, complementary to the peak damping of Section~\ref{subsec:harmonic_n1}.

% -----------------------------------------------
\subsubsection{Which Parameter Moves Which Harmonic} \label{subsec:harmonic_matrix}

% Figure~\ref{fig:harmonic_parameter_tour} varies each hyperparameter about the fiducial model and shows its effect on every harmonic term; Table~\ref{tab:harmonic_matrix} summarizes the pattern. Reading down a column shows what one order can constrain; reading across a row shows where a parameter is identifiable.

% \begin{figure}[ht!]
%     \centering
%     \script{plot_kernel_psd_harmonic_decomposition.py}ro
%     \includegraphics[width=0.95\linewidth]{figures/fig_kernel_psd_harmonic_decomposition.pdf}
%     \caption{Harmonic decomposition of the kernel (left) and PSD (right) as each hyperparameter is varied about the fiducial model ($P_{\rm eq} = 5$\,d, $\kappa = 0.2$, $I = 60^\circ$, $\ell_{\rm spot} = 15$\,d, $\tau_{\rm spot} = 3$\,d, $\sigma_k = 0.05$; parameter values are encoded light $\to$ dark). $P_{\rm eq}$ and $\kappa$ move only the $n \geq 1$ terms and leave the $n = 0$ curve untouched; $I$ redistributes power between the orders without changing any shape; $\ell_{\rm spot}$ and $\tau_{\rm spot}$ stretch the common envelope of all orders; $\sigma_k$ rescales everything equally.}
%     \label{fig:harmonic_parameter_tour}
% \end{figure}

\begin{table}[ht!]
    \centering
    \caption{Which hyperparameter affects which harmonic term, and how. ``Amplitude'' entries change only the term's overall scale; ``shape'' entries change its dependence on $\tau$.}
    \label{tab:harmonic_matrix}
    \begin{tabular}{lccll}
        \hline
        Parameter & $n = 0$ & $n = 1,2$ & Effect \\
        \hline
        $\sigma_k$ & amplitude & amplitude & rescales all orders; shares and shapes unchanged \\
        $\ell_{\rm spot}$ & shape + amp. & shape + amp. & common envelope; support ends at $\ell_{\rm spot} + 2\tau_{\rm spot}$ \\
        $\tau_{\rm spot}$ & shape + amp. & shape + amp. & as $\ell_{\rm spot}$, at roughly half the leverage \\
        $I$ & amplitude & amplitude & redistributes variance between orders; no shape change \\
        $p(\Phi)$ & amplitude & amplitude$^{\dagger}$ & same lever as $I$; degenerate with it \\
        $P_{\rm eq}$ & -- & shape & recurrences at $P_{\rm eq}$ and $P_{\rm eq}/2$ \\
        $\kappa$ & -- & shape & damps recurrences below the $n = 0$ envelope; skews them later \\
        \hline
    \end{tabular}
    \tablecomments{$^{\dagger}$For $\kappa \neq 0$ the latitude distribution also enters the $n \geq 1$ shapes through the sampled range of $\omega_0(\Phi)$ (Section~\ref{subsec:harmonic_kappa}).}
\end{table}

Four consequences for inference follow directly:
\begin{enumerate}[nosep]
    \item \textbf{Spot lifetime and shear are separable, but only jointly.} Both damp the ACF, but only the shear makes the $n \geq 1$ terms decay faster than the $n = 0$ term. The comparison between the orders (not either curve alone) carries $\kappa$.
    \item \textbf{Inclination and the latitude distribution are one lever, not two.} Both act purely through $|c_n|^2$: an equatorial band viewed at $I = 60^\circ$ and a full-latitude star viewed edge-on put nearly identical power in the fundamental ($n = 1$ share $\approx 0.50$ in both cases). Without a latitude prior, a fitted inclination is really a fitted geometry weight (cf.\ the Kepler-63 analysis of Paper~II, where transit-mapped spot latitudes provide exactly such a prior).
    \item \textbf{$\kappa$ is never constrained alone.} Dephasing needs a lever arm in latitude: at $\kappa = 0.4$ an equatorial band barely damps the first recurrence at all (0.712 versus 0.741 for solid body), while a polar band moves it from 5.0 to 7.4\,d. What the data constrain is $\kappa$ times the occupied latitude spread: the same degeneracy noted for individual-spot modeling by \citet{basri_information_2020}.
    \item \textbf{Rotation is measurable only while the envelope survives it.} When $\ell_{\rm spot} + 2\tau_{\rm spot} \lesssim P_{\rm eq}$, the first recurrence lands beyond the envelope's support and $P_{\rm eq}$, $\kappa$, and $I$ all lose their handles at once, leaving only the weak half-period recurrence of $n = 2$. This is the regime, identified independently by the Cram\'er--Rao analysis of Section~\ref{sec:cranmer_rao}, in which the kernel degrades gracefully into a purely aperiodic model.
\end{enumerate}

\vspace{\baselineskip}
% ========================================================================================
\section{GP Solver Implementation} \label{sec:gp_solver}

In order to solve for the GP posterior and perform inference on the hyperparameters $\theta$, we need to evaluate the GP marginal log-likelihood (Eq.~\ref{eq:gp_logL}), which requires computing the kernel matrix inversion. In this section, we show how the structure of our kernel allows for efficient computation using a banded Cholesky solver, which reduces the computational cost from $\mathcal{O}(N^3)$ to $\mathcal{O}(Nb^2)$, where $b$ is the bandwidth determined by the spot lifetime and observing cadence. We implement this solver in JAX, enabling fast automatic differentiation for MCMC sampling of the hyperparameters.

% -----------------------------------------------
\subsection{Naive Cholesky Solver} 

We implement a JAX-accelerated GP solver that uses the analytic kernel for Bayesian inference on stellar lightcurves. The GP marginal log-likelihood for $N$ observations $\mathbf{y}$ at times $\mathbf{x}$ is
\begin{equation}
    \ln p(\mathbf{y} \mid \theta) = -\frac{1}{2}\left[
        (\mathbf{y} - \boldsymbol{\mu})^\top K_\theta^{-1} (\mathbf{y} - \boldsymbol{\mu})
        + \ln |K_\theta|
        + N \ln 2\pi
    \right],
    \label{eq:gp_logL}
\end{equation}
where $K_\theta$ is the $N \times N$ covariance matrix with entries $[K_\theta]_{ij} = k(|x_i - x_j|;\, \theta) + (\sigma_{{\rm err},i}^2 + \sigma_n^2)\,\delta_{ij}$, $\boldsymbol{\mu}$ is the mean vector, and $\sigma_n$ is an optional white-noise amplitude.

Evaluating Eq.~\eqref{eq:gp_logL} involves two computationally distinct steps:
\begin{enumerate}[nosep]
    \item \textit{Kernel evaluation}: computing $k(|x_i - x_j|;\, \theta)$ for all pairs, which requires the latitude integral in Eq.~\eqref{eq:kernel_full}. We perform this integral using Gauss-Legendre quadrature with $n_{\rm lat} = 64$ nodes. To avoid materializing all latitude contributions simultaneously (which would require an $n_{\rm lat} \times M$ intermediate tensor, where $M$ is the number of unique lags), we accumulate the weighted latitude contributions sequentially, reducing peak memory from $\mathcal{O}(n_{\rm lat} \times M)$ to $\mathcal{O}(M)$. We further reduce memory by exploiting the symmetry of the covariance matrix, evaluating the kernel only on the $N(N+1)/2$ upper-triangular lags rather than the full $N^2$ matrix.
    \item \textit{Cholesky factorization and solve}: decomposing $K_\theta = L L^\top$ and solving for $L^{-1}(\mathbf{y} - \boldsymbol{\mu})$, which dominates at $\mathcal{O}(N^3)$.
\end{enumerate}

Our implementation provides two interfaces. The \code{GPSolver} class pre-computes and caches the Cholesky factorization at construction time, enabling fast repeated queries (predictions, log-likelihood evaluation) for fixed hyperparameters. For MCMC sampling, where hyperparameters change at every step, we provide a pure-functional JIT-compiled log-likelihood function that recomputes the full pipeline (kernel evaluation, Cholesky, and solve) from scratch on each call. JAX's just-in-time compilation and automatic differentiation provide efficient gradient computation for Hamiltonian Monte Carlo (HMC) and related samplers.

\begin{table}[ht!]
    \centering
    \caption{GPSolver wall-clock time and peak memory as a function of the number of data points $N$. \textit{Build}: covariance matrix + Cholesky factorization. \textit{Log-likelihood}: JIT-compiled evaluation (kernel + Cholesky + solve) as called during MCMC. \textit{Predict}: predictive mean at 200 test points. Times are the median of 5 runs after JIT warmup. Memory is the peak RSS increase during each operation, measured in an isolated subprocess.}
    \label{tab:gp_solver_speed}
    \begin{tabular}{rcccccc}
    \toprule
    & \multicolumn{2}{c}{Build} & \multicolumn{2}{c}{Log-likelihood} & \multicolumn{2}{c}{Predict} \\
    \cmidrule(lr){2-3} \cmidrule(lr){4-5} \cmidrule(lr){6-7}
    $N$ & Time [s] & Memory & Time [s] & Memory & Time [s] & Memory \\
    \midrule
    10 & 0.0810 & $17\,\text{MB}$ & 0.0001 & $0.1\,\text{MB}$ & 0.1749 & $29\,\text{MB}$ \\
    100 & 0.0900 & $17\,\text{MB}$ & 0.0057 & $0.1\,\text{MB}$ & 0.2116 & $29\,\text{MB}$ \\
    1000 & 0.3546 & $137\,\text{MB}$ & 0.1416 & $0.1\,\text{MB}$ & 0.2892 & $38\,\text{MB}$ \\
    5000 & 8.4480 & $3.57\,\text{GB}$ & 3.0122 & $503\,\text{MB}$ & 0.7084 & $246\,\text{MB}$ \\
    10000 & 31.9826 & $14.16\,\text{GB}$ & 12.3798 & $1.97\,\text{GB}$ & 1.7840 & $882\,\text{MB}$ \\
    \bottomrule
    \end{tabular}
\end{table}

Table~\ref{tab:gp_solver_speed} benchmarks the wall-clock time and peak memory usage of the GP solver as a function of the number of data points $N$. The ``Build'' column measures the \texttt{GPSolver} construction (kernel evaluation + Cholesky), while the ``Log-likelihood'' column measures the JIT-compiled pure-functional evaluation used during MCMC. After JIT warmup, the log-likelihood evaluation is faster than the class-based build because XLA reuses pre-allocated buffers. At $N = 1000$, the full log-likelihood evaluation takes ${\sim}0.1$\,s on a single CPU core, making MCMC chains of $\sim$10$^4$ samples feasible within hours. The $\mathcal{O}(N^2)$ memory scaling of the covariance matrix is the primary limitation for larger datasets.

% -----------------------------------------------
\subsection{Banded Cholesky Solver}

The compact support of the spot-lifetime envelope $R_\Gamma(\tau)$ (Section~\ref{sec:hyperparameters}) guarantees that the kernel vanishes exactly for lags $|\tau| > \tau_{\max} \equiv \ell_{\rm spot} + 2\tau_{\rm spot}$.  For uniformly sampled data with cadence $\Delta t$, this means the covariance matrix $K_\theta$ is banded with bandwidth
\begin{equation}
    b = \left\lceil \frac{\ell_{\rm spot} + 2\tau_{\rm spot}}{\Delta t} \right\rceil,
    \label{eq:bandwidth}
\end{equation}
i.e., $[K_\theta]_{ij} = 0$ whenever $|i - j| > b$.  For typical stellar monitoring campaigns where the observing baseline greatly exceeds the spot lifetime ($T_{\rm obs} \gg \tau_{\max}$), the bandwidth satisfies $b \ll N$, and the covariance matrix is sparse.

We exploit this banded structure to reduce the cost of the Cholesky factorization from $\mathcal{O}(N^3)$ to $\mathcal{O}(Nb^2)$ and the memory from $\mathcal{O}(N^2)$ to $\mathcal{O}(Nb)$ \citep[Section~4.3.5]{golub_matrix_2013}.  The lower Cholesky factor $L$ inherits the banded structure of $K_\theta$ (i.e., $L_{ij} = 0$ for $i - j > b$), so both the factorization and the subsequent forward/backward substitution can be restricted to within the band.

We store $L$ in compact banded form as a $(b+1) \times N$ array $\tilde{L}$, where
\begin{equation}
    \tilde{L}_{d,\,j} = L_{j+d,\,j}, \qquad d = 0, 1, \ldots, b,
    \label{eq:banded_storage}
\end{equation}
so that the diagonal occupies row $d = 0$ and the $d$-th sub-diagonal occupies row $d$.  The kernel is evaluated directly into this compact storage, computing only the $\mathcal{O}(Nb)$ entries within the band rather than the full $N^2$ matrix.  The log-determinant is then $\ln |K_\theta| = 2 \sum_{j} \ln \tilde{L}_{0,j}$.

\begin{figure}[ht!]
    \centering
    \script{plot_banded_covariance.py}
    \includegraphics[width=\linewidth]{figures/fig_banded_covariance.pdf}
    \caption{Scaling of the GP solver with bandwidth $b$ for fixed $N = 1000$. The banded Cholesky solver scales as $\mathcal{O}(Nb^2)$, while the full Cholesky scales as $\mathcal{O}(N^3)$. For typical spot lifetimes and observing cadences, the bandwidth is small ($b \ll N$), leading to significant speedups.}
    \label{fig:banded_covariance}
\end{figure}

\begin{table}[ht!]
    \centering
    \caption{GPSolver wall-clock time: \texttt{cholesky\_full} vs.\ \texttt{cholesky\_banded} (JAX backend). \textit{Build}: kernel evaluation + Cholesky factorization. \textit{Predict}: GP predictive mean at 200 test points. Times are the median of 3 runs after JIT warmup. The full Cholesky solver scales as $\mathcal{O}(N^3)$ and becomes impractical for $N \gtrsim 10^4$; the banded solver scales as $\mathcal{O}(Nb^2)$ and handles $N = 10^5$ in seconds.}
    \label{tab:banded_vs_full}
    \begin{tabular}{rr cc cc c}
    \toprule
    & & \multicolumn{2}{c}{Full Cholesky} & \multicolumn{2}{c}{Banded Cholesky} & \\
    \cmidrule(lr){3-4} \cmidrule(lr){5-6}
    $x_{\max}$ [d] & $N$ & Build & Predict & Build & Predict & Speedup \\
    \midrule
    10 & 100 & $333\,\text{ms}$ & $493\,\text{ms}$ & $535\,\text{ms}$ & $589\,\text{ms}$ & $0.6\times$ \\
    50 & 500 & $458\,\text{ms}$ & $669\,\text{ms}$ & $587\,\text{ms}$ & $639\,\text{ms}$ & $0.8\times$ \\
    100 & 1{,}000 & $732\,\text{ms}$ & $761\,\text{ms}$ & $595\,\text{ms}$ & $734\,\text{ms}$ & $1.2\times$ \\
    500 & 5{,}000 & $16.9\,\text{s}$ & $2.66\,\text{s}$ & $810\,\text{ms}$ & $773\,\text{ms}$ & $20.8\times$ \\
    1{,}000 & 10{,}000 & $36.2\,\text{s}$ & $9.55\,\text{s}$ & $957\,\text{ms}$ & $1.02\,\text{s}$ & $37.8\times$ \\
    5{,}000 & 50{,}000 & --- & --- & $2.52\,\text{s}$ & $2.52\,\text{s}$ & --- \\
    10{,}000 & 100{,}000 & --- & --- & $4.47\,\text{s}$ & $4.56\,\text{s}$ & --- \\
    \bottomrule
    \end{tabular}
\end{table}


We implement the banded Cholesky factorization using \code{jax.lax.scan} over columns: at each column $j$, the diagonal and the $b$ sub-diagonal entries are updated using only the $(b \times b)$ block of previously computed entries within the band.  The entire factorization is JIT-compiled and end-to-end differentiable via JAX's automatic differentiation, enabling gradient-based sampling with no additional implementation effort.  The forward and backward triangular solves ($L\mathbf{y} = \mathbf{r}$ and $L^\top \mathbf{x} = \mathbf{y}$) similarly operate on the compact storage, restricting each dot product to at most $b$ terms.


% % -----------------------------------------------
% \subsection{Software Architecture} \label{sec:software}

% \begin{figure*}[ht!]
%     \centering
%     \includegraphics[width=\linewidth]{figures/architecture.pdf}
%     \caption{Architecture of the \code{spotgp} package. The core components are the \texttt{Kernel} class, which computes the covariance matrix from the kernel function; the \texttt{GPSolver} class, which performs Cholesky factorization and solves for the GP posterior; and the MCMC sampling module, which uses NUTS to sample the kernel hyperparameters. The kernel is constructed from modular components: the envelope autocorrelation $R_\Gamma(\tau)$, the amplitude parameterization (e.g., Fourier coefficients), and the latitude distribution $p(\Phi)$. Each component can be independently modified or extended without affecting the overall architecture.}
%     \label{fig:spotgp_flowchart}
% \end{figure*}

We implement the kernel derivation and GP inference framework described above in the open-source Python package \code{spotgp}. The package is designed with a modular architecture so that each physical assumption (envelope shape, amplitude parameterization, latitude distribution) can be swapped independently without modifying the kernel or solver code. For a detailed description of the code modules and how to implement custom kernel components, we refer the reader to our software paper \xxx{Birky et al. 2026 [add JOSS citation]}.

% Figure~\ref{fig:spotgp_flowchart} illustrates the data flow through the main components, which we describe below.

% % ------------------------------------------------------------------
% \subsection{Custom Model Components} \label{sec:custom_components}

% The modularity of the implementation mirrors the structure of the kernel itself: Eq.~\eqref{eq:kernel_full} consumes each physical ingredient only through a narrow interface: the envelope autocorrelation $R_\Gamma(\tau)$, the squared visibility coefficients $|c_n(I,\Phi)|^2$ with the rotation law $\omega_0(\Phi)$, and the latitude weights $p(\Phi)$, so each ingredient is represented by an abstract base class, and swapping a physical assumption means passing a different subclass to \code{SpotEvolutionModel}. Neither the kernel evaluation, the latitude quadrature, nor the solver needs to know which concrete model is in use.

% \begin{itemize}[leftmargin=2em]
%     \item \textbf{Envelope functions} subclass \code{EnvelopeFunction} and provide the normalized squared spot-size profile $\Gamma(t)$ together with a \code{param\_dict} naming the free shape parameters (e.g., $\ell_{\rm spot}$, $\tau_{\rm spot}$), which are thereby exposed to the sampler. A subclass may supply a closed-form, JAX-traceable autocorrelation $R_\Gamma(\tau)$; if it does not, the package falls back to evaluating $R_\Gamma$ numerically from $\Gamma(t)$ via FFT, so any profile shape is usable at the cost of the analytic fast path.
%     \item \textbf{Latitude distribution functions} subclass \code{LatitudeDistributionFunction} and override the density $p(\Phi)$ and, optionally, the latitude range and free parameters. Distributions with free parameters (e.g., the edges of an active-latitude band) are sampled alongside the kernel hyperparameters; the latitude quadrature then applies the current weights at every likelihood evaluation.
%     \item \textbf{Visibility functions} subclass \code{VisibilityFunction}, which provides the analytic small-spot coefficients of Eq.~\eqref{eq:cn_general} and the differential rotation law $\omega_0(\Phi)$. A subclass overrides \code{cn\_squared} and, to participate in gradient-based fitting, defines the \code{cn\_sq\_jax} hook: a JAX-traceable map from the visibility's slice of the parameter vector to $|c_n|^2$, which the kernel detects and routes the latitude-averaged covariance through automatically. Coefficients without a closed form (limb darkening beyond the quadratic law, finite spot size) are evaluated by sampling the visibility profile over one rotation and taking its discrete Fourier transform, as described in Paper~II.
% \end{itemize}

% The base \code{VisibilityFunction} and \code{LatitudeDistributionFunction} classes are themselves usable defaults (the analytic small-spot visibility and the uniform latitude distribution); \code{EnvelopeFunction} is abstract. The concrete implementations of each component type available in the current release of \code{spotgp}, together with the location where each is derived, are tabulated in Paper~II, whose composite kernels (the \code{Term} layer: several spot populations sharing one geometry, or additive non-spot components) consume these components through the same interfaces, so custom components propagate to those models without modification.


% ------------------------------------------------------------------
% \vspace{\baselineskip}
% ========================================================================================
% \vspace{\baselineskip}
\clearpage
% ========================================================================================
\section{MCMC Sampling the Kernel Hyperparameters} \label{sec:mcmc}

Figure~\ref{fig:pgm} summarizes the probabilistic structure of the model, showing how the physical parameters determine the photometric kernel and how the flux observations derive from it.

\begin{figure*}
\centering
\begin{tikzpicture}[
  param/.style={circle, draw, semithick, minimum size=10mm, inner sep=0pt, font=\small},
  obs/.style={circle, draw, semithick, fill=gray!20, minimum size=10mm, inner sep=0pt, font=\small},
  det/.style={rectangle, draw, semithick, rounded corners=2pt, minimum height=9mm, inner sep=3pt, font=\small},
  plate/.style={draw, semithick, dashed, rounded corners=3pt, inner sep=6pt},
  arr/.style={-{Stealth[length=4pt]}, thin},
]

% === Row 1: Physical parameters (y=0) ===
\node[param, minimum size=11mm] (Peq) at (0, 0)    {\scriptsize $P_{\rm eq}$};
\node[param] (kap) at (2.2, 0)   {\scriptsize $\kappa$};
\node[param] (inc) at (4.0, 0)   {\scriptsize $I$};
\node[param] (ell) at (6.5, 0)   {\scriptsize $\ell_{\rm sp}$};
\node[param] (tau) at (8.5, 0)   {\scriptsize $\tau_{\rm sp}$};
\node[param] (sig) at (11.5, 0)  {\scriptsize $\sigma_k$};

% Group labels
\node[font=\scriptsize\itshape, text=gray!60!black] at (2.1, 1.05)  {rotation};
\node[font=\scriptsize\itshape, text=gray!60!black] at (7.5, 1.05)  {evolution};
\node[font=\scriptsize\itshape, text=gray!60!black] at (11.5, 1.05) {amplitude};

% Brackets around parameter groups
\draw[thin, gray!60!black] (-0.55, 0.55) -- (-0.55, 0.72) -- (4.55, 0.72) -- (4.55, 0.55);
\draw[thin, gray!60!black] (5.95, 0.55) -- (5.95, 0.72) -- (9.05, 0.72) -- (9.05, 0.55);
\draw[thin, gray!60!black] (10.95, 0.55) -- (10.95, 0.72) -- (12.05, 0.72) -- (12.05, 0.55);

% Fixed nodes (to the left of the kernel block)
\node[circle, fill=black, minimum size=3pt, inner sep=0pt] (pPhi) at (2.0, -2.4) {};
\node[font=\scriptsize, anchor=east] at (1.8, -2.4) {$p(\Phi)$};
\node[circle, fill=black, minimum size=3pt, inner sep=0pt] (ulimb) at (2.0, -3.2) {};
\node[font=\scriptsize, anchor=east] at (1.8, -3.2) {$u$};

% === Row 2: Photometric kernel (y=-2.8) ===
\node[det] (K) at (5.5, -2.8) {$\mathcal{K}_{\rm spot}(\tau)$};

\foreach \p in {Peq, kap, inc, ell, tau, sig}
  \draw[arr] (\p) -- (K);
\draw[arr] (pPhi) -- (K);
\draw[arr] (ulimb) -- (K);

% === Row 3: Observable (y=-5.5) ===
\node[param] (sn)  at (3.0, -5.5) {\footnotesize $\sigma_n$};
\node[obs]   (F)   at (5.5, -5.5) {\footnotesize $F_i$};
\node[plate, fit=(F), label={[font=\tiny]below right:{$N_t$}}] {};

\draw[arr] (K) -- (F);
\draw[arr] (sn) -- (F);

% === Node-type legend ===
\node[param, minimum size=6mm]  (L1) at (1.0, -7.5) {};
\node[font=\scriptsize, anchor=west] at (1.45, -7.5) {free parameter};
\node[obs, minimum size=6mm]    (L2) at (4.5, -7.5) {};
\node[font=\scriptsize, anchor=west] at (4.95, -7.5) {observed};
\node[det, minimum height=6mm, inner sep=2pt] (L3) at (8.0, -7.5) {};
\node[font=\scriptsize, anchor=west] at (8.45, -7.5) {deterministic};
\node[circle, fill=black, minimum size=3pt, inner sep=0pt] (L4) at (11.5, -7.5) {};
\node[font=\scriptsize, anchor=west] at (11.7, -7.5) {fixed};

% === Variable definitions ===
\draw[gray!40, thin] (-0.5, -8.3) -- (13.0, -8.3);

\node[font=\scriptsize, anchor=west, text=gray!70!black] at (-0.5, -8.8)  {$P_{\rm eq}$: equatorial rotation period};
\node[font=\scriptsize, anchor=west, text=gray!70!black] at (6.5, -8.8)   {$\sigma_k$: variability amplitude};

\node[font=\scriptsize, anchor=west, text=gray!70!black] at (-0.5, -9.3)  {$\kappa$: diff.\ rotation shear};
\node[font=\scriptsize, anchor=west, text=gray!70!black] at (6.5, -9.3)   {$\sigma_n$: photometric noise};

\node[font=\scriptsize, anchor=west, text=gray!70!black] at (-0.5, -9.8)  {$I$: stellar inclination};
\node[font=\scriptsize, anchor=west, text=gray!70!black] at (6.5, -9.8)   {$p(\Phi)$: spot latitude distribution};

\node[font=\scriptsize, anchor=west, text=gray!70!black] at (-0.5, -10.3) {$\ell_{\rm sp}$: spot active lifetime};
\node[font=\scriptsize, anchor=west, text=gray!70!black] at (6.5, -10.3)  {$u$: limb darkening coefficient};

\node[font=\scriptsize, anchor=west, text=gray!70!black] at (-0.5, -10.8) {$\tau_{\rm sp}$: spot decay timescale};

\node[font=\scriptsize, anchor=west, text=gray!70!black] at (-0.5, -11.3) {$F_i$: photometric flux};

\end{tikzpicture}
\caption{Probabilistic graphical model for the photometric starspot GP. Free parameters (open circles) determine the photometric kernel $\mathcal{K}_{\rm spot}(\tau)$ (Equation~\ref{eq:kernel_full}), which governs the temporal covariance of the flux observations $F_i$ across $N_t$ epochs. The rotation parameters ($P_{\rm eq}$, $\kappa$, $I$) control the harmonic structure, the evolution parameters ($\ell_{\rm sp}$, $\tau_{\rm sp}$) control the temporal envelope $R_\Gamma(\tau)$, and the amplitude parameter $\sigma_k$ controls the overall variability level. The spot latitude distribution $p(\Phi)$ and limb darkening coefficient $u$ (filled dots) enter the kernel as fixed inputs from stellar models.}
\label{fig:pgm}
\end{figure*}

Given a set of observations $\{t_i, y_i, \sigma_i\}_{i=1}^N$, we sample the posterior distribution of the kernel hyperparameters $\boldsymbol{\theta} = (P_{\rm eq}, \kappa, I, \ell_{\rm spot}, \tau_{\rm spot}, \sigma_k)$ using the No-U-Turn Sampler \citep[NUTS;][]{hoffman_no-u-turn_2014}, a variant of Hamiltonian Monte Carlo (HMC) that automatically tunes the trajectory length. We implement this using the \code{BlackJAX} library \citep{lao_blackjax_2020}, which provides a JAX-native NUTS kernel that is fully compatible with the JIT-compiled, autodifferentiable log-posterior described below.

% -----------------------------------------------
\subsection{Log-Posterior and Gradients}

The GP marginal log-likelihood is
\begin{equation}
    \ln \mathcal{L}(\boldsymbol{\theta}) = -\frac{1}{2}\left[\mathbf{r}^\top K_\theta^{-1} \mathbf{r} + \ln|K_\theta| + N\ln(2\pi)\right],
    \label{eq:gp_loglik}
\end{equation}
where $\mathbf{r} = \mathbf{y} - \boldsymbol{\mu}$ is the residual vector, $K_\theta = K(\boldsymbol{\theta}) + \mathrm{diag}(\boldsymbol{\sigma}^2 + \sigma_n^2) + \epsilon \mathbf{I}$ is the noise-augmented covariance matrix, and $\epsilon = 10^{-8}$ is a small jitter for numerical stability. The Cholesky factorization $K_\theta = LL^\top$ is used for efficient computation of both the data-fit term and the log-determinant: $\ln|K_\theta| = 2\sum_i \ln L_{ii}$.

The entire computational graph, from hyperparameters $\boldsymbol{\theta}$ through the kernel evaluation (Eq.~\ref{eq:kernel_full}), covariance matrix assembly, Cholesky decomposition, and log-likelihood, is implemented in JAX and is end-to-end differentiable. The gradient $\nabla_{\boldsymbol{\theta}} \ln p(\boldsymbol{\theta} | \mathbf{y})$ is computed via reverse-mode automatic differentiation, providing exact gradients at roughly the cost of a single likelihood evaluation.

% -----------------------------------------------
\subsection{Prior}

We adopt a soft uniform prior that enforces parameter bounds while maintaining differentiability. For each parameter $\theta_i$ with bounds $[a_i, b_i]$:
\begin{equation}
    \ln p(\boldsymbol{\theta}) = \sum_i \left[\ln\sigma(k(\theta_i - a_i)) + \ln\sigma(k(b_i - \theta_i)) - \ln(b_i - a_i)\right],
\end{equation}
where $\sigma(\cdot)$ is the sigmoid function and $k=100$ controls the steepness of the boundary (Figure~\ref{fig:prior}). This provides effectively uniform support within the bounds while keeping the gradient finite everywhere, which is critical for HMC-based samplers.

\begin{figure}
    \centering
    \script{plot_soft_uniform_prior.py}
    \includegraphics[width=\linewidth]{figures/fig_soft_uniform_prior.pdf}
    \caption{Soft uniform prior for a single parameter $\theta$ with bounds $[a, b]$. The sigmoid-based formulation provides a smooth transition at the boundaries while maintaining an approximately uniform density within the bounds.}
    \label{fig:prior}
\end{figure}

\subsection{Mass Matrix Estimation} \label{subsec:mass_matrix}

The efficiency of NUTS depends on the choice of mass matrix $M$, which preconditions the Hamiltonian dynamics to account for differences in the scale and correlation of the parameters. An identity mass matrix performs poorly when the hyperparameters span vastly different scales (e.g., $P_{\rm eq} \sim 5$~d vs.\ $\alpha_{\max} \sim 0.1$~rad). We provide three methods for estimating the inverse mass matrix $M^{-1}$:

\begin{enumerate}
    \item \textbf{Hessian at the MAP.} We first find the maximum a posteriori estimate $\boldsymbol{\theta}_{\rm MAP}$ via L-BFGS-B using JAX-computed gradients, then compute the Hessian of the negative log-likelihood at that point using \code{jax.hessian}:
    \begin{equation}
        M^{-1} = \left[\nabla^2_{\boldsymbol{\theta}}(-\ln\mathcal{L})\Big|_{\boldsymbol{\theta}_{\rm MAP}}\right]^{-1}.
    \end{equation}
    The Hessian is regularized by clamping its eigenvalues to be at least $10^{-6}$, ensuring positive definiteness.

    \item \textbf{Fisher information.} The expected Fisher information matrix for a GP with covariance $K_\theta$ is
    \begin{equation}
        \mathcal{I}_{ij} = \frac{1}{2}\mathrm{tr}\left(K_\theta^{-1}\frac{\partial K_\theta}{\partial\theta_i} K_\theta^{-1}\frac{\partial K_\theta}{\partial\theta_j}\right).
        \label{eq:fisher}
    \end{equation}
    The kernel derivatives $\partial K_\theta / \partial \theta_i$ are computed via forward-mode autodiff (\code{jax.jacfwd}), which is memory-efficient for the case of few parameters mapping to a large covariance matrix. The inverse mass matrix is then $M^{-1} = \mathcal{I}^{-1}$.

    \item \textbf{Laplace approximation.} This is equivalent to the Hessian method but additionally provides a Gaussian approximation to the posterior, $p(\boldsymbol{\theta}|\mathbf{y}) \approx \mathcal{N}(\boldsymbol{\theta}_{\rm MAP}, M^{-1})$, which can be used for rapid approximate inference or as a proposal distribution.
\end{enumerate}

In practice, for NUTS we extract the diagonal of the full inverse mass matrix and use a diagonal mass matrix, which is more numerically robust. Extreme values are clamped to within a factor of $10^4$ of the median diagonal element to prevent instability.

\subsection{Sampler Backends} \label{subsec:sampler_backends}

Because the \code{spotgp} likelihood and its gradients are implemented as pure JAX functions, they can be readily coupled to any sampler that accepts a callable log-probability (and, optionally, its gradient). We currently provide built-in interfaces to two complementary backends:

\begin{itemize}
    \item \textbf{BlackJAX NUTS.} The default sampler is the NUTS kernel from \code{BlackJAX} \citep{lao_blackjax_2020}, which exploits the exact gradients provided by JAX autodifferentiation to efficiently explore the posterior. This is the recommended choice when the posterior is unimodal and the mass matrix estimation described in Section~\ref{subsec:mass_matrix} provides adequate preconditioning.

    \item \textbf{Dynesty nested sampling.} For problems where the posterior may be multimodal or where one additionally requires a Bayesian evidence estimate, \code{spotgp} also supports the \code{dynesty} dynamic nested sampler \citep{speagle_dynesty_2020}. Because nested sampling does not require gradient information, this backend uses only the log-likelihood and prior transform, making it a useful complement to NUTS when gradient-based sampling encounters difficulties such as highly correlated or multimodal posteriors.
\end{itemize}

Users may select the sampler backend at runtime, and the interface is designed so that additional samplers can be integrated with minimal effort by supplying the appropriate log-probability and prior functions.


\vspace{\baselineskip}
% ========================================================================================
\section{Numerical Results} \label{sec:results}

\subsection{Analytic vs. Numerical Kernel} \label{sec:analytic_vs_numerical}

In this Section, we compare the analytic kernel derived in Section~\ref{sec:kernel_derivation} to a numerical kernel computed by simulating many random spot realizations and empirically estimating their autocorrelation (Eq.~\ref{eq:numerical_kernel}). For the numerical kernel, we simulate $N_{\rm sim} = 500$ realizations of the flux deficit $\delta F(t)$, each with $N_{\rm spot} = 10$ spots drawn from the same parameter distribution, computed with a cadence of 0.05 days for a duration of 200 days. We then compute the sample autocorrelation function across all simulations to obtain an empirical estimate of the kernel. For the analytic kernel, we evaluate Eq.~\eqref{eq:kernel_full} numerically for the same set of parameters.

Table~\ref{tab:kernel_mse} summarizes the parameter combinations tested and the mean squared error (MSE) between the analytic and numerical autocorrelation functions. For each row, a single parameter is varied from its baseline value ($P_{\rm eq} = 5$~d, $\kappa = 0$, $I = 90^\circ$, $\ell_{\rm spot} = 15$~d, $\tau_{\rm spot} = 3$~d, $\alpha_{\max} = 0.1$~rad), with $N_{\rm spot} = 10$ held fixed throughout. 

\variable{output/kernel_accuracy_benchmark.tex}
\begin{table}[ht!]
    \centering
    \caption{A comparison showing the kernel computation time (in seconds) using the analytic and numerical approaches. The analytic kernel is evaluated at varying lag-array sizes (with $N$ simulated datapoints); the numerical (Monte Carlo) kernel uses 200 and 500 simulations over 100~d.}
    \label{tab:kernel_speed}
    \begin{tabular}{rccccc}
    \toprule
    & & \multicolumn{2}{c}{$N_{\rm sim}=200$} & \multicolumn{2}{c}{$N_{\rm sim}=500$} \\
    \cmidrule(lr){3-4} \cmidrule(lr){5-6}
    $N$ & Analytic [s] & Numerical [s] & Speedup & Numerical [s] & Speedup \\
    \midrule
    50 & 0.0053 & 0.23 & 42$\times$ & 0.33 & 62$\times$ \\
    100 & 0.0048 & 0.22 & 46$\times$ & 0.34 & 71$\times$ \\
    500 & 0.0050 & 0.26 & 53$\times$ & 0.47 & 94$\times$ \\
    1000 & 0.0054 & 0.36 & 66$\times$ & 0.48 & 89$\times$ \\
    5000 & 0.0075 & 0.50 & 67$\times$ & 0.86 & 114$\times$ \\
    \bottomrule
    \end{tabular}
\end{table}

Figures~\ref{fig:kernel_vary_grid} shows the analytic vs.\ numerical kernel comparison for each parameter variation. In all cases, the analytic kernel closely matches the numerical kernel, with MSE values on the order of $10^{-3}$ or smaller. Overall, these results validate our analytic kernel against numerical simulations across a range of physically relevant parameters.

Table~\ref{tab:kernel_speed} benchmarks the computational speed of evaluating the analytic kernel vs.\ computing the numerical kernel. The analytic kernel is at least an order of magnitude faster to evaluate than the numerical kernel, which requires simulating many spot realizations and computing their autocorrelation. This speed advantage makes the analytic kernel practical for use in GP regression on large datasets, where repeated kernel evaluations are required.

\begin{figure}[p]
    \centering
    \script{plot_analytic_vs_numerical_kernel_single_panel.py}
    \includegraphics[width=0.95\linewidth]{figures/fig_kernel_comparison_grid.pdf}
    \caption{Analytic (dashed) vs.\ numerical Monte Carlo (solid) kernel comparison, parameter sweeps with varying $P_{\rm eq}$ (upper left), $\kappa$ (upper right), inclination (middle left), $\ell_{\rm spot}$ (middle right), and $\tau_{\rm spot}$ (lower left). The numerical kernel is computed using $N_{\rm sim} = 500$ lightcurve simulations with a baseline of 200d, and cadence of 0.05d. Vertical dotted lines mark the rotation frequency and its first harmonic.}
    \label{fig:kernel_vary_grid}
\end{figure}

% \begin{figure}[ht!]
%     \centering
%     \script{analytic_vs_numerical_kernel.py}
%     \includegraphics[width=0.95\linewidth]{figures/fig_kernel_comparison_vary_peq.pdf}
%     \caption{Analytic (dashed) vs.\ numerical Monte Carlo (solid) kernel comparison, varying $P_{\rm eq} \in \{3, 5, 10\}$~d. The numerical kernel is computed using $N_{\rm sim} = 500$ simulations. \textit{Left:} Autocorrelation function. \textit{Right:} Power spectral density. Vertical dotted lines mark the rotation frequency and its first harmonic.}
%     \label{fig:kernel_vary_peq}
% \end{figure}

% \begin{figure}[ht!]
%     \centering
%     \script{analytic_vs_numerical_kernel.py}
%     \includegraphics[width=0.95\linewidth]{figures/fig_kernel_comparison_vary_kappa.pdf}
%     \caption{Same as Fig.~\ref{fig:kernel_vary_peq}, varying differential rotation shear $\kappa \in \{-1.0, 0.0, 1.0\}$.}
%     \label{fig:kernel_vary_kappa}
% \end{figure}

% \begin{figure}[ht!]
%     \centering
%     \script{analytic_vs_numerical_kernel.py}
%     \includegraphics[width=0.95\linewidth]{figures/fig_kernel_comparison_vary_inc.pdf}
%     \caption{Same as Fig.~\ref{fig:kernel_vary_peq}, varying stellar inclination $I \in \{30^\circ, 60^\circ, 90^\circ\}$.}
%     \label{fig:kernel_vary_inc}
% \end{figure}

% \begin{figure}[ht!]
%     \centering
%     \script{analytic_vs_numerical_kernel.py}
%     \includegraphics[width=0.95\linewidth]{figures/fig_kernel_comparison_vary_lspot.pdf}
%     \caption{Same as Fig.~\ref{fig:kernel_vary_peq}, varying spot plateau duration $\ell_{\rm spot} \in \{5, 15, 30\}$~d.}
%     \label{fig:kernel_vary_lspot}
% \end{figure}

% \begin{figure}[ht!]
%     \centering
%     \script{analytic_vs_numerical_kernel.py}
%     \includegraphics[width=0.95\linewidth]{figures/fig_kernel_comparison_vary_tau_spot.pdf}
%     \caption{Same as Fig.~\ref{fig:kernel_vary_peq}, varying spot rise/decay timescale $\tau_{\rm spot} \in \{1, 3, 8\}$~d.}
%     \label{fig:kernel_vary_tau}
% \end{figure}

\clearpage
% -----------------------------------------------
\subsection{Comparison to Kernels from Literature} \label{sec:comparison_literature}

In this Section, we compare our physically motivated kernel to commonly used kernels in the literature, such as the quasi-periodic (QP) kernel and the simple harmonic oscillator (SHO) kernel. We show that our kernel can reproduce the behavior of these kernels in certain limits, but also captures additional features such as harmonic structure and non-sinusoidal periodicity that are not present in the QP or SHO kernels.

Quasi-periodic kernels typically have the form:
\begin{equation}
    k_{\rm QP}(\tau) = A^2 \exp\left(-\frac{\tau^2}{2\lambda^2}\right) \exp\left(-\frac{2\sin^2(\pi\tau/P)}{\omega^2}\right),
\end{equation}
where $P$ is the period, $\lambda$ is the decay timescale, and $\omega$ controls the smoothness of the periodic component. In the limit of large $\omega$, the periodic component approaches a pure cosine, and the kernel resembles a product of an exponential decay and a cosine. 

The simple harmonic oscillator kernel has the form:
\begin{equation}
    k_{\rm SHO}(\tau) = A^2 \exp\left(-\frac{\tau}{\lambda}\right) \cos\left(\frac{2\pi\tau}{P}\right),
\end{equation}
which is a damped cosine with a single period and an exponential decay.

\begin{figure}[p]
    \centering
    \script{plot_kernel_literature_comparison.py}
    \includegraphics[width=\textwidth]{figures/fig_kernel_literature_comparison.pdf}
    \caption{Comparison of the starspot kernel (this work) with the quasi-periodic (QP) and simple harmonic oscillator (SHO) kernels. \textbf{Left:} Normalized autocorrelation functions. \textbf{Right:} Power spectral densities. \textbf{Top row:} Edge-on geometry with short-lived spots ($P_{\rm eq}=5$\,d, $I=90^\circ$, $\ell=10$\,d, $\tau=3$\,d). \textbf{Bottom row:} Inclined geometry with long-lived spots ($P_{\rm eq}=5$\,d, $I=60^\circ$, $\ell=25$\,d, $\tau=5$\,d). Vertical dotted lines in the PSD panels mark integer harmonics of the rotation frequency. The starspot kernel naturally produces harmonic peaks and non-sinusoidal ACF shapes that the QP and SHO kernels cannot capture.}
    \label{fig:kernel_literature_comparison}
\end{figure}

Figure~\ref{fig:kernel_literature_comparison} compares the normalized autocorrelation function (left) and power spectral density (right) of our starspot kernel against the QP and SHO kernels for three different parameter regimes. The top panel shows a case with short-lived spots and edge-on geometry. Here, the starspot kernel exhibits a rapidly decaying ACF with prominent harmonic peaks in the PSD, while the QP and SHO kernels show a more sinusoidal ACF and a single dominant peak at the rotation frequency. 
The second panel shows a case with long-lived spots and inclined geometry, where the starspot kernel has a more complex ACF shape and multiple harmonic peaks in the PSD, which are not captured by the QP or SHO kernels.

The third panel shows a case with strong solar-like differential rotation ($\kappa = 1$). Here, the starspot kernel is able to capture the ``skewed'' behavior around the ACF peaks (i.e., slower moving spots at higher latitudes have a longer period, thus skewing the ACF around the equitorial period peaks to slightly longer lags). For the QP and SHO kernels, shifts in the ACF peaks can only be made by adjusting the single period parameter, which cannot account for the range of rotation periods due to differential rotation. Differential rotation in the starspot kernel also produces a broadened peak in the PSD around the rotation frequency, reflecting the range of rotation frequencies across different latitudes. In contrast, the QP and SHO kernels cannot capture these features, as they assume a single period and do not account for differential rotation. 

Overall, this comparison demonstrates that our physically motivated starspot kernel can capture a richer variety of behaviors in the autocorrelation and power spectrum than commonly used kernels, making it a more flexible, interpretable (and potentially more accurate) choice for modeling stellar variability due to starspots. 

% -----------------------------------------------
\subsection{Comparison to Other GP Packages} \label{sec:comparison_packages}

While our physically motivated kernel provides richer structure than standard QP or SHO kernels, practical adoption depends on computational performance. We benchmark the \code{spotgp} GP solver against two widely used packages: \code{tinygp} \citep{foreman-mackey_tinygp_2024}, which implements the \code{celerite} $\mathcal{O}(N)$ semi-separable solver \citep{foreman-mackey_fast_2017}, and \code{george} \citep{ambikasaran_fast_2015}, which uses a standard $\mathcal{O}(N^3)$ Cholesky factorization.

For a fair comparison, we use a kernel that is representable in all three frameworks: the one-sided exponential envelope $\Gamma(t) = e^{-t/\tau_{\rm spot}}\,H(t)$, which gives $R_\Gamma(\tau) = (2a)^{-1}\,e^{-a|\tau|}$ with $a = 1/\tau_{\rm spot}$. With $N_{\rm harm} = 2$ harmonics, this decomposes into $J = 3$ \code{celerite} terms (one real + two complex), making it compatible with the $\mathcal{O}(N)$ solver while also being evaluable by the standard dense and banded solvers.

Figure~\ref{fig:benchmark_packages} shows the wall-clock time for a single log-likelihood evaluation as a function of the number of datapoints $N$, with fixed cadence $\Delta t = 1$\,d. We compare four solver configurations:
(1) \code{tinygp} with the \code{celerite} $\mathcal{O}(N)$ solver,
(2) \code{george} with a dense $\mathcal{O}(N^3)$ Cholesky,
(3) \code{spotgp} with a banded $\mathcal{O}(Nb^2)$ Cholesky, and
(4) \code{spotgp} with a full $\mathcal{O}(N^3)$ Cholesky.
All JAX-based solvers (\code{tinygp} and \code{spotgp}) are JIT-compiled and run on the GPU, while \code{george} uses compiled C/LAPACK and can only run on the CPU. 

The top panel of Figure~\ref{fig:benchmark_packages} shows that \code{tinygp}'s \code{celerite} solver scales linearly with $N$ and has the lowest overhead at small dataset sizes. The \code{george} solver scales as $N^3$ but benefits from highly optimized C/LAPACK routines, making it competitive for $N \lesssim 500$. The \code{spotgp} banded solver carries higher per-evaluation overhead than \code{tinygp} at small $N$, but because the bandwidth $b$ is set by the kernel correlation length and does not grow with $N$, its $\mathcal{O}(Nb^2)$ scaling is effectively linear for long baselines. At $N = 10{,}000$, the banded solver reaches wall-clock times comparable to \code{tinygp} (${\sim}0.5$\,s), while \code{george} exceeds 30\,s. The banded solver overtakes the full $\mathcal{O}(N^3)$ Cholesky around $N \sim 500$.

The middle panel shows matrix storage requirements. The \code{celerite} solver stores only $\mathcal{O}(NJ)$ elements (where $J = 2N_{\rm harm} + 1$ is the number of semi-separable terms), keeping memory below a few MB even at $N = 10{,}000$. The banded solver stores the $(2b+1) \times N$ band, which grows linearly with $N$ for fixed $b$. In contrast, the dense solvers (\code{george} and \code{spotgp} full) store the full $N \times N$ covariance matrix, reaching ${\sim}800$\,MB at $N = 10{,}000$. For the full \code{spotgp} solver run on the GPU, the max matrix storage size is limited to the memory capacity of the GPU, which is typically 16--32\,GB for consumer hardware (in this case, an Nvidia RTX 5070 TI with 16\,GB of memory), thus limiting the maximum $N$ to a few thousand. Similarly for the full \code{george} solver, the max matrix size is limited by the RAM capacity of the CPU, which is typically 16--64\,GB for consumer hardware, thus also limiting the maximum $N$ to a few thousand. In comparison, the banded solver's linear memory scaling allows it to handle much larger datasets without hitting memory limits.

The bottom panel shows the sparsity fraction $(2b+1)/N$ of the banded covariance matrix. Because the bandwidth is determined by the kernel's correlation length (a few rotation periods), it remains fixed as the observation baseline grows. The sparsity fraction therefore drops as $1/N$, falling below 10\% for $N \gtrsim 1000$, which is the regime where the banded solver's advantage is most pronounced.

The \code{celerite} solver's $\mathcal{O}(N)$ scaling is only available for kernels that decompose into a sum of (damped) exponentials and cosines. Our trapezoidal and skew-normal envelopes do not admit this decomposition, so the banded Cholesky solver is the only option for the general kernel. For applications where speed is critical and the one-sided exponential envelope is a sufficient approximation (e.g., for noise modeling,  or applications like lightcurve detrending where the user does not care about interpreting the hyperparameters), the \code{celerite} decomposition provides a practical fast path.

\begin{figure}[ht!]
    \centering
    \script{plot_benchmark_tinygp_vs_spotgp.py}
    \includegraphics[width=0.8\linewidth]{figures/fig_benchmark_tinygp_vs_spotgp.pdf}
    \caption{Benchmark of GP solver performance as a function of the number of datapoints $N$, using a one-sided exponential envelope with $N_{\rm harm} = 2$ harmonics and cadence $\Delta t = 1$\,d. \emph{Top:} Wall-clock time for a single log-likelihood evaluation. \code{tinygp} (c\'{e}lerite, green) scales as $\mathcal{O}(N)$; \code{george} (red) and \code{spotgp} full Cholesky (teal) scale as $\mathcal{O}(N^3)$; \code{spotgp} banded (dark blue) scales as $\mathcal{O}(Nb^2)$. Horizontal gray lines mark 1\,s and 1\,min. \emph{Middle:} Matrix storage in MB. The \code{celerite} and banded solvers scale linearly with $N$, while the dense solvers scale as $N^2$. \emph{Bottom:} Sparsity fraction $(2b+1)/N$ of the banded covariance matrix, showing that the bandwidth $b$ remains fixed as $N$ grows. All JAX solvers are JIT-compiled and warmed up before timing.}
    \label{fig:benchmark_packages}
\end{figure}

\subsection{Comparison with \code{starry\_process}}
\label{sec:starry_comparison}

An alternative GP framework for stellar variability is the \code{starry\_process} of \citet{luger_mapping_2021,luger_mapping_2021}, which models the stellar surface as a linear combination of spherical harmonics $Y_{\ell m}$ and derives the covariance of the disk-integrated flux analytically. Our approach differs in several fundamental respects, which we summarize here. 

\textit{Surface representation: }
\code{starry\_process} represents the stellar surface intensity as a truncated spherical harmonic expansion, $I(\theta,\phi) = \sum_{\ell,m} y_{\ell m}\,Y_{\ell m}(\theta,\phi)$, with the flux computed via the linear design-matrix relation $\mathbf{f} = \mathbf{A}\,\mathbf{y}$, where $\mathbf{A}$ encodes the rotational geometry and limb darkening \citep{luger_starry_2019}. This yields a non-parametric surface model: the coefficients $y_{\ell m}$ are treated as Gaussian random variables, and the GP covariance of the flux follows from the prior covariance of the harmonic coefficients, $\boldsymbol{\Sigma}_f = \mathbf{A}\,\boldsymbol{\Sigma}_y\,\mathbf{A}^\top$. By contrast, our model parameterizes the surface in terms of individual circular spots (each described by a latitude $\Phi$, longitude evolution $\lambda(t)$, and trapezoidal size envelope $\alpha(t)$), and derives the kernel by taking the autocorrelation of the resulting flux analytically.

\textit{Temporal evolution: }
A key distinction is the treatment of time dependence. The \code{starry\_process} covariance is \emph{non-stationary}: it depends on absolute time through the rotation matrix $\mathbf{A}(t)$, since the surface pattern is treated as fixed (or drawn once) and the star merely rotates. There is no built-in mechanism for spot emergence, evolution, or decay; the surface is effectively a static random field viewed from different angles. Our kernel, by contrast, is \emph{stationary} (depending only on the time lag $\tau = |t_2 - t_1|$) and explicitly encodes spot lifetimes through the envelope autocorrelation $R_\Gamma(\tau)$, which models the birth, plateau, and decay of each spot. This stationarity arises because we marginalize over the random arrival times of spots (Section~\ref{subsec:multi_spot}), treating the spot process as a temporal Poisson process. As a consequence, our kernel naturally captures the finite coherence time of the variability signal: the autocorrelation decays to zero at $\tau = \ell_{\rm spot} + 2\tau_{\rm spot}$, a feature that has no analog in \code{starry\_process}.

\textit{Differential rotation: }
Our model includes differential rotation through the shear parameter $\kappa$, which allows spots at different latitudes to rotate at different angular velocities (Section~\ref{sec:hyperparameters}). This produces observable effects including peak broadening, decoherence, and the coherence lag $\tau_{\rm coh}$ (Figure~\ref{fig:decoherence_samples}). The \code{starry\_process} framework assumes solid-body rotation.

\textit{Spot parameterization: }
In \code{starry\_process}, the surface statistics are controlled by the angular power spectrum of the spherical harmonic coefficients (parameterized by a characteristic angular scale and latitude-dependent modulation). These parameters describe the \emph{spatial} statistics of the surface but do not map directly onto physical spot properties. In our model, the hyperparameters have direct physical interpretations: $P_{\rm eq}$ is the equatorial rotation period, $\ell_{\rm spot}$ is the spot plateau lifetime, $\tau_{\rm spot}$ is the emergence/decay timescale, $\alpha_{\max}$ is the peak angular spot radius, and $\kappa$ is the differential rotation shear.

\textit{Computational structure: }
The \code{starry\_process} covariance matrix $\boldsymbol{\Sigma}_f$ is dense and generally $O(N^3)$ to factorize, as it is for any GP. The computational advantage of that framework lies in the analytic evaluation of the matrix entries (avoiding numerical integration over the stellar surface). Our kernel shares this advantage -- its evaluation reduces to sums of cosines weighted by tabulated Fourier coefficients -- but additionally benefits from sparsity: $K(\tau) = 0$ for $|\tau| > \ell_{\rm spot} + 2\tau_{\rm spot}$, which produces a banded covariance matrix with bandwidth $b = (\ell_{\rm spot} + 2\tau_{\rm spot})/\Delta t$. This enables $O(Nb^2)$ Cholesky factorization via banded linear algebra, a significant speedup for long time series with fine cadence.

\textit{Degeneracies and information content: }
\citet{luger_mapping_2021} showed that there exist fundamental degeneracies in mapping stellar surfaces from disk-integrated photometry alone: infinitely many surface maps can produce identical lightcurves. These degeneracies persist in both frameworks. In \code{starry\_process}, they manifest as a null space of the design matrix $\mathbf{A}$; in our model, they appear as the well-known degeneracy between $N_{\rm spot}$, $f_{\rm spot}$, and $\alpha_{\max}$ (which are absorbed into the single amplitude parameter $\sigma_k^2$; Eq.~\ref{eq:kernel_amplitude}), and the fact that the kernel constrains only the latitude-averaged squared Fourier coefficients $\langle |c_n|^2 \rangle_\Phi$, not the surface map itself. However, our kernel provides direct constraints on \emph{temporal} properties (spot lifetime, evolution timescale, differential rotation) that are not accessible from the static surface model of \code{starry\_process}.


\vspace{\baselineskip}
% \clearpage
% -----------------------------------------------
\section{Inference Conditions} \label{sec:inference_conditions}

Before applying our GP model to real data, we first test its ability to recover known parameters from synthetic datasets. We generate mock lightcurves by simulating the flux using the same spot model (Equation~\ref{eqn:flux_model}) that underlies our kernel, with known hyperparameters $\boldsymbol{\theta}_{\rm true}$. We then add Gaussian noise (with a standard deviation $\sigma_n$) to the simulated fluxes to create realistic observations. 


\subsubsection{Kernel Convergence and Data Requirements}

As a first test, we estimate the minimum observation duration ($T_{\rm obs}$) and sampling cadence ($\Delta t$) required for the ACF to converge within a certain error of the true kernel function. To perform this test, we generate a grid of synthetic timeseries with varying $T_{\rm obs}$, $\Delta t$, and $\sigma_n$, and numerically compute the ACF for each individual lightcurve. We then compare the empirical ACF to the true kernel function (computed using the known parameters) and compute the mean squared error (MSE) as a function of $T_{\rm obs}$, $\Delta t$, and $\sigma_n$. This analysis allows us to identify the regimes in which the ACF provides a reliable estimate of the kernel, which is critical for informing the design of observing campaigns and setting expectations for parameter recovery.

Figures~\ref{fig:mse_vs_ndata} shows the results of this test for two sampling cadences: $\Delta t = 0.5$\,d (representative of space-based photometry such as \textit{Kepler} or \textit{TESS}) and $\Delta t = 5.0$\,d (representative of ground-based monitoring campaigns). In each panel, the $x$-axis gives the number of datapoints in the lightcurve $N = T_{\rm obs}/\Delta t$, and the $y$-axis gives the MSE between the sample ACF (computed from a single simulated lightcurve) and the analytic kernel, averaged over 10 random realizations with $1\sigma$ errorbars. Different colored lines correspond to different Gaussian noise levels $\sigma_{\rm noise}$, spanning from $10^{-4}$ (photon-noise-limited space photometry) to $10^{-1}$ (heavily contaminated data). The spot emergence rate is held fixed at $\dot{N}_{\rm spot} = 0.25$\,spots\,day$^{-1}$, so that longer baselines contain proportionally more spots and the kernel amplitude is independent of $T_{\rm obs}$ (Section~\ref{subsec:multi_spot}).

Several trends are apparent. First, for low noise levels ($\sigma_{\rm noise} \leq 10^{-3}$), the MSE decreases steadily with increasing $N$, roughly as a power law, reflecting the improved statistical convergence of the sample ACF as more independent spot events contribute to the estimate. At the highest noise levels ($\sigma_{\rm noise} = 10^{-1}$), the noise dominates the signal and the MSE saturates at ${\sim}10^{-1}$ regardless of the number of datapoints, indicating that the spot-induced correlations are buried in the noise floor.

Second, comparing the two cadences, the coarser sampling ($\Delta t = 5.0$\,d) yields systematically higher MSE at fixed $N$ for the intermediate noise levels ($\sigma_{\rm noise} = 10^{-2}$), because fewer samples per rotation period degrade the ACF estimate at short lags. For the fine cadence ($\Delta t = 0.5$\,d), the low-noise curves ($\sigma_{\rm noise} \leq 10^{-3}$) converge to MSE $\lesssim 10^{-3}$ with $N \gtrsim 10^3$ datapoints (corresponding to $T_{\rm obs} \gtrsim 500$\,d), while the coarse cadence requires $N \gtrsim 10^2$--$10^3$ points ($T_{\rm obs} \gtrsim 500$--$5000$\,d) to reach comparable accuracy.

These results suggest that for typical spot parameters, reliable kernel estimation from a single lightcurve requires observation baselines spanning at least ${\sim}10$--$20$ spot lifetimes ($\ell_{\rm spot} + 2\tau_{\rm spot} = 21$\,d in our test case), with noise levels $\sigma_{\rm noise} \lesssim 10^{-2}$ in relative flux units.

\begin{figure}[ht!]
    \centering
    \script{plot_kernel_mse_vs_ndata.py}
    % TODO: replace with 0p5d and 5p0d once generated
    \includegraphics[width=0.49\linewidth]{figures/fig_kernel_mse_vs_ndata_0p5d.pdf}
    \includegraphics[width=0.49\linewidth]{figures/fig_kernel_mse_vs_ndata_5p0d.pdf}
    \caption{MSE between the analytic kernel and the sample ACF computed from individual simulated lightcurves, as a function of the number of datapoints $N$, for a sampling cadence of $\Delta t = 0.5$\,d (left) and $\Delta t = 5.0$\,d (right). Each line corresponds to a different Gaussian noise level $\sigma_{\rm noise}$ added to the simulated flux. Points show the mean over 10 random realizations; errorbars indicate the $1\sigma$ scatter. The spot emergence rate is $\dot{N}_{\rm spot} = 0.25$\,spots\,day$^{-1}$, so $N_{\rm spot}$ scales with the observation duration.}
    \label{fig:mse_vs_ndata}
\end{figure}

% -----------------------------------------------
\subsection{Testing the Cranmer--Rao Bound of Different Parameter Regimes} \label{sec:cranmer_rao}

Even though we are only modeling a six parameter kernel model, there are many possible combinations of parameters that can produce similar kernel shapes, leading to degeneracies and non-identifiability. For example, if the spot lifetime is very short compared to the rotation period, the kernel will be dominated by the envelope autocorrelation and may not contain enough information about the periodic structure to constrain $P_{\rm eq}$ or $\kappa$. Conversely, if the spot lifetime is relatively long compared to the rotation period or observational baseline, the kernel may be dominated by the periodic structure and have a very slowly decaying ACF, which could also make it difficult to constrain the envelope parameters $\ell_{\rm spot}$ and $\tau_{\rm spot}$. 

Before performing MCMC sampling (which is computationally intensive to do across a wide parameter space), we can first test the theoretical identifiability of the parameters by computing the Fisher information matrix (Eq.~\ref{eq:fisher}) across a grid of parameters. The inverse of the Fisher information matrix gives the Cramér--Rao lower bound on the covariance of any unbiased estimator of the parameters, which provides a theoretical limit on how well we can constrain the parameters from data generated by our model. By analyzing the eigenvalues and eigenvectors of the Fisher information matrix, we can identify which combinations of parameters are most and least identifiable, and how this depends on the underlying spot parameters. This analysis will inform our expectations for MCMC sampling and help us interpret the results in terms of parameter degeneracies and uncertainties.

To systematically explore parameter space, we draw $N = 4096$ parameter combinations from a scrambled Sobol quasi-random sequence spanning the bounds listed in Table~\ref{tab:crb_bounds}. For each sample $\boldsymbol{\theta}_k$, we construct a synthetic lightcurve of duration $T_{\rm sim} = 365$~d sampled at a cadence of $\Delta t = 0.5$~d (i.e., $N_t = 730$ observations) with white noise $\sigma_n = 10^{-4}$ and a mean spot emergence rate of $\dot{N}_{\rm spot} = 0.25$~d$^{-1}$. We then evaluate the expected Fisher information matrix $\mathcal{I}(\boldsymbol{\theta}_k)$ via Eq.~\eqref{eq:fisher} and invert it to obtain the CRB covariance matrix $\Sigma_{\rm CRB} = \mathcal{I}^{-1}$. Cases where any eigenvalue of $\mathcal{I}$ falls below a regularization floor of $10^{-6}$ are flagged as numerically singular and excluded from the analysis, as the corresponding CRB is unreliable. From the CRB covariance we extract the marginal standard deviations $\sigma_{\rm CRB}(\theta_i) = \sqrt{\Sigma_{ii}}$ and the pairwise Pearson correlation coefficients $r_{ij} = \Sigma_{ij}/(\sigma_i \sigma_j)$ for all parameter pairs.

\begin{table}[ht!]
    \centering
    \caption{Parameter bounds used for the Sobol quasi-random sweep of the CRB analysis.}
    \label{tab:crb_bounds}
    \begin{tabular}{lcc}
        \hline
        Parameter & Lower & Upper \\
        \hline
        $P_{\rm eq}$ [d] & 0.1 & 50 \\
        $\kappa$ & $-1.0$ & 1.0 \\
        $I$ [deg] & 0 & 90 \\
        $\ell_{\rm spot}$ [d] & 1.0 & 100 \\
        $\tau_{\rm spot}$ [d] & 1.0 & 100 \\
        $\sigma_k$ & $10^{-3}$ & 1 \\
        \hline
    \end{tabular}
\end{table}

\begin{figure}[ht!]
    \centering
    \script{plot_crb_param_sweep.py}
    \includegraphics[width=\linewidth]{figures/crb_sweep_corner_contour_kappa.png}
    \caption{Cramér--Rao bound parameter sweep corner plot showing the theoretical lower bounds for $\kappa$ uncertainty across different parameter regimes. To compute the average percent uncertainty for each parameter region, we computed 4096 simulations within the bounds of Table~\ref{tab:crb_bounds} using Sobol quasi-random sampling. The average percent uncertainty on $\kappa$ is indicated by the color scale, with dark blue indicating tighter constraints and dark red indicating looser constraints. For some regions, with short periods ($P_{\rm eq} \lesssim 20$), long spot lifetimes ($\ell_{\rm spot} + 2 \tau_{\rm spot} \gtrsim 10 P_{\rm eq}$), and high inclinations ($I \gtrsim 45\,$deg), the differential rotation shear $\kappa$ can theoretically be constrained within $\lesssim 10\%$ uncertainty.}
    \label{fig:crb_sweep_corner_kappa}
\end{figure}

\begin{figure}[ht!]
    \centering
    \script{plot_crb_param_sweep.py}
    \includegraphics[width=\linewidth]{figures/crb_sweep_corner_contour_inc.png}
    \caption{Same as Fig.~\ref{fig:crb_sweep_corner_kappa}, colored by inclination $I$.}
    \label{fig:crb_sweep_corner_inc}
\end{figure}

\begin{figure}[ht!]
    \centering
    \script{plot_crb_param_sweep.py}
    \includegraphics[width=\linewidth]{figures/crb_sweep_corner_contour_peq.png}
    \caption{Same as Fig.~\ref{fig:crb_sweep_corner_kappa}, colored by $P_{\rm eq}$. $P_{\rm eq}$ is by far the easiest of the parameters to constrain. With the exception of a few regions, e.g., systems with short spot lifetimes relative to their period ($\ell_{\rm spot} + 2\tau_{\rm spot} \lesssim P_{\rm eq}$), at pole-on inclinations ($I \lesssim 30\,$deg), the equitorial period can generally be constrained to better than 10\% uncertainty,and in many cases better than 1\% uncertainty.}
    \label{fig:crb_sweep_corner_peq}
\end{figure}

\begin{figure}[ht!]
    \centering
    \script{plot_crb_param_sweep.py}
    \includegraphics[width=\linewidth]{figures/crb_sweep_corner_contour_sigma_k.png}
    \caption{Same as Fig.~\ref{fig:crb_sweep_corner_kappa}, colored by $\sigma_k$.}
    \label{fig:crb_sweep_corner_sigma_k}
\end{figure}

\clearpage

% -----------------------------------------------
\subsection{MCMC Inference} \label{sec:mcmc_results}

While Section~\ref{sec:inference_conditions} gives a sense of the minimum observation requirements for reliable ACF estimation, it doesn't directly test the ability of our GP model to recover the underlying parameters $\boldsymbol{\theta}$ and their uncertainties from noisy data. To assess this, we perform MCMC sampling of the posterior distribution $p(\boldsymbol{\theta} | \mathbf{y})$ using the sampling methods described in Section~\ref{sec:mcmc}. More specifically, we investigate how the noise level and observation duration affect the ability to recover the true parameters, and how this compares to the theoretical limits set by the Cramér--Rao bound. We also analyze the posterior distributions to identify any degeneracies or correlations between parameters, and to understand how well the model can constrain each parameter under different conditions.

As in the CRB analysis (Section \ref{sec:cranmer_rao}), we simulate lightcurve data with a white noise level of $\sigma_n = 10^{-4}$ ($\sim 100$~ppm), but using a sampling cadence of $\Delta t = 1.0$~d (i.e., $N_t = 365$ and $1460$ observations for $T_{\rm sim} = 365$ and $1460$~d). Given that \kepler has a finer cadence (1 minute for the short cadence or 30 min for the long cadence data), and typical instrumental noise levels of less than 50~ppm \cite{gilliland2010}, we assume that these simulated values are reasonable for \kepler-like data.

For each synthetic lightcurve, we ran both NUTS (with \code{BlackJAX}) and nested sampling (with \code{dynesty}) to obtain posterior samples of the parameters $\boldsymbol{\theta}$ and compare the results to the true parameters $\boldsymbol{\theta}_{\rm true}$ used to generate the data. In practice, we found that the NUTS sampler tended to get stuck in small local modes of the posterior and was less effective at exploring the full parameter space, and so we primarily report results from the nested sampling runs, which were more robust to multimodality and provided better coverage of the posterior.

% Generated from src/data by src/scripts/table_simulated_posteriors.py (Snakefile rule simulated_posterior_table)
\variable{output/simulated_posterior_table.tex}

\begin{figure}[ht!]                                                                           
    \centering 
    \script{plot_simulated_posteriors.py}
    \begin{tikzpicture}   
        \node[anchor=south west, inner sep=0] (img90) at (0,0)                                         
            {\includegraphics[width=0.65\linewidth]{figures/simulated_corner_dynesty_kde_90.pdf}};                 
        \node[anchor=south west, inner sep=0] (img30) at (6.5,6.5)
            {\includegraphics[width=0.65\linewidth]{figures/simulated_corner_dynesty_kde_30.pdf}};          
    \end{tikzpicture}
    \caption{Posterior distributions for different inclinations: $I = 90$ (bottom left) and $30$ (upper right) degrees based on nested sampling runs on synthetic lightcurves. The blue posterior outlines the constraints using 1 year of data, while the green posterior outlines the constraints using 4 years of data (comparable to the Kepler mission duration), and the true parameters are indicated by the dashed black lines. The posterior distributions show that the model can recover the true parameters with tight precision given a typical noise level of $\sigma_n = 10^{-4}$ and a spot emergence rate of $\dot{N}_{\rm spot} = 0.25$\,d$^{-1}$.}
    \label{fig:simulated_corners}
\end{figure}

To see how the parameter uncertainties carry through to the covariance structure, Figure~\ref{fig:simulated_kernel_psd_posterior} shows the kernel and PSD evaluated at random draws from the posteriors in Figure~\ref{fig:simulated_corners}. At $I = 90^\circ$, both baselines recover the true kernel and PSD. At $I = 30^\circ$, the 1-year posterior underestimates the correlation at short lags, which matches its inclination posterior lying above the truth ($I = 36 \pm 3^\circ$): a higher inclination deepens the rotational modulation. With 4 years of data, the posterior kernel follows the true kernel.

\begin{figure}[ht!]
    \centering
    \script{plot_simulated_kernel_posteriors.py}
    \includegraphics[width=\linewidth]{figures/simulated_kernel_psd_posterior_90.pdf}
    \includegraphics[width=\linewidth]{figures/simulated_kernel_psd_posterior_30.pdf}
    \caption{Kernel and PSD evaluated at 2000 random draws from the nested sampling posteriors in Figure~\ref{fig:simulated_corners}, for $I = 90^\circ$ (top) and $I = 30^\circ$ (bottom). The left panels show the normalized kernel $K(\tau)/K(0)$, with each draw normalized by its own $K(0)$. The right panels show the PSD up to the Nyquist frequency of the $\Delta t = 1$\,d sampling. Solid lines show the posterior median, and the dark and light shaded regions show the pointwise 68\% and 95\% intervals for 1 year (blue) and 4 years (green) of data. The dashed black lines show the kernel and PSD at the true parameters. The lower panels show the residuals from the truth: $\Delta K/K(0)$ for the kernel and $P/P_{\rm true} - 1$ for the PSD. The fractional PSD residual spikes near the nulls of the PSD, where $P_{\rm true} \rightarrow 0$.}
    \label{fig:simulated_kernel_psd_posterior}
\end{figure}


\clearpage
% ========================================================================================
\section{Discussion} \label{sec:discussion}

\subsection{Summary of Results} \label{sec:summary_results}

In this work, we derived a physically motivated GP kernel for stellar photometric variability caused by rotating starspots and systematically characterized the performance and limitations of this methodology using simulations; the application to real data is the subject of Paper~II. We summarize the main results below.

\textbf{Analytic kernel derivation (Section~\ref{sec:kernel_derivation}).}
Starting from a forward model of the flux deficit produced by circular spots on a differentially rotating star, we decomposed the single-spot signal into a product of an envelope term $\Gamma(t)$ (capturing spot emergence and decay) and a visibility term $\mathcal{V}(t)$ (capturing rotational modulation and limb geometry). By taking the Fourier transform of each component and applying the Wiener--Khinchin theorem, we obtained a closed-form expression for the autocovariance kernel $\mathcal{K}(\tau)$ as a sum over rotation harmonics weighted by latitude-averaged visibility coefficients $\langle |c_n|^2 \rangle_\Phi$ and modulated by the envelope autocorrelation $R_\Gamma(\tau)$ (Eq.~\ref{eq:kernel_full}). We provided analytic $R_\Gamma(\tau)$ expressions for trapezoidal, exponential, and skew-normal spot profiles, and validated all derivations symbolically with \code{SymPy} and numerically against brute-force simulations.

\textbf{Physical interpretability of hyperparameters (Section~\ref{sec:hyperparameters}).}
The six free kernel hyperparameters, $P_{\rm eq}$, $\kappa$, $I$, $\ell_{\rm spot}$, $\tau_{\rm spot}$, and $\sigma_k$, map directly onto stellar and spot properties, unlike the phenomenological parameters of standard quasi-periodic or SHO kernels. We showed that $P_{\rm eq}$ and $\kappa$ control the periodicity and harmonic structure of the kernel; $I$ controls the balance between the DC and oscillating components via the visibility coefficients; $\ell_{\rm spot}$ and $\tau_{\rm spot}$ set the correlation lifetime; and the amplitude parameters ($\alpha_{\max}$, $f_{\rm spot}$, $N_{\rm spot}$) are degenerate and collapse into a single scale parameter $\sigma_k^2$. We further demonstrated that our kernel captures signatures, such as PSD broadening from differential rotation and harmonic ratio variations with inclination, that standard kernels cannot reproduce (Section~\ref{sec:comparison_literature}).

\textbf{Efficient GP solver (Section~\ref{sec:gp_solver}).}
Because the kernel $\mathcal{K}(\tau)$ has compact support (set by the spot lifetime $\ell_{\rm spot} + 2\tau_{\rm spot}$), the covariance matrix is banded with bandwidth $b \propto (\ell_{\rm spot} + 2\tau_{\rm spot})/\Delta t$. Exploiting this structure, we implemented a banded Cholesky solver that reduces the GP inference cost from $\mathcal{O}(N^3)$ to $\mathcal{O}(Nb^2)$. The solver is fully JIT-compiled and autodifferentiable via JAX, enabling gradient-based sampling with no additional implementation effort.

\textbf{Inference conditions and Cram\'{e}r--Rao bounds (Sections~\ref{sec:inference_conditions}--\ref{sec:cranmer_rao}).}
Using synthetic lightcurves, we showed that reliable kernel estimation requires observation baselines spanning at least $\sim$10--20 spot lifetimes with noise levels $\sigma_{\rm noise} \lesssim 10^{-2}$ in relative flux. We computed the Cram\'{e}r--Rao lower bound across a grid of 4096 parameter combinations and found that the differential rotation coefficient $\kappa$ is best constrained for short-period stars ($P_{\rm eq} \lesssim 20$\,d) with long-lived spots and moderate-to-high inclinations, where $\kappa$ can theoretically be measured to $\lesssim 10\%$ uncertainty.

\textbf{Application to real data (Paper~II).}
The companion paper extends the kernel into an additive composite model, with several spot populations sharing one stellar geometry and with granulation and white-noise terms, and applies the framework to the four-year \kepler lightcurve of Kepler-63, comparing the inferred inclination, differential rotation, and spot lifetimes with independent constraints from spectroscopy and transit spot mapping.

% -----------------------------------------------
\subsection{Applications and Future Work} \label{sec:applications}

\subsubsection{Direct Applications for this model}

There are numerous potential applications for this starspot kernel. Firstly, it can be used as a tool for inferring star spot parameters using only photometric lightcurves, without the need for spectroscopic or transit data. This is particularly useful for stars that do not have transiting planets or for which high-resolution spectroscopy is not available. By fitting the kernel to the lightcurve, one can extract information about the star's rotation period, differential rotation, inclination, and spot lifetimes, which are important for understanding stellar magnetic activity and its impact on exoplanet habitability. 

Additionally, the kernel can be used to model and mitigate stellar variability in the context of exoplanet transit searches and radial velocity measurements, improving the precision of planet detection and characterization. The physically motivated nature of the kernel also allows for more interpretable modeling of stellar variability compared to phenomenological kernels, enabling insights into the underlying stellar physics driving the observed lightcurve features.

\subsubsection{Extensions to Spectroscopy}

This model can also be used in combination with spectroscopic data to jointly constrain stellar and spot properties. For example, our GP kernel can be used to infer the inclination angle of a star just from the photometric variability. This means that we can combine this information with spectroscopic measurements of $v_*\sin I$, allowing us to infer the equatorial rotation velocity $v_{\rm eq}$ and thus the stellar radius $R_* = v_{\rm eq} P_{\rm eq} / 2\pi$. This methodology would allow us to constrain stellar radii independently of stellar evolution/isochrone models.

Another natural extension of this framework is to high-resolution spectroscopic time series, where starspots induce wavelength-dependent distortions in absorption line profiles as they rotate across the stellar disk. In the Doppler imaging framework \citep[e.g.,][]{luger_mapping_2021-2}, a spot at latitude $\Phi$ and longitude $\Lambda(t)$ perturbs the line profile at a radial velocity $v(t) = v_{\rm eq} \sin I \cos\Phi \sin\Lambda(t)$, which is directly related to the same visibility geometry (Eq.~\ref{eq:cos_beta}) that underpins our photometric kernel. The Fourier decomposition and Wiener--Khinchin machinery developed in Section~\ref{sec:kernel_derivation} should therefore generalize to a two-dimensional spectral-temporal kernel $\mathcal{K}(\tau, \Delta v)$, describing the covariance of spectral perturbations across both time lag and velocity separation. Such an extension could help break degeneracies that are present in photometry alone: in particular, the inclination $I$ and spot latitude $\Phi$ enter differently through the line-of-sight velocity than through the projected area, and differential rotation $\kappa$ produces directly observable velocity-track separations for spots at different latitudes. Joint photometric--spectroscopic inference with a physically motivated kernel could therefore improve constraints on stellar surface properties, particularly in the context of extreme precision radial velocity (EPRV) surveys where disentangling stellar activity from planetary signals remains a central challenge.

\subsubsection{Further Development}

We make the kernel implementation publicly available as a Python package, \code{spotgp}, which provides a user-friendly interface for evaluating the kernel and performing GP regression with it. The package includes functions for computing the kernel given a set of hyperparameters, as well as utilities for sampling the hyperparameters from data using MCMC. We also provide example Jupyter notebooks demonstrating how to use the kernel for modeling stellar lightcurves and inferring starspot properties. 

Furthermore, we have implemented the package in a modular way such that users can easily adapt the kernel to their specific needs and assumptions (such as different spot latitude distributions, envelope functions, or spot parameter distributions). We hope that \code{spotgp} will be a valuable resource for the community and facilitate further research into stellar variability and starspot modeling.

% -----------------------------------------------
\subsection{Limitations} \label{sec:limitations} 

Our kernel model is based on several simplifying assumptions, such as circular spots with uniform contrast, a fixed spot latitude distribution, and independent spot evolution. In reality, starspots can have more complex shapes, non-uniform contrasts, and form in groups or active regions, which could introduce additional correlations in the lightcurve that are not captured by our model. Additionally, the single-population kernel derived here does not account for other forms of astrophysical variability, such as asteroseismic pulsations, flares, or granulation, which may be present in the lightcurve; additive kernel terms for granulation and white noise, and for several coexisting spot populations, are introduced in Paper~II.

In the analysis of this paper, we have focused on dark starspots, however bright faculae can also contribute to photometric variability, particularly in Sun-like stars. In the next paper of this series, we will derive a kernel for multiband photometric data and show how dark spots and bright faculae can be disentangled using their different wavelength-dependent contrasts.

Our kernel is also derived under the assumption that all spots are drawn from the same distributions for $\Phi$, $\ell_{\rm spot}$, $\tau_{\rm spot}$, and $\alpha_{\max}$. In reality, there may be correlations between these parameters (e.g., larger spots may have longer lifetimes), which could affect the shape of the kernel. We also assume that these starspot properties are constant in time, such that the kernel is stationary. In practice, spot properties may evolve over time (e.g., due to activity cycles), leading to non-stationary behavior in the lightcurve. Future work could explore more complex models that allow for such correlations. However, the current model provides a useful starting point for understanding the key physical parameters that shape the photometric variability of spotted stars and can be extended in future work to relax some of these assumptions. We also provide several examples on the online documentation for how users can test non-stationary kernels and more complex spot models using the modular implementation of \code{spotgp}.

\vspace{\baselineskip}
% ========================================================================================
\section{Conclusion} \label{sec:conclusion}

In conclusion, we have derived an analytic, physically motivated Gaussian process kernel for stellar photometric variability driven by rotating starspots. Starting from a forward model of circular spots on a differentially rotating star, we obtained a closed-form autocovariance function whose six hyperparameters, $P_{\rm eq}$, $\kappa$, $I$, $\ell_{\rm spot}$, $\tau_{\rm spot}$, and $\sigma_k$, map directly onto measurable stellar and spot properties. This stands in contrast to the phenomenological kernels widely used in the literature, whose parameters lack clear physical interpretation.

Using simulated lightcurves and MCMC/nested sampling, we demonstrated that the kernel can recover injected stellar parameters, including the stellar inclination and differential rotation coefficient $\kappa$, given sufficient observational baseline ($\gtrsim$10 spot lifetimes) and photometric precision ($\sigma_{\rm noise} \lesssim 10^{-2}$). Using Cram\'{e}r--Rao analysis we further showed that $\kappa$ is best constrained for short-period, high-inclination stars with long-lived spots. The application of the method to the \kepler lightcurve of Kepler-63 is presented in Paper~II.

The compact support of the kernel enables a banded Cholesky GP solver that scales as $\mathcal{O}(Nb^2)$ rather than $\mathcal{O}(N^3)$, making inference tractable for the long-baseline, densely sampled lightcurves produced by \kepler\ and \tess. Combined with a JAX-based implementation that supports automatic differentiation and gradient-based sampling, this framework is well suited for population-level studies of stellar surface properties.

Looking ahead, systematic application of this kernel to large photometric surveys could yield differential rotation and spot lifetime measurements for thousands of stars across a range of spectral types and ages: quantities that are currently available for only a handful of targets. Such measurements would place new empirical constraints on stellar dynamo theory, the evolution of surface magnetic activity, and the impact of stellar variability on exoplanet detection and habitability assessments.

% ========================================================================================

\software{
\code{JAX} \citep{bradbury_jax_2018},
\code{NumPy} \citep{harris_numpy_2020},
\code{SciPy} \citep{virtanen_scipy_2020},
\code{Matplotlib} \citep{hunter_matplotlib_2007},
\code{Astropy} \citep{astropy_collaboration_astropy_2013, astropy_collaboration_astropy_2018},
\code{BlackJAX} \citep{lao_blackjax_2020},
\code{SymPy} \citep{meurer_sympy_2017}
}

\begin{acknowledgments}
    EA was supported by National Aeronautics and Space Administration Exoplanet Research Program (NASA XRP) grant 80NSSC21K1111.  This work was performed in part by the Virtual Planetary Laboratory Team, which is a member of the NASA Nexus for Exoplanet System Science, and funded via the NASA Astrobiology Program ICAR Grant 80NSSC23K1398.
\end{acknowledgments}

\clearpage
\bibliographystyle{aasjournal}
\bibliography{references.bib}


% ========================================================================================
\appendix

\section{Fourier Transform Properties and Pairs} \label{sec:fourier_appendix}

\begin{table}[ht!]
\centering
\caption{Properties of the Fourier transform. We use the convention $\FT{f}(\omega) = \int_{-\infty}^{\infty} f(t)\,e^{-i\omega t}\,dt$. Here $a, b$ are constants, $t_0$ and $\omega_0$ are fixed shifts, and $f$, $g$ are functions with transforms $\FT{f}$, $\FT{g}$.}
\label{tab:fourier_properties}
\setcounter{ftprop}{0}
\begin{tabular}{cllr}
    \hline
    \textbf{No.} & \textbf{Property} & \textbf{Time domain $\longleftrightarrow$ Frequency domain} \\
    \hline
    \refstepcounter{ftprop}\label{ft:linearity}\theftprop &
    Linearity &
    $a\,f(t) + b\,g(t) \;\longleftrightarrow\; a\,\FT{f}(\omega) + b\,\FT{g}(\omega)$ \\[8pt]
    \refstepcounter{ftprop}\label{ft:timeshift}\theftprop &
    Time shifting &
    $f(t - t_0) \;\longleftrightarrow\; e^{-i\omega t_0}\,\FT{f}(\omega)$ \\[8pt]
    \refstepcounter{ftprop}\label{ft:freqshift}\theftprop &
    Frequency shifting &
    $e^{i\omega_0 t}\,f(t) \;\longleftrightarrow\; \FT{f}(\omega - \omega_0)$ \\[8pt]
    \refstepcounter{ftprop}\label{ft:convolution}\theftprop &
    Convolution &
    $(f \ast g)(t) \;\longleftrightarrow\; \FT{f}(\omega)\cdot\FT{g}(\omega)$ \\[8pt]
    \refstepcounter{ftprop}\label{ft:multiplication}\theftprop &
    Multiplication &
    $f(t)\cdot g(t) \;\longleftrightarrow\; \dfrac{1}{2\pi}\,\FT{f}(\omega) \ast \FT{g}(\omega)$ \\[12pt]
    \refstepcounter{ftprop}\label{ft:even}\theftprop &
    Even symmetry &
    $f(t)$ real and even $\;\Longrightarrow\; \FT{f}(\omega)$ purely real \\[8pt]
    \refstepcounter{ftprop}\label{ft:odd}\theftprop &
    Odd symmetry &
    $f(t)$ real and odd $\;\Longrightarrow\; \FT{f}(\omega)$ purely imaginary \\[8pt]
    \hline
\end{tabular}
\end{table}

\begin{table}[ht!]
\centering
\caption{Fourier transform pairs used in this work. We use the convention $\FT{f}(\omega) = \int_{-\infty}^{\infty} f(t)\,e^{-i\omega t}\,dt$ and $\mathrm{sinc}(x) \equiv \sin(x)/x$.}
\label{tab:fourier_pairs}
\setcounter{ftpair}{0}
\begin{tabular}{clll}
    \hline
    \textbf{No.} & \textbf{Function} & $f(t)$ & $\FT{f}(\omega)$ \\
    \hline
    \refstepcounter{ftpair}\label{fp:rect}\theftpair &
    Rectangular pulse &
    $\mathrm{rect}_W(t) = \begin{cases} 1 & |t| \leq W/2 \\ 0 & \text{otherwise} \end{cases}$ &
    $W\,\mathrm{sinc}\!\left(\dfrac{\omega W}{2}\right)$ \\[16pt]
    \refstepcounter{ftpair}\label{fp:cosine}\theftpair &
    Cosine &
    $\cos(\omega_0 t + \phi)$ &
    $\pi\left[e^{i\phi}\,\delta(\omega - \omega_0) + e^{-i\phi}\,\delta(\omega + \omega_0)\right]$ \\[10pt]
    \refstepcounter{ftpair}\label{fp:exp}\theftpair &
    Complex exponential &
    $e^{i\omega_0 t}$ &
    $2\pi\,\delta(\omega - \omega_0)$ \\[10pt]
    \refstepcounter{ftpair}\label{fp:sqwave}\theftpair &
    Square wave &
    $\mathrm{sgn}(\cos\omega_0 t)$ &
    $\displaystyle 2\pi\sum_{n\,\text{odd}} \frac{(-1)^{(n-1)/2}}{n}\;\delta(\omega - n\omega_0)$ \\[14pt]
    \refstepcounter{ftpair}\label{fp:delta}\theftpair &
    Dirac delta &
    $\delta(t - t_0)$ &
    $e^{-i\omega t_0}$ \\[8pt]
    \refstepcounter{ftpair}\label{fp:comb}\theftpair &
    Dirac comb &
    $\displaystyle\sum_n \delta(t - nT)$ &
    $\dfrac{2\pi}{T}\displaystyle\sum_n \delta(\omega - n\omega_0)$, \; $\omega_0 = 2\pi/T$ \\[14pt]
    \hline
\end{tabular}
\end{table}

% ========================================================================================

\section{Fourier Coefficients of the Visibility Function} \label{app:fourier_visibility}

Here we derive the Fourier coefficients $c_n$ of the visibility function $\mathcal{V}(t) = \max\{\cos\beta(t),\,0\}$ from first principles.\footnote{Symbolic verification: \variable{output/derivations/derive_fourier_cosine.tex}}

\subsection{Setup}

From Equation~\ref{eq:cos_beta}, we have:
\begin{equation}
    \cos\beta(t) = \underbrace{\cos I\sin\Phi}_{b_0} + \underbrace{\sin I\cos\Phi}_{b_1}\cos(\Lambda_0 + \omega_0 t).
\end{equation}
This is a shifted cosine oscillating between $b_0 - b_1$ and $b_0 + b_1$. The spot is visible when $\cos\beta(t) > 0$, which (for $b_1 > 0$, $|b_0| < b_1$) occurs when:
\begin{equation}
    |\omega_0 t + \Lambda_0| < \theta_{\rm vis}, \qquad \theta_{\rm vis} = \arccos\!\left(-\frac{b_0}{b_1}\right).
    \label{eq:theta_vis}
\end{equation}
This defines the \emph{duty cycle} $\Delta = 2\theta_{\rm vis}/(2\pi) = \theta_{\rm vis}/\pi$ of the visibility window. When $b_0 + b_1 < 0$ the spot is never visible; when $b_0 - b_1 > 0$ the spot is always visible ($\theta_{\rm vis} = \pi$). Figure~\ref{fig:visibility_function_annotated} illustrates the components of the visibility function $\mathcal{V}(t)$t.

\begin{figure}[ht!]
    \centering
    \script{plot_spot_visibility_annotated.py}
    \includegraphics[width=\linewidth]{figures/fig_visibility_annotated.pdf}
    \caption{Visibility function $\mathcal{V}(t)$ for a single spot. The spot is visible when $\cos\beta(t) > 0$, which occurs for a fraction $\Delta = 2\theta_{\rm vis}/(2\pi)$ of the rotation period. The visibility function is a truncated cosine, with amplitude $b_1$ and offset $b_0$.}
    \label{fig:visibility_function_annotated}
\end{figure}

Since $\mathcal{V}(t)$ is periodic with period $P = 2\pi/\omega_0$, we expand it in a Fourier series:
\begin{equation}
    \mathcal{V}(t) = \sum_{n=-\infty}^{\infty} c_n\, e^{in\omega_0 t},
\end{equation}
where the Fourier coefficients are:
\begin{equation}
    c_n = \frac{1}{T}\int_0^T \mathcal{V}(t)\, e^{-in\omega_0 t}\, dt = \frac{1}{2\pi}\int_{-\theta_{\rm vis}}^{\theta_{\rm vis}} \big(b_0 + b_1\cos\theta\big)\, e^{-in\theta}\, d\theta,
    \label{eq:cn_integral}
\end{equation}
with the substitution $\theta = \omega_0 t + \Lambda_0$ (and without loss of generality setting $\Lambda_0 = 0$, since a phase shift only affects the complex phase of $c_n$ but not $|c_n|^2$).

\subsection{Evaluating the Integral}

We split the integrand into the constant and cosine parts:
\begin{equation}
    c_n = \frac{1}{2\pi}\left[b_0\int_{-\theta_{\rm vis}}^{\theta_{\rm vis}} e^{-in\theta}\, d\theta + b_1\int_{-\theta_{\rm vis}}^{\theta_{\rm vis}} \cos\theta\; e^{-in\theta}\, d\theta\right].
    \label{eq:cn_split}
\end{equation}

\paragraph{First integral (constant part).}
For $n \neq 0$:
\begin{equation}
    \int_{-\theta_{\rm vis}}^{\theta_{\rm vis}} e^{-in\theta}\, d\theta = \left[\frac{e^{-in\theta}}{-in}\right]_{-\theta_{\rm vis}}^{\theta_{\rm vis}} = \frac{e^{-in\theta_{\rm vis}} - e^{in\theta_{\rm vis}}}{-in} = \frac{2\sin(n\,\theta_{\rm vis})}{n}.
    \label{eq:cn_int_const}
\end{equation}
For $n = 0$ this gives $2\theta_{\rm vis}$.

\paragraph{Second integral (cosine part).}
Using $\cos\theta = \tfrac{1}{2}(e^{i\theta} + e^{-i\theta})$:
\begin{equation}
    \int_{-\theta_{\rm vis}}^{\theta_{\rm vis}} \cos\theta\; e^{-in\theta}\, d\theta = \frac{1}{2}\int_{-\theta_{\rm vis}}^{\theta_{\rm vis}} e^{i(1-n)\theta}\, d\theta + \frac{1}{2}\int_{-\theta_{\rm vis}}^{\theta_{\rm vis}} e^{-i(1+n)\theta}\, d\theta.
\end{equation}
Each integral is of the same form as the first. For $m \neq 0$:
\begin{equation}
    \int_{-\theta_{\rm vis}}^{\theta_{\rm vis}} e^{im\theta}\, d\theta = \frac{2\sin(m\,\theta_{\rm vis})}{m}.
\end{equation}
Therefore, for $n \neq \pm 1$:
\begin{equation}
    \int_{-\theta_{\rm vis}}^{\theta_{\rm vis}} \cos\theta\; e^{-in\theta}\, d\theta = \frac{\sin\big((1-n)\theta_{\rm vis}\big)}{1-n} + \frac{\sin\big((1+n)\theta_{\rm vis}\big)}{1+n}.
    \label{eq:cn_int_cos}
\end{equation}

\paragraph{Combining.}
Substituting back into Eq.~\eqref{eq:cn_split}, for $n \neq 0, \pm 1$:
\begin{equation}
    c_n = \frac{1}{\pi}\left[\frac{b_0\sin(n\,\theta_{\rm vis})}{n} + \frac{b_1}{2}\left(\frac{\sin\big((n-1)\theta_{\rm vis}\big)}{n-1} + \frac{\sin\big((n+1)\theta_{\rm vis}\big)}{n+1}\right)\right],
    \label{eq:cn_result}
\end{equation}
where we used $\sin\big((1-n)x\big) = -\sin\big((n-1)x\big)$ and $1/(1-n) = -1/(n-1)$.


\vspace{\baselineskip}
% ========================================================================
\section{Fourier Transform of the Trapezoidal Envelope} \label{app:fourier_trapezoid}

Here we derive the Fourier transform and autocorrelation function of the trapezoidal envelope from Section~\ref{sec:spot_profile}; the transform of the unsquared trapezoid $\alpha(t)$ is derived in \NotebookLink{ft_trapezoid.ipynb}. Symbolic verification of these equations using \code{sympy} can be found at: \variable{output/derivations/derive_fourier_gamma.tex}

% ----------------------------------------------------------------
\subsection{Piecewise Definition of $\Gamma(t)$}

The trapezoidal spot evolution $\alpha(t)$, centered at $t = 0$, has a flat plateau of duration $\ell_{\rm spot}$ at amplitude $\alpha_{\max}$ and linear ramps of duration $\tau_{\rm spot}$ on each side:
\begin{equation}
    \alpha(t) = \alpha_{\max}
    \begin{cases}
        0 & t < -\tfrac{\ell_{\rm spot}}{2} - \tau_{\rm spot} \\[4pt]
        \displaystyle\frac{t + \tfrac{\ell_{\rm spot}}{2} + \tau_{\rm spot}}{\tau_{\rm spot}}
            & -\tfrac{\ell_{\rm spot}}{2} - \tau_{\rm spot} \leq t < -\tfrac{\ell_{\rm spot}}{2} \\[6pt]
        1 & -\tfrac{\ell_{\rm spot}}{2} \leq t \leq \tfrac{\ell_{\rm spot}}{2} \\[4pt]
        \displaystyle\frac{-t + \tfrac{\ell_{\rm spot}}{2} + \tau_{\rm spot}}{\tau_{\rm spot}}
            & \tfrac{\ell_{\rm spot}}{2} < t \leq \tfrac{\ell_{\rm spot}}{2} + \tau_{\rm spot} \\[6pt]
        0 & t > \tfrac{\ell_{\rm spot}}{2} + \tau_{\rm spot}
    \end{cases}
    \label{eq:alpha_piecewise}
\end{equation}

Squaring and normalizing the trapezoid~\eqref{eq:alpha_piecewise}:
\begin{equation}
    \Gamma(t) = \frac{\alpha^2(t)}{\alpha_{\max}^2} =
    \begin{cases}
        0 & |t| > \ell/2 + \tau \\[4pt]
        \displaystyle\left(\frac{|t| - \ell/2 - \tau}{-\tau}\right)^{\!2}
            & \ell/2 < |t| \leq \ell/2 + \tau \\[8pt]
        1 & |t| \leq \ell/2
    \end{cases}
    \label{eq:Gamma_piecewise}
\end{equation}
where $\ell \equiv \ell_{\rm spot}$ and $\tau \equiv \tau_{\rm spot}$ for brevity. The ramps are now \emph{parabolic} (concave) rather than linear.

% ----------------------------------------------------------------
\subsection{Direct Integration}

Since $\Gamma(t)$ is real and even, $\FT{\Gamma}(\omega)$ is purely real:
\begin{equation}
    \FT{\Gamma}(\omega) = 2\int_0^{\infty} \Gamma(t)\cos(\omega t)\,dt
    = 2\left[\underbrace{\int_0^{\ell/2}\cos(\omega t)\,dt}_{I_1}
    + \underbrace{\int_{\ell/2}^{\ell/2+\tau}\left(\frac{\ell/2 + \tau - t}{\tau}\right)^{\!2}\cos(\omega t)\,dt}_{I_2}\right].
\end{equation}

\paragraph{Integral $I_1$ (plateau):}
\begin{equation}
    I_1 = \frac{\sin(\omega\ell/2)}{\omega}.
    \label{eq:I1_plateau}
\end{equation}

\paragraph{Integral $I_2$ (parabolic ramp):}
With $a = \ell/2$ and $b = \ell/2 + \tau$, we integrate by parts twice. Setting $f(t) = \big((b-t)/\tau\big)^2$:

\emph{First integration by parts} ($u = f(t)$, $dv = \cos(\omega t)\,dt$):
\begin{equation}
    I_2 = \left[\frac{f(t)\sin(\omega t)}{\omega}\right]_a^b + \frac{2}{\tau^2\omega}\int_a^b (b-t)\sin(\omega t)\,dt.
\end{equation}
The boundary term gives $0 - f(a)\sin(\omega a)/\omega = -\sin(\omega\ell/2)/\omega$ (since $f(a) = 1$).

\emph{Second integration by parts} on $J = \int_a^b (b-t)\sin(\omega t)\,dt$ ($u = b-t$, $dv = \sin(\omega t)\,dt$):
\begin{equation}
    J = \left[-\frac{(b-t)\cos(\omega t)}{\omega}\right]_a^b - \frac{1}{\omega}\int_a^b \cos(\omega t)\,dt
    = \frac{\tau\cos(\omega\ell/2)}{\omega} - \frac{\sin(\omega b) - \sin(\omega\ell/2)}{\omega^2}.
    \label{eq:J_ibp}
\end{equation}

Combining:
\begin{equation}
    I_2 = -\frac{\sin(\omega\ell/2)}{\omega} + \frac{2\cos(\omega\ell/2)}{\tau\omega^2} - \frac{2\big[\sin\!\big(\frac{\omega\ell}{2}+\omega\tau\big) - \sin\!\big(\frac{\omega\ell}{2}\big)\big]}{\tau^2\omega^3}.
    \label{eq:I2_ramp}
\end{equation}

\paragraph{Combining $I_1 + I_2$:}
The $\sin(\omega\ell/2)/\omega$ terms from $I_1$ and $I_2$ cancel:
\begin{equation}
    I_1 + I_2 = \frac{2\cos(\omega\ell/2)}{\tau\omega^2} + \frac{2\sin(\omega\ell/2)}{\tau^2\omega^3} - \frac{2\sin\!\big(\frac{\omega\ell}{2}+\omega\tau\big)}{\tau^2\omega^3}.
    \label{eq:I1_plus_I2}
\end{equation}

Therefore:
\begin{equation}
    \FT{\Gamma}(\omega) = \frac{4}{\tau^2\omega^3}\left[\tau\omega\cos\!\left(\frac{\omega\ell}{2}\right) + \sin\!\left(\frac{\omega\ell}{2}\right) - \sin\!\left(\frac{\omega\ell}{2} + \omega\tau\right)\right].
    \label{eq:Gamma_hat_direct}
\end{equation}

\vspace{\baselineskip}
% ========================================================================
\section{Autocorrelation of the Squared Envelope} \label{app:acf_trapezoid}
% ========================================================================

Next we derive the closed-form autocorrelation $R_\Gamma(\tau) = \int_{-\infty}^{\infty} \Gamma(s)\,\Gamma(s+\tau)\,ds$ of the normalized squared trapezoidal envelope $\Gamma(t) = \alpha^2(t)/\alpha_{\max}^2$ defined in Eq.~\eqref{eq:Gamma_piecewise}. The \code{sympy} symbolic verification can be found at: \variable{output/derivations/derive_R_Gamma.tex}

% ----------------------------------------------------------------
\subsection{Setup}

The function $\Gamma(t)$ has three regions: a plateau of height $1$ for $|t| \leq \ell/2$, parabolic ramps for $\ell/2 < |t| \leq S$, and zero otherwise. The autocorrelation integral
\begin{equation}
    R_\Gamma(\tau) = \int_{-S}^{S-\tau} \Gamma(s)\,\Gamma(s+\tau)\,ds, \qquad \tau \geq 0,
\end{equation}
must be split into sub-intervals where both $\Gamma(s)$ and $\Gamma(s+\tau)$ have known polynomial forms. The breakpoints of $\Gamma(s)$ are at $s = \pm\ell/2$ and $s = \pm S$, while those of $\Gamma(s+\tau)$ are shifted by $-\tau$. The relevant sub-intervals and the character of the overlap change as $\tau$ increases through the four intervals $[0,\,\tau_{\rm s}]$, $[\tau_{\rm s},\,\ell]$, $[\ell,\,\ell+\tau_{\rm s}]$, and $[\ell+\tau_{\rm s},\,\ell+2\tau_{\rm s}]$. The result below holds whenever $\ell \geq \tau_{\rm s}$ (plateau at least as long as one ramp), which is typical for observed starspots; for $\ell < \tau_{\rm s}$ the intervals reorder and the closed form no longer applies.

% ----------------------------------------------------------------
\subsection{Result}

Evaluating the polynomial integrals in each sub-interval and summing (verified symbolically and numerically), the autocorrelation is:
\begin{equation}
    \boxed{
    R_\Gamma(\tau) =
    \begin{cases}
        \displaystyle \ell + \frac{2\tau_{\rm s}}{5}
        - \frac{4\tau^2}{3\tau_{\rm s}}
        + \frac{2\tau^3}{3\tau_{\rm s}^2}
        - \frac{\tau^5}{15\tau_{\rm s}^4}
        & 0 \leq \tau \leq \tau_{\rm s} \\[12pt]
        \displaystyle \ell + \frac{2\tau_{\rm s}}{3} - \tau
        & \tau_{\rm s} \leq \tau \leq \ell \\[12pt]
        P_3(\tau)
        & \ell \leq \tau \leq \ell + \tau_{\rm s} \\[12pt]
        \displaystyle \frac{(\ell + 2\tau_{\rm s} - \tau)^5}{30\,\tau_{\rm s}^4}
        & \ell + \tau_{\rm s} \leq \tau \leq \ell + 2\tau_{\rm s} \\[12pt]
        0 & \tau > \ell + 2\tau_{\rm s}
    \end{cases}
    }
    \label{eq:R_Gamma_closed}
\end{equation}
where the third piece is the degree-5 polynomial:
\begin{align}
    P_3(\tau) &= \frac{\tau^5}{30\tau_{\rm s}^4}
    - \frac{(\ell + 2\tau_{\rm s})\,\tau^4}{6\tau_{\rm s}^4}
    + \frac{(\ell^2 + 4\ell\tau_{\rm s} + 2\tau_{\rm s}^2)\,\tau^3}{3\tau_{\rm s}^4} \notag \\
    &\quad - \frac{\ell(\ell^2 + 6\ell\tau_{\rm s} + 6\tau_{\rm s}^2)\,\tau^2}{3\tau_{\rm s}^4}
    + \frac{(\ell^4 + 8\ell^3\tau_{\rm s} + 12\ell^2\tau_{\rm s}^2 - 6\tau_{\rm s}^4)\,\tau}{6\tau_{\rm s}^4} \notag \\
    &\quad + \frac{-\ell^5 - 10\ell^4\tau_{\rm s} - 20\ell^3\tau_{\rm s}^2 + 30\ell\tau_{\rm s}^4 + 20\tau_{\rm s}^5}{30\tau_{\rm s}^4}.
    \label{eq:P3}
\end{align}

\begin{figure*}
    \centering
    \script{plot_envelope_autocorrelation.py}
    \includegraphics[width=\textwidth]{figures/fig_envelope_autocorrelation.pdf}
    \caption{Piecewise structure of the envelope autocorrelation. \emph{Left:} Fourier transform $\FT{\Gamma}(\omega)$ (black) decomposed into the plateau sinc-like term $I_1$ (blue dashed) and the parabolic ramp term $I_2$ (orange dashed). \emph{Right:} Autocorrelation $R_\Gamma(\tau)$ from Eq.~\eqref{eq:R_Gamma_closed}, with each of the four piecewise intervals shown in a different color. The thin gray curve shows the result of direct numerical convolution, confirming the closed-form expressions. Parameters: $\ell = 15$\,d, $\tau_{\rm s} = 3$\,d.}
    \label{fig:envelope_autocorrelation}
\end{figure*}

Figure~\ref{fig:envelope_autocorrelation} illustrates these results. The left panel shows the Fourier transform $\FT{\Gamma}(\omega)$ decomposed into the plateau contribution $I_1$ and the parabolic ramp contribution $I_2$. The right panel shows $R_\Gamma(\tau)$ with the four piecewise intervals color-coded, overlaid on a numerical convolution for verification.


% ========================================================================
\section{Verified Equations} \label{app:equation_checks}
% ========================================================================

Table~\ref{tab:equation_checks} lists the equations marked with \EquationCheckIcon in the margin. Each is transcribed once into \code{SymPy}, in an equation registry (\texttt{src/derivations/equations.py}), and three checks run against that transcription in every build of this paper. The notebook that derives the equation must reproduce it symbolically. Its \LaTeX\ source in this manuscript must be the text the transcription was last reviewed against. Where the table names a \code{spotgp} function, that function must match the transcription in value and in gradient with respect to every continuous argument, at random parameter values and on both sides of every piecewise breakpoint, with the transcription evaluated in 50-digit arithmetic. The build fails if any check fails, so every equation marked in this version of the paper has passed.

\begin{table*}[ht!]
\centering
\caption{Equations verified symbolically and, where a function is named, against \code{spotgp}. The errors are the largest differences between \code{spotgp} and the transcribed equation, relative to the largest value (gradient) over each sweep of the equation's variable.}
\label{tab:equation_checks}
\footnotesize
\setlength{\tabcolsep}{5pt}
\variable{output/equation_checks.tex}
\end{table*}

\end{document}
